\documentclass[11pt,a4paper,oneside]{amsbook}
\usepackage[T1]{fontenc}
\usepackage[utf8]{inputenc}
\usepackage{lmodern}
\usepackage[american]{babel}
\usepackage{microtype}
\usepackage{geometry}
\geometry{margin=1.1in}
\makeatletter
\def\input@path{{../}}
\makeatother
\usepackage{sga1-en}
\usepackage[hidelinks]{hyperref}

\addto\captionsamerican{\renewcommand{\chaptername}{Expos\'e}}

\title[SGA~1, Expos\'e XIII (English draft)]{Cohomological properness of sheaves of sets and of sheaves of non-commutative groups\\
{\normalsize SGA~1, Expos\'e XIII\\English draft}}
\author{M.\ Raynaud}
\date{}
\hypersetup{
  pdftitle={SGA 1, Expose XIII: Cohomological properness of sheaves of sets and of sheaves of non-commutative groups (English draft)},
  pdfauthor={M. Raynaud},
  pdfsubject={Unofficial English translation of SGA 1, Expose XIII}
}

\begin{document}
\frontmatter
\maketitle

\chapter*{About this draft}
This is an unofficial English translation of Expos\'e~XIII,
``Propret\'e cohomologique des faisceaux d'ensembles et des faisceaux de groupes non commutatifs'', from SGA~1 (after unpublished notes of A.\ Grothendieck). The source volume is
A.\ Grothendieck and M.\ Raynaud,
\textit{Rev\^etements \'etales et groupe fondamental} (SGA~1),
S\'eminaire de g\'eom\'etrie alg\'ebrique du Bois Marie, 1960--61.
It follows the slightly corrected SMF recomposition,
\href{https://arxiv.org/abs/math/0206203v2}{arXiv:math/0206203v2}.

This draft contains the whole expos\'e, including proofs, footnotes,
diagrams, the updating remarks added in 2003 by M.~Raynaud (between bold
brackets and marked (MR)) where the expos\'e has them, and its bibliography.
Statement numbering and mathematical notation follow the source.
References to other expos\'es print the source's numbers. Apparent
misprints in the source have been retained; they are recorded in the
accompanying README, separately from the translated text. Scholarly
proofreading remains outstanding.

The translator's contribution is licensed under
the MIT License.
The French original remains copyright of its authors and publishers.

\tableofcontents
\mainmatter
\setcounter{chapter}{12}
\chapter[Cohomological properness of sheaves of sets...]{Cohomological properness\\ of sheaves of sets and\\ of sheaves of non-commutative groups}
\label{XIII}
% original p. 344

\begin{center}
{by Mme M.\ \textsc{Raynaud}\footnote{After unpublished notes of
A\ptbl Grothendieck.}}
\end{center}
\vspace*{1cm}

This expos\'e proposes to use \'etale cohomology in order to
generalize certain results of \Ref{IX} and \Ref{X}. It also shows
how one can extend to sheaves of not necessarily commutative groups
those results of~SGA~5~II which still make sense for such
sheaves. The notions of \'etale cohomology expounded in~SGA~4 are
assumed known.

The main result \eqref{XIII.2.4} gives an important example
of a non-proper morphism $f\colon~U\to~S$ which is ``cohomologically
proper in dimension~$\leq1$'', that is to say such that, for
certain sheaves of groups $F$ on~$U$ (in the sense of the \'etale
topology), the formation of~$f_*F$ and $\R^1f_*F$ commutes with every
change of base~$S'\to~S$. This property is indeed
satisfied by the open set $U$ of a scheme $X$ proper over~$S$,
complementary to a divisor $D$ with normal crossings
relative to~$S$, at least if one requires $F$ to be
constant finite, of order prime to the residue
characteristics of~$S$. If one no longer assumes $F$ of order prime
to the residue characteristics of~$S$, one has an
analogous result on replacing $\R^1f_*F$ by the subsheaf
$\R^1_{\tame}f_*F$ obtained by restricting oneself to considering the torsors
under~$F$ ``tamely ramified on~$X$ relative
to~$S$''. In particular this makes it possible to show that the tamely
ramified fundamental group of a proper and smooth algebraic
curve over a separably closed field, with a finite number of
closed points removed, is topologically of finite type
\eqref{XIII.2.12}.

% original p. 345
No.~\Ref{XIII.4} is devoted to the homotopy exact sequence and to
the K\"unneth formula.

Finally an appendix gives useful variants of Abhyankar's lemma
proved in X.\Ref{X.3.6}.

\setcounter{section}{-1}

\section{Recollections on the theory of stacks}
\label{XIII.0}

We shall use in what follows the theory of stacks
expounded in \cite{XIII.1} and \cite{XIII.2}. We restrict ourselves to the case of the
\'etale site of a scheme. Given a scheme $X$, let us denote by
$X_\et$ the \'etale site of~$X$. Recall that a
stack $F$
on~$X$ is a fibered category over~$X_\et$ such
that, for every scheme $X'$ \'etale over~$X$ and for every pair
of objects
$x$, $y$ of the fiber $F_{X'}$, the presheaf
$\SheafHom_{X'}(x,y)$ is a sheaf, and such that, for every
surjective \'etale morphism $X'' \to X'$, every object of~$F_{X''}$
endowed with a descent datum relative to $X'' \to X'$ is the
inverse image of an object of~$F_{X'}$.

We denote by $F(X')$ the category of cartesian sections of
$F/X'$. More generally, if $\Sch_X$ is the category
of schemes over~$X$ endowed with the \'etale topology, the
stack $F$ can be extended to a stack $\cal{F}$ on~$\Sch_X$ and,
for every morphism $f\colon X' \to X$, we again denote by $F(X')$ the
category of cartesian sections of this stack $\cal{F}$
over~$X'$.

A gerbe
is a stack such that, for every scheme $X'$ \'etale
over~$X$ and for every pair of objects $x$, $y$ of~$F_{X'}$, every
morphism from~$x$ to $y$ is an isomorphism and $x$ and $y$ are
locally isomorphic, and such that the set of objects $X'$ of
$X_\et$ such that $F_{X'}$ is nonempty is a refinement of
$X_\et$. For example the stack of torsors under a sheaf of groups
is a gerbe which, moreover, has a cartesian section.
Conversely a gerbe which has a section, \ie such that there
exists an object $x$ of~$F_X$, is equivalent to the stack of
torsors under the sheaf
% original p. 346
of groups $\SheafAut_X(x)$.

There is an obvious notion of subgerbe and of maximal subgerbe
of a stack~$F$. Given a cartesian section $x$ of
$F(X)$, there exists a unique maximal subgerbe $G_x$ of~$F$ such
that $x$ factors through~$G_x$. We call $G_x$ the
subgerbe generated by~$x$; it is by definition a trivial
gerbe.
The presheaf $SF$
defined by
$$
SF(X') = \{ \text{maximal subgerbes of }F|X' \}
$$
is a sheaf, called the sheaf of maximal subgerbes
of~$F$.
Let $O$ be the presheaf defined by
$$
O(X') = \{ \text{classes of objects of }F_{X'}\text{ mod. isomorphism}
\} \text{.}
$$
By associating to every object $x$ of~$F_{X'}$ the maximal subgerbe
of~$F|X'$ generated by~$x$, we obtain a morphism
$$
O \to SF \text{;}
$$
by \cite[III 2.1.4]{XIII.2}, this morphism makes~$SF$ a sheaf
associated to~$O$.

A stack $F$ is said to be \emph{constructible}
(\emph{\resp ind-$\LL$-finite},
$\LL$ being a set of prime numbers) if, for every
scheme $X'$ \'etale over~$X$ and for every object $x$ of~$F_{X'}$,
the same holds for the sheaf $\SheafAut_{X'}(x)$ \cite[VII 2.2.1]{XIII.2}. A stack is
said to be $1$-constructible
if it is constructible and if the sheaf of maximal subgerbes is
constructible.

\section{Cohomological properness}
\label{XIII.1}
%
\setcounter{subsection}{-1}
\subsection{}
\label{XIII.1.0}
%
Let $S$ be a scheme, $f \colon X \to Y$ a morphism of
$S$-schemes. If $S'$ is an $S$-scheme, we consider the
following diagram, all of whose squares are cartesian:
% original p. 347
\begin{equation*}
\label{eq:XIII.1.0.1}
\tag{\thesubsection.1}
\begin{array}{c}
\xymatrix{ X \ar[d]_f & \ar[l]_h X'
\ar[d]^{f'} \\ Y \ar[d] & \ar[l]_g Y' \ar[d] \\ S & \ar[l] \,S'. }
\end{array}
\end{equation*}

If $Y_1$ is a scheme \'etale over~$Y$, we set $X_1 =
X \times_Y Y_1$, $Y_1' = Y' \times_Y Y_1$, and we consider the
cartesian square
\begin{equation*}
\label{eq:XIII.1.0.2}
\tag{\thesubsection.2}
\begin{array}{c}
\xymatrix{ X_1 \ar[d]_{f_1} & \ar[l]_{h_1}
X_1' \ar[d]^{f_1'} \\ Y_1 & \ar[l]^{g_1} \,Y_1'. }
\end{array}
\end{equation*}

\begin{definition}
\label{XIII.1.1}
Let $F$ be a stack on~$X$. We say that $(F,f)$ is
\emph{cohomologically proper relative to~$S$ in dimension
$\leqslant -1$}
(\resp in dimension $\leqslant 0$, \resp in dimension $\leqslant 1$)
if, for every $S$-scheme $S'$, the canonical functor (defined
in the obvious way by the universal property of
the inverse image of stacks):
$$
g^* f_* F \to f'_* h^* F \qquad\text{ (\cf 1.0.1)}
$$
is faithful (\resp fully faithful, \resp an equivalence
of categories).
\end{definition}

If there is no possible confusion about~$S$, in particular if $S=Y$,
we say cohomologically proper instead of cohomologically proper
relative to~$S$.

\subsection{}
\label{XIII.1.2}
Let $F$ be a sheaf of sets on~$X$; let $\Phi$ be the stack in
discrete categories associated to $F$, \ie the stack whose
fiber over every scheme $X_1$ \'etale over~$X$ is the
discrete category having $F(X_1)$ as its set of objects. We
say that $(F,f)$ is cohomologically proper relative to~$S$ in
dimension $\leqslant -1$ (\resp in dimension $\leqslant 0$) if $(\Phi,
f)$ is cohomologically proper relative to~$S$ in dimension
$\leqslant 0$ (\resp in dimension $\leqslant 1$).

The
% original p. 348
canonical morphism
\begin{equation*}
\label{eq:XIII.1.2.1}
\tag{\thesubsection.1} { g^* f_* F \to f'_* h^* F }
\end{equation*}
gives, on passing to the associated stacks in discrete
categories, the canonical morphism
$$
g^* f_* \Phi \to f'_* h^* \Phi \text{.}
$$
Consequently, to say that $(F,f)$ is cohomologically proper relative
to~$S$ in dimension $\leqslant -1$ (\resp in dimension $\leqslant
0$) is equivalent to saying that, for every $S$-scheme $S'$, the
morphism \eqref{eq:XIII.1.2.1} is injective (\resp bijective).

\subsection{}
\label{XIII.1.3}
Let $F$ be a sheaf of groups on~$X$ and $\Phi$ the stack of
torsors on~$X$ with group $F$ \cite[II 2.3.2]{XIII.1}. We say that $(F,f)$ is
cohomologically proper relative to~$S$ in dimension
$\leqslant -1$ (\resp $\leqslant 0$, \resp $\leqslant 1$) if
$(\Phi,f)$ is cohomologically proper relative to~$S$ in
dimension $\leqslant -1$ (\resp $\leqslant 0$, \resp $\leqslant
1$). The condition of cohomological properness can be made explicit
as follows.

\begin{subproposition}
\label{XIII.1.3.1}
The notation is that of \eqref{eq:XIII.1.0.1} and
\eqref{eq:XIII.1.0.2}. Let $F$ be a sheaf of groups on~$X$. We
denote by $F'$ (\resp $F_1$, \resp $F_1'$, etc.) the inverse
image of~$F$ on~$X'$ (\resp on~$X_1$, \resp on~$X_1'$,
etc.). Then the following conditions are equivalent:
\begin{enumerate}
\item[(i)] $(F,f)$ is cohomologically proper relative to~$S$
in dimension $\leqslant -1$ (\resp $\leqslant 0$, \resp $\leqslant
1$).
\item[(ii)] For every morphism $S' \to S$, for every scheme $Y_1$
\'etale over~$Y$, and for every torsor $P$ on~$X_1$ with
group $F_1$, if $\leftexp{P\mkern-4mu}{F}_1$ denotes the group $F_1$
twisted by $P$ \textup{\cite[II 4.1.2.3]{XIII.1}}, the canonical morphism
$$
a_0 \colon g_1^*(f_{1 *}(\leftexp{P\mkern-4mu}{F}_1)) \to f'_{1
*}(\leftexp{P'\mkern-6mu}{F}'_1)
$$
is injective (\resp $a_0$ is bijective and the canonical morphism
$$
a_1 \colon g^*(\R^1 f_* F) \to \R^1 f'_* F'
$$
is
% original p. 349
injective, \resp $a_0$ and $a_1$ are bijective).
\item[(ii~bis)] For every morphism $S' \to S$, for every scheme
$Y_1$ \'etale over~$Y$, for every torsor $P$ on~$X_1$ with
group $F_1$, and for every torsor $R$ under $\leftexp{P\mkern-4mu}{F}_1$, the
canonical morphism
$$
\alpha_0 \colon g_1^*(f_{1 *}R) \to f'_{1 *} R'
$$
is injective (\resp $\alpha_0$ is bijective, \resp the morphisms
$\alpha_0$ and
$$
\alpha_1 \colon g_1^*(\R^1 f_{1 *}(\leftexp{P\mkern-4mu}{F}_1)) \to \R^1
f'_{1 *}(\leftexp{P'\mkern-6mu}{F}'_1)
$$
are bijective).
\end{enumerate}
\end{subproposition}

\subsubsection*{Proof}
(i)$\To$(ii~bis). By \cite[II 4.2.5]{XIII.1} every torsor $R$
with group $\leftexp{P\mkern-4mu}{F}_1$ is of the form $R = Q
\overset{F_1}{\wedge} P^\circ$, where $Q$ is a torsor with group
$F_1$ and $P^\circ$ the opposite of~$P$. We then have $R' \simeq Q'
\overset{F'_1}{\wedge} {P'}^\circ$. Let $\Phi$ be the stack of torsors
under $F$ and let $x$, $y$ (\resp~$x'$,~$y'$) be the objects of the
fiber category $(g^* f_* \Phi)_{Y_1'}$ (\resp $(f'_*
\Phi')_{Y_1'}$) associated to $P$, $Q$ (\resp $P'$,~$Q'$). We have
the relation
$$
Q\overset{F_1}{\wedge} P^\circ \simeq \SheafHom_{F_1}(P,Q) \text{,}
$$
and it follows that we have canonical isomorphisms
$$
\SheafHom_{Y'_1}(x,y) \simeq g_1^* f_{1 *}(Q
\overset{F_1}{\wedge} P^\circ) \text{,}\quad \SheafHom_{Y_1'}(x',y')
\simeq f'_{1 *}(Q' \overset{F'_1}{\wedge}{P'}^\circ)\text{.}
$$
Consequently the morphism $\alpha_0$ is identified with the morphism
$$
\SheafHom_{Y'_1}(x,y) \to \SheafHom_{Y'_1}(x',y') \text{,}
$$
from which it follows that, if $(F,f)$ is cohomologically proper
relative to~$S$ in dimension $\leqslant -1$, $\alpha_0$ is
injective and that, if $(F,f)$ is cohomologically proper in dimension
$\leqslant 0$, $\alpha_0$ is bijective.

Suppose now that $(F,f)$ is cohomologically proper
relative to~$S$ in dimension $\leqslant 1$, \ie that the
canonical morphism
% original p. 350
$$
\varphi \colon g^* f_* \Phi \to f'_* \Phi'
$$
is an equivalence. Let $G$ be the sheaf of maximal subgerbes
of the stack $f_* \Phi$ \hbox{\cite[III~2.1.8]{XIII.1}}; we then have an
isomorphism $G \simeq \R^1 f_* F$. Since $g^* G$ is the
sheaf of maximal subgerbes of~$g^* f_* \Phi$ \cite[III
2.1.5.5]{XIII.2}, the morphism $\alpha_1$ is obtained from
$\varphi|Y'_1$ by taking the sheaves of maximal subgerbes,
hence is an isomorphism.

(ii~bis)$\To$(ii). It suffices to show that, if the
morphisms $\alpha_0$ are bijective, then the morphisms $a_1$ are
injective. Let $Y'_1$ be a scheme \'etale over~$Y'$, and
$s$ and $t$ two elements of~$g^*(\R^1 f_* F)(Y'_1)$
having the same image in $\R^1 f'_* F'(Y'_1)$, and let us show that
we have $s=t$. The assertion is local for the \'etale topology of
$Y'_1$ and, taking into account the definition of the inverse image
$g^*(\R^1 f_* F)$, we may assume that $Y'_1$ is the inverse
image of a scheme $Y_1$ \'etale over~$Y$ and that $s$ and
$t$ come from torsors $P$ and $Q$ on~$X_1$. The hypothesis
made on~$s$ and $t$ then means that the inverse images $P'$ and
$Q'$ of~$P$ and $Q$ on~$X'_1$ are isomorphic locally for the
\'etale topology of~$Y'_1$. If we set $R =
\SheafHom_{F_1}(P,Q)$, the fact that the morphism
$$
g_1^* f_{1 *} R \to f'_{1 *} R'
$$
is bijective proves that $P$ and $Q$ are isomorphic locally for the
\'etale topology of~$Y_1$, hence that we have $s=t$.

(ii)$\To$(i). To prove that $\varphi$ is faithful
(\resp fully faithful), it suffices to show that, if $Y_1$ is
a scheme \'etale over~$Y$, if $P$, $Q$ are two torsors on~$X_1$ with group~$F_1$, and if $x$, $y$ (\resp $x'$, $y'$) are the objects
of~$(g^* f_* \Phi)_{Y'_1}$ (\resp $(f'_* \Phi')_{Y'_1}$)
associated to $P$, $Q$ (\resp $P'$, $Q'$), then the morphism
$$
a \colon \Hom(x,y) \to \Hom(x',y')
$$
is injective (\resp bijective). Now $a$ is identified with the
canonical morphism
$$
\H^0(Y'_1,g^*_1 f_{1 *} (Q \overset{F_1}{\wedge} P^\circ)) \to
\H^0(Y'_1, f'_{1 *}(Q' \overset{F'_1}{\wedge} {P'}^\circ))\text{.}
$$
If
% original p. 351
we have $\Hom(x,y) \neq \emptyset$, then $Q \overset{F_1}{\wedge}
P^\circ$ is a torsor under $\leftexp{P\mkern-4mu}{F}_1$ which is locally trivial over~$Y_1$; it follows that $f_{1 *}(Q \overset{F_1}{\wedge}
P^\circ)$ is a torsor under $f_{1*}(\leftexp{P\mkern-4mu}{F}_1)$, and that
$g^*_1 f_{1 *} (Q \overset{F_1}{\wedge} P^\circ)$ is a trivial
torsor. The morphism $a$ is then identified with the canonical morphism
$$
\H^0(Y_1, g^*_1 f_{1 *}(\leftexp{P\mkern-4mu}{F}_1)) \to \H^0 (Y'_1,
f'_{1 *}(\leftexp{P'\mkern-6mu}{F}'_1)) \text{.}
$$
The same holds if we have $\Hom(x',y') \neq \emptyset$ and if
$a_1$ is injective, for then $Q' \overset{F'_1}{\wedge} {P'}^\circ$
is trivial, and it follows from the injectivity of~$a_1$ that $P$
and $Q$ are locally isomorphic over~$Y_1$. We conclude that, if
$a_0$ is injective (\resp if $a_0$ is bijective and $a_1$ injective),
$\varphi$ is faithful (\resp fully faithful).

It remains to show that, if $a_0$ and $a_1$ are bijective, the
functor $\varphi$ is essentially surjective. Let $Y''$ be a
scheme \'etale over~$Y'$, $X'' = X' \times_{Y'} Y''$, and
let $P''$ be a torsor on~$X''$ with group $F'' = F'|X''$. We shall
show that there exists an element~$x$ of~$(g^* f_*
\Phi)_{Y''}$ whose image in $(f'_* \Phi')_{Y''}$ is isomorphic
to $P''$. Let $p''$ be the class of~$P''$. From the fact that $a_1$ is
surjective it follows that one can find a surjective \'etale
morphism $Y''_1 \to Y''$, an \'etale morphism $Y_1 \to Y$ such that
there is a morphism $Y''_1 \to Y'_1$, and a torsor $P_1$ on~$X_1$
with group $F_1$ whose inverse image $P''_1$ on~$X''_1$ is
isomorphic to the inverse image of~$P''$. Using the fact that
$\varphi$ is fully faithful, one sees that the object $x_1$ of
$(g^* f_* \Phi)_{Y''_1}$ which corresponds to $P''_1$ is endowed
with a descent datum relative to $Y''_1 \to Y''$, hence
comes from an element $x$ of~$(g^* f_*
\Phi)_{Y''}$. Since the image of~$x$ in $(f'_* \Phi')_{Y''}$ is
$P''$, this proves that $\varphi$ is essentially surjective and
completes the proof.

\begin{example}
\label{XIII.1.4}
Let $f \colon X \to Y$ be a proper morphism. It follows from \cite[VII
2.2.2]{XIII.2} that, for every ind-finite stack $F$ on~$X$, the pair $(F,f)$
is cohomologically proper (relative to $Y$) in dimension
$\leqslant 1$. In particular, for every sheaf of sets
(\resp every sheaf of groups, \resp every ind-finite sheaf of
groups) $F$ on~$X$, $(F,f)$ is cohomologically
% original p. 352
proper in dimension $\leqslant 0$ (\resp in dimension $\leqslant 0$,
\resp in dimension $\leqslant 1$).
\end{example}

\begin{remarks}
\label{XIII.1.5}
a) Let $F$ be a sheaf of groups on~$X$ such that $(F,f)$ is
cohomologically proper relative to~$S$ in dimension
$\leqslant -1$ (\resp $\leqslant 0$). If one regards $F$ as a
sheaf of sets, $(F,f)$ is cohomologically proper
relative to~$S$ in dimension $\leqslant -1$ (\resp $\leqslant
0$), but the converse is false.

For example, let $Y$ be the spectrum of a strictly local discrete
valuation ring with closed point $t$ and generic point
$s$, $f \colon X \to Y$ a nonempty scheme over~$Y$ whose closed
fiber is empty, $F$ a nontrivial constant sheaf of groups
on~$X$, and $P$ a torsor under $F$ such that we have $\H^0(X_s,
\leftexp{P\mkern-4mu}{F}|X_s)=1$. Then $(\leftexp{P\mkern-4mu}{F},f)$ is
cohomologically proper relative to $Y$ in dimension
$\leqslant -1$ when one regards $\leftexp{P\mkern-4mu}{F}$ as a
sheaf of sets. If one regards $\leftexp{P\mkern-4mu}{F}$ as a sheaf
of groups, one has an isomorphism ${\vphantom{F}}^{P^\circ\mkern-7mu}(\leftexp{P\mkern-4mu}{F})
\simeq F$; since the canonical morphism
$$
\H^0(X,F)\to\H^0(X_t,F|X_t)=1
$$
is not injective, this proves that $(\leftexp{P\mkern-4mu}{F},f)$ is not
cohomologically proper relative to $Y$ in dimension
$\leqslant -1$.


b) Suppose $f$ coherent (\ie quasi-compact and
quasi-separated). Let $F$ be a stack on~$X$. For every
geometric point $\overline{y}$ of~$Y'$, we denote by $\overline{Y}$
(\resp $\overline{Y'}$) the strict localization of~$Y$ (\resp $Y'$) at
$\overline{y}$, and we set $\overline{X} = X \times_Y \overline{Y}$,
$\overline{X'} = X' \times_{Y'} \overline{Y'}$, etc. In order that $(F,f)$
be cohomologically proper relative to~$S$ in dimension
$\leqslant -1$ (\resp $\leqslant 0$, \resp $\leqslant 1$), it is necessary and
sufficient that, for every $S$-scheme $S'$ and for every geometric
point $\overline{y}$ of~$Y'$, the canonical functor
$$
\overline{F}(\overline{X})\to\overline{F'}(\overline{X'})
$$
be faithful (\resp fully faithful, \resp an
equivalence).

Indeed,
% original p. 353
if $S'$ is an $S$-scheme, in order that the functor
$$
g^* f_* F \to f'_* F'
$$
be faithful (\resp fully faithful, \resp an
equivalence), it is necessary and sufficient that the same hold for the functor
induced on the fibers at the various geometric points
$\overline{y}'$ of~$\overline{Y'}$ \cite[III 2.1.5.9]{XIII.2}. The assertion
therefore follows from the computation of the geometric fibers of the direct
image of a stack by a coherent morphism \cite[VII 2.1.5]{XIII.2}.


c) Let $F$ be a stack on~$X$. The fact that $(F,f)$ is
cohomologically proper relative to~$S$ in dimension
$\leqslant -1$ (\resp $\leqslant 0$, \resp $\leqslant 1$) is local
on~$Y$ for the \'etale topology.

Let $S'$ be an $S$-scheme, and $F'$ the inverse image of~$F$ on~$X'$
(\cf \eqref{eq:XIII.1.0.1}). If $(F,f)$ is cohomologically proper
relative to~$S$ in dimension $\leqslant 1$, the same holds
for~$(F', f')$. But if $(F,f)$ is cohomologically proper
relative to~$S$ in dimension $\leqslant -1$ (\resp $\leqslant
0$), the same does not necessarily hold for~$(F', f')$.

For example, let $S'$ be a discrete valuation ring, $f' \colon
E_{S'} \to S'$ the affine space over~$S'$, $x$ a closed
point of~$E_{S'}$ lying over the generic point of~$S'$, and
$F'$ the sheaf of sets on~$E_{S'}$ whose restriction to
$E_{S'}-\{x\}$ is the constant sheaf with one element and
whose fiber at a geometric point over~$x$ has two
elements. Then $(F',f')$ is not cohomologically proper
relative to~$S'$ in dimension $\leqslant -1$. Let $S =
S'[Z]$, $f \colon E_S \to S$ the affine space over~$S$, and $T$ a
closed subset of~$X = E_S$ which does not meet the closed set $Z=0$ and
such that $f(T)$ contains the generic point of~$S$. Let
$G$ be the inverse image of~$F'$ on~$X$ and $F$ the sheaf on~$X$ obtained
by extending $G|X-T$ by the empty set. Then $(F,f)$ is
cohomologically proper relative to~$S$ in dimension
$\leqslant -1$, but this is no longer so after the
change of base $S' \to S$ defined by $Z=0$.


d)
% original p. 354
Let $F$ be a stack on~$X$ such that $(F,f)$ is cohomologically
proper relative to $Y$ in dimension $\leqslant -1$ (\resp $\leqslant 0$, \resp $\leqslant 1$). Then, for every
geometric point $\overline{y}$ of~$Y$, the canonical functor
$$
(f_* F)_{\overline{y}} \to F(X_{\overline{y}})
$$
is faithful (\resp fully faithful, \resp an equivalence
of categories).
\end{remarks}

\begin{proposition}
\label{XIII.1.6}
Let $f \colon X \to Y$ and $g \colon Y \to Z$ be two $S$-morphisms, and
$\Phi$ a stack on~$X$.
\begin{enumerate}
\item[1)] Suppose that $(\Phi,f)$ and $(f_* \Phi, g)$ are
cohomologically proper relative to~$S$ in dimension
$\leqslant -1$ (\resp $\leqslant 0$, \resp $\leqslant 1$). Then the
same holds for~$(\Phi, gf)$.
\item[2)] Suppose that $(\Phi, gf)$ is cohomologically proper
relative to~$S$ in dimension $\leqslant -1$ (\resp that $(\Phi,
gf)$ is cohomologically proper relative to~$S$ in dimension
$\leqslant 0$ and $(\Phi,f)$ cohomologically proper relative
to~$S$ in dimension $\leqslant -1$, \resp that $(\Phi, gf)$ is
cohomologically proper relative to~$S$ in dimension
$\leqslant 1$ and $(\Phi, f)$ cohomologically proper relative
to~$S$ in dimension $\leqslant 0$). Then $(f_* \Phi, g)$ is
cohomologically proper relative to~$S$ in dimension
$\leqslant -1$ (\resp in dimension $\leqslant 0$, \resp in dimension
$\leqslant 1$).
\end{enumerate}
\end{proposition}

For every $S$-scheme $S'$, we consider the following diagram,
all of whose squares are cartesian:
\begin{equation*}
\label{eq:XIII.1.6.1}
\tag{\thesubsection.1}
\begin{array}{c}
\xymatrix{ X' \ar[r]^{f'} \ar[d]_h & Y'
\ar[r]^{g'} \ar[d]_k & Z' \ar[r] \ar[d]_{m} & S' \ar[d] \\ X \ar[r]^f
& Y \ar[r]^g & Z \ar[r] & \,S. }
\end{array}
\end{equation*}
The
canonical morphism
$$
m^* (g_* f_* \Phi) \to g'_* f'_* (h^* \Phi)
$$
is identified
% original p. 355
with the composite of the canonical morphisms
$$
m^*(g_* f_* \Phi) \lto{i} g'_*(k^*f_* \Phi) \lto{j} g'_* f'_*(h^* \Phi).
$$
\begin{enumerate}
\item[1)] The hypothesis implies that $i$ and $j$ are
faithful (\resp fully faithful, \resp
equi\-va\-lences); hence the same holds for~$ji$.
\item[2)] The hypothesis implies that $ji$ is faithful
(\resp that $ji$ is
fully
faithful and $j$ faithful, \resp that $ji$ is an equivalence
and $j$ fully faithful); it follows that $i$ is
faithful (\resp fully faithful, \resp an equivalence).
\end{enumerate}

\begin{corollary}
\label{XIII.1.7}
Let $f \colon X \to Y$ and $g \colon Y \to Z$ be two $S$-morphisms,
and let $F$ be a sheaf of groups on~$X$. Suppose that $(F,gf)$
is cohomologically proper relative to~$S$ in dimension
$\leqslant -1$ (\resp that $(F,gf)$ is cohomologically proper
relative to~$S$ in dimension $\leqslant 0$ and that $(F,f)$ is
cohomologically proper relative to~$S$ in dimension
$\leqslant -1$). Then $(f_* F, g)$ is cohomologically proper
relative to~$S$ in dimension $\leqslant -1$ (\resp in dimension
$\leqslant 0$).
\end{corollary}

Let us take up again the notation of \eqref{eq:XIII.1.6.1} and, for every
scheme $Y_1$ \'etale over~$Y$, let us denote by $f_1$, $F_1$ the
respective inverse images of~$f$, $F$ by the morphism $Y_1 \to
Y$. Let $\Phi$ be the stack of torsors under $F$ and $\Psi$ the stack
of torsors under $f_* F$. We have a canonical functor
$$
\varphi \colon \Psi \to f_* \Phi \text{,}
$$
obtained by associating to every scheme $Y_1$ \'etale over~$Y$ and
to every torsor $P$ on~$Y_1$ with group $f_{1 *}F_1$ the torsor
$\tilde{P}$ on~$X_1$ deduced from~$f_1^* P$ by extension of the
structure group $f_1^* f_{1 *}F_1 \to F_1$. The functor
$\varphi$ is fully faithful. Indeed, if $P$ and $Q$ are two
torsors on~$Y_1$ with group $f_{1 *} F_1$, we have a canonical
morphism
$$
\SheafIsom_{f_{1 *}F_1}(P,Q) \to
f_{1*}(\SheafIsom_{F_1}(\tilde{P},\tilde{Q}))
$$
which
% original p. 356
is an isomorphism since this is so locally. We deduce from this
that the canonical morphism
$$
\Isom_{f_{1 *}F_1}(P,Q) \to \Isom_{F_1}(\tilde{P},\tilde{Q})
$$
is an isomorphism, hence that $\varphi$ is fully faithful.

We have a commutative diagram
$$
\xymatrix{ g'_* k^* \Psi \ar[r]^i \ar[d]_{g'_*
k^*(\varphi)} & m^*(g_* \Psi) \ar[d]^{m^*
g_*(\varphi)} \\ g'_* k^*(f_* \Phi) \ar[r]^j &
m^*(g_* f_* \Phi)\text{,} }
$$
where $i$ and $j$ are the change of base morphisms. It
follows from~\Ref{XIII.1.6} 2) that $j$ is faithful
(\resp fully faithful). Since $g'_* k^* (\varphi)$ and
$m^* g_* (\varphi)$ are fully faithful, we deduce
from the above diagram that $i$ is faithful (\resp fully
faithful).

\begin{corollary}
\label{XIII.1.8}
Let $f \colon X \to Y$ be a coherent $S$-morphism, $g \colon Y
\to Z$ a proper $S$-morphism, and $\Phi$ an ind-finite stack on~$X$ \textup{\cite[VII 2.2.1]{XIII.2}}. Suppose that $(\Phi,f)$ is cohomologically proper
relative to~$S$ in dimension $\leqslant -1$ (\resp $\leqslant
0$, \resp $\leqslant 1$); then the same holds for~$(\Phi, gf)$.
\end{corollary}

Since $f$ is coherent, $f_* \Phi$ is an ind-finite stack (SGA~4 IX
1.6~(ii)). The corollary therefore follows from \Ref{XIII.1.6}~1)
and~\Ref{XIII.1.4}.

\begin{corollary}
\label{XIII.1.9}
Let $f \colon X \to Y$ be an integral $S$-morphism, and $g \colon Y \to Z$
an $S$-morphism. If $F$ is a sheaf of sets on~$X$, in order that
$(f_* F, g)$ be cohomologically proper relative to~$S$
in dimension $\leqslant -1$ (\resp $\leqslant 0$), it is necessary and
sufficient that the same hold for~$(F, gf)$. If $F$ is a sheaf of
groups on~$X$, in order that $(f_* F, g)$ be cohomologically
proper relative to~$S$ in dimension $\leqslant -1$
(\resp $\leqslant 0$, \resp $\leqslant 1$), it is necessary and sufficient that
the same hold for~$(F, gf)$.
\end{corollary}

The assertion
% original p. 357
relating to the case of a sheaf of sets follows
from~\Ref{XIII.1.6} and from the fact that $(F,f)$ is cohomologically proper
relative to~$S$ in dimension $\leqslant 0$. Let $F$ be a
sheaf of groups on~$X$ and $\Phi$ the stack of torsors under $F$.
By SGA~4 VIII~5.8, every torsor under $F$ is locally
trivial over~$Y$. It follows that the stack $f_* \Phi$ is
equivalent to the stack of torsors under $f_* F$, the equivalence
being obtained by associating to every scheme $Y_1$ \'etale
over~$Y$ and to every torsor $P$ on~$X_1 = X \times_Y Y_1$ with group
$F|X_1$ the torsor $f_* P$ with group $f_* F | Y_1$. Since $(F,f)$ is
cohomologically proper relative to~$S$ in dimension
$\leqslant 1$, the corollary therefore follows from~\Ref{XIII.1.6}.

\begin{definitions}
\label{XIII.1.10}
\leavevmode
\subsubsection{}
\label{XIII.1.10.1}
Let $E$ be a category and consider a diagram
$$
\xymatrix@C=.6cm{ \Phi \ar[r]^p & \Phi_1 \ar@<2pt>[r]^{p_1}
\ar@<-2pt>[r]_{p_2} & \Phi_2 } \text{,}
$$
where $\Phi$, $\Phi_1$, $\Phi_2$ are fibered categories
over~$E$ and the arrows are morphisms of fibered
categories, and let
$$
a \colon p_1 p \isomto p_2 p
$$
be an isomorphism of functors.

One says that the above diagram is exact if the following condition
is satisfied

a) For every pair of objects $x$, $y$ of~$\Phi$ and every morphism $f
\colon p(x) \to p(y)$ such that $p_1(f) = p_2(f)$ ($p_1p$ and
$p_2p$ being identified by means of $a$), there exists a
unique morphism $g \colon x \to y$ such that $p(g)=f$.

\subsubsection{}
\label{XIII.1.10.2}
Consider the diagram
$$
\xymatrix@C=.5cm{ \Phi \ar[r]^p & \Phi_1 \ar@<2pt>[rr]^{p_1,p_2}
\ar@<-2pt>[rr] && \Phi_2 \ar@<4pt>[rrr]^{p_{12},p_{23},p_{31}} \ar[rrr]
\ar@<-4pt>[rrr] &&& \Phi_3 } \text{,}
$$
where
% original p. 358
$\Phi$, $\Phi_i$, $1 \leqslant i \leqslant 3$, are fibered
categories over~$E$ and the arrows are morphisms of fibered
categories. Suppose given isomorphisms of
functors
$$
a \colon p_1 p \isomto p_2 p
$$
$$
a_1 \colon p_{31} p_2 \isomto p_{12} p_1 \text{,}\quad a_2 \colon p_{12}
p_2 \isomto p_{23} p_1 \text{,}\quad a_3 \colon p_{23} p_2 \isomto
p_{31}p_1
$$
such that the following diagram is commutative:
$$
\xymatrix@C=1.2cm{ p_{23} p_1 p \ar[r]^{\id.a} & p_{23} p_2 p \ar[r]^{a_3.\id}
& p_{31} p_1 p \ar[d]^{\id.a} \\ p_{12} p_2 p \ar[u]^{a_2.\id} &
p_{12} p_1 p \ar[l]^{\id.a} & p_{31} p_2 p \ar[l]^{a_1.\id} } \text{.}
$$
One identifies $p_1 p$ and $p_2 p$, $p_{31} p_2$ and $p_{12} p_1$, etc.

One says that the above diagram is exact if the following
conditions are satisfied:
\begin{enumerate}
\item[a)] Analogous to condition a) of~\Ref{XIII.1.10.1}.
\item[b)] For every object $x_1$ of~$\Phi_1$ and every isomorphism
$u \colon p_1(x_1) \isomto p_2(x_1)$ such that
\begin{equation*}
\label{eq:XIII.1.10.2.1}
\tag{\thesubsubsection.1} { p_{23}(u) p_{31}(u) = p_{12}(u)^{-1}
\text{,} }
\end{equation*}
there exists an object $x$ of~$\Phi$ such that there is an isomorphism $i
\colon p(x) \isomto x_1$ making the diagram
\begin{equation*}
\label{eq:XIII.1.10.2.2}
\tag{\thesubsubsection.2}
\begin{array}{c}
{ \xymatrix{ p_1p(x) \ar@{=}[r]
\ar[d]_{p_1(i)} & p_2p(x) \ar[d]^{p_2(i)} \\ p_1(x_1) \ar[r]^{\sim}_u
& p_2(x_1)
} }
\end{array}
\end{equation*}
commutative.
\end{enumerate}

\subsubsection{}
\label{XIII.1.10.3}
One defines in the obvious way the notion of a morphism of
exact diagrams of fibered categories over a
category $E$.
\end{definitions}

\subsubsection{}
\label{XIII.1.10.4}
We shall use
% original p. 359
in particular the notion of exact diagram in the
case where $E$ is a site and $\Phi$, $\Phi_i$, $1 \leqslant i
\leqslant 3$, are stacks over~$E$.

Let $f \colon E \to E'$ be a morphism of sites and let
\begin{equation*}
\label{eq:XIII.1.10.4.1}
\tag{\thesubsubsection.1} { \xymatrix@C=.5cm{ \Phi \ar[r]^p & \Phi_1
\ar@<2pt>[rr]^{p_1,p_2} \ar@<-2pt>[rr] && \Phi_2
\ar@<4pt>[rrr]^{p_{12},p_{23},p_{31}} \ar[rrr] \ar@<-4pt>[rrr] &&&
\Phi_3 } }
\end{equation*}
be an exact diagram of stacks over~$E$. By direct image one obtains a
diagram
$$
\xymatrix@C=.5cm{ f_* \Phi \ar[r] & f_* \Phi_1 \ar@<2pt>[r]
\ar@<-2pt>[r] & f_* \Phi_2 \ar@<4pt>[r]\ar[r] \ar@<-4pt>[r] &
f_* \Phi_3 }
$$
which is obviously exact.

If $h \colon E'' \to E$ is a morphism of sites, one likewise has an
exact diagram
$$
\xymatrix@C=.5cm{ h^* \Phi \ar[r]^{p''} & h^* \Phi_1
\ar@<2pt>[rr]^{p''_1,p''_2} \ar@<-2pt>[rr] && h^* \Phi_2
\ar@<4pt>[rrr]^{p''_{12},p''_{23},p''_{31}} \ar[rrr] \ar@<-4pt>[rrr] &&&
h^* \Phi_3 }
$$

Let us first verify condition~a) of~\Ref{XIII.1.10.2}. Let
$F''$ be an object of~$E''$, $x''$ and $y''$ two objects of~$(h^*
\Phi)_{F''}$, $x''_1$ and $y''_1$ their respective images in
$h^* \Phi_1$, and $x''_2$, $y''_2$ their images in $h^*
\Phi_2$. Let $u''_1 \colon x''_1 \to y''_1$ be a morphism such that
$p''_1(u''_1) = p''_2(u''_1)$, and let us prove that $u''_1$ comes
from a unique morphism $u'' \colon x'' \to y''$. The question being
local over~$F''$, one may assume that one has an object $F_1$ of~$E$, a
morphism from~$F''$ into the inverse image $F''_1$ of~$F_1$ by $h$, and
objects $x$, $y$ of~$\Phi_{F_1}$ whose inverse images over~$F''$
are $x''$ and $y''$. Let $x_1$, $y_1$ (\resp $x_2$, $y_2$) be the
images of~$x$, $y$ in $\Phi_1$ (\resp~$\Phi_2$). One may assume
that $u''_1$ comes from a morphism $u_1 \colon x_1 \to y_1$ such that
$p_1(u_1) = p_2(u_1)$. In view of the exactness of
\eqref{eq:XIII.1.10.4.1}, one obtains a unique morphism $u \colon x
\to y$ whose inverse image by $h$ is the required morphism
$u''$.

Condition~b)
% original p. 360
of~\Ref{XIII.1.10.2} is verified in an analogous
way. Let $x''_1$ be an object of~$(h^* \Phi_1)_{F''}$ and $u''
\colon p''_1(x''_1) \to p''_2(x''_1)$ a morphism satisfying
the relation
$$
p''_{23}(u'')p''_{31}(u'')=p''_{12}(u'')^{-1} \text{,}
$$
and let us prove that there exist an object $x''$ of~$(h^* \Phi)_{F''}$ and
an isomorphism $i'' \colon p''(x'') \simeq x''_1$ making commutative
a diagram analogous to \eqref{eq:XIII.1.10.2.2}. Since the
question is local over~$F''$, one may assume that one has an object
$F_1$, a morphism $F'' \to F''_1$ as above, and an object $x_1$
of~$(\Phi_1)_{F_1}$ whose inverse image in $(h^* \Phi_1)_{F''}$
is $x''_1$. Likewise one may assume that $u''$ comes from a
morphism $u \colon p_1(x_1) \to p_2(x_1)$ satisfying
\eqref{eq:XIII.1.10.2.2}. The existence of an object $x$ of~$\Phi_{F_1}$
whose inverse image by $h$ is an element $x''$
answering the question then follows from the exactness of
\eqref{eq:XIII.1.10.4.1}.

\begin{examples}
\label{XIII.1.11}
\begin{enumerate}
\item[1)] Let $f \colon X_1 \to X$ be a \emph{descent morphism}
for the category of \'etale sheaves over variable
schemes (for example a universally submersive morphism (SGA~4
VIII~9.3)). Let $X_2 = X_1 \times_X X_1$, $g \colon X_2 \to X$ the
canonical projection, and $F$ a sheaf of sets over~$X$. It then
follows from \loccit that one has an exact sequence of
sheaves of sets
\begin{equation*}
\label{eq:XIII.1.11.1}
\tag{\thesubsection.1} { \xymatrix@C=.5cm{ F \ar[r] & f_* f^* F
\ar@<2pt>[r] \ar@<-2pt>[r] & g_* g^* F \text{.}}
}
\end{equation*}
If $\Phi$ is the stack in discrete categories associated
to $F$ and $\Phi_3$ the final stack over~$X$, \ie the stack all of whose
fibers are reduced to a single element
having as its only morphism the identity morphism, to say that the sequence
\eqref{eq:XIII.1.11.1} is exact amounts to saying that the same
holds for the diagram of stacks
$$
\xymatrix@C=.5cm{ \Phi \ar[r] & f_* f^* \Phi \ar@<2pt>[r]
\ar@<-2pt>[r] & g_* g^* \Phi \ar[r] \ar@<4pt>[r]
\ar@<-4pt>[r] & \Phi_3 } \text{.}
$$

\item[2)]
Let
% original p. 361
$f \colon X_1 \to X$ be an \emph{effective descent morphism}
for the category of \'etale sheaves over variable
schemes (for example a proper surjective morphism, or an integral
surjective morphism, or a faithfully flat morphism locally of
finite presentation (SGA~4 VIII~9.4)). Let $X_2 = X_1 \times_X
X_1$, $g \colon X_2 \to X$ the canonical projection, $X_3 = X_1
\times_X X_1 \times_X X_1$, $h \colon X_3 \to X$ the canonical
morphism. Let $\Phi$ be a stack over~$X$, $\Phi_1 = f_* f^* \Phi$,
$\Phi_2 = g_* g^* \Phi$, $\Phi_3 = h_* h^* \Phi$. One then has an
exact diagram
$$
\xymatrix@C=.5cm{ \Phi \ar[r] & \Phi_1 \ar@<2pt>[r] \ar@<-2pt>[r] &
\Phi_2 \ar[r] \ar@<4pt>[r] \ar@<-4pt>[r] & \Phi_3 \text{,} }
$$
where the arrows are the canonical morphisms associated with the
projections.

Indeed, consider $\Phi$ as a stack over the category
$\Sch_X$ of schemes over~$X$, endowed with the \'etale
topology. Then, by \cite[VII 2.2.8]{XIII.2}, $\Phi$ is also a
stack for the finest topology on~$\Sch_X$ for which the
covering morphisms are the effective descent morphisms for the
category of \'etale sheaves. The exactness of the above diagram
follows at once.
\end{enumerate}
\end{examples}

\begin{proposition}
\label{XIII.1.12}
Let $S$ be a scheme and $f\colon X \to Y$ an $S$-morphism.
\begin{enumerate}
\item[1)] Let
$$
\xymatrix@C=.5cm{ \Phi \ar[r] & \Phi_1 \ar@<2pt>[r] \ar@<-2pt>[r] &
\Phi_2 }
$$
be an exact diagram of stacks over~$X$. Suppose that $(\Phi_1, f)$ is
cohomologically proper relative to~$S$ in dimension
$\leqslant 0$ and that $(\Phi_2, f)$ is cohomologically proper relative to~$S$ in dimension $\leqslant -1$. Then $(\Phi,f)$
is cohomologically proper relative to~$S$ in dimension
$\leqslant 0$.

\item[2)] Let
$$
\xymatrix@C=.5cm{ \Phi \ar[r] & \Phi_1 \ar@<2pt>[r] \ar@<-2pt>[r] &
\Phi_2 \ar@<4pt>[r] \ar[r] \ar@<-4pt>[r] & \Phi_3 }
$$
be an exact diagram of stacks over~$X$. Suppose that $(\Phi,f)$ is
cohomologically proper relative to~$S$ in dimension
$\leqslant 1$, that $(\Phi_2, f)$ is cohomologically proper
% original p. 362
relative to~$S$ in dimension $\leqslant 0$, and that
$(\Phi_3, f)$ is cohomologically proper
relative
to~$S$ in dimension $\leqslant -1$. Then $(\Phi,f)$ is
cohomologically proper relative to~$S$ in dimension
$\leqslant 1$.
\end{enumerate}
\end{proposition}

For every $S$-scheme $S'$, one considers the following commutative
diagram, all of whose squares are cartesian:
$$
\xymatrix{ X' \ar[r]^{f'} \ar[d]_h & Y' \ar[r] \ar[d]_g & S' \ar[d] \\
X \ar[r]^f & Y \ar[r] & S }
$$

Let us prove 2), the proof of 1) being
analogous. Since the direct image and inverse image functors
transform exact diagrams of stacks into exact diagrams
\eqref{XIII.1.10.4}, one has the following morphism of exact diagrams of stacks:
$$
\xymatrix{ g^* f_* \Phi \ar[r]^\pi \ar[d]_{\varphi} & g^*
f_* \Phi_1 \ar@<2pt>[r]^{\pi_1} \ar@<-2pt>[r]_{\pi_2}
\ar[d]_{\varphi_1} & g^* f_* \Phi_2 \ar@<4pt>[r] \ar[r]
\ar@<-4pt>[r] \ar[d]_{\varphi_2} & g^* f_* \Phi_3
\ar[d]_{\varphi_3} \\ f'_* h^* \Phi \ar[r] & f'_* h^*
\Phi_1 \ar@<2pt>[r] \ar@<-2pt>[r] & f'_* h^* \Phi_2
\ar@<4pt>[r] \ar[r] \ar@<-4pt>[r] & f'_* h^* \Phi_3 }
\text{.}
$$
By hypothesis $\varphi_1$ is an equivalence of
categories, $\varphi_2$ is fully faithful and $\varphi_3$
faithful. It therefore follows from the preceding diagram that
$\varphi$ is an equivalence.

\begin{proposition}
\label{XIII.1.13}
Let $f \colon X \to Y$ be an $S$-morphism.
\begin{enumerate}
\item[1)] Let
$$
\xymatrix@C=.5cm{ F \ar[r] & G \ar@<2pt>[r] \ar@<-2pt>[r] & H }
$$
be an exact diagram of sheaves of sets over~$X$. Suppose that
$(G,f)$ is cohomologically proper relative to~$S$ in
dimension $\leqslant 0$ and that $(H,f)$ is cohomologically proper
relative to~$S$ in dimension $\leqslant -1$. Then $(F,f)$ is
cohomologically proper relative to~$S$ in dimension
$\leqslant 0$.
\item[2)]
Let
% original p. 363
$F \to G$ be a monomorphism of sheaves of groups over~$X$. If
$Y_1$ is a scheme \'etale over~$Y$, one sets $X_1 = Y_1 \times_Y
X$, and one denotes by $f_1$ (\resp $F_1$, \resp $G_1$) the inverse image of~$f$
(\resp $F$, \resp $G$) over~$Y_1$ (\cf \Ref{eq:XIII.1.0.2}). Suppose
that $(G,f)$ is cohomologically proper relative to~$S$ in
dimension $\leqslant 0$ (\resp in dimension $\leqslant 1$) and that,
for every scheme $Y_1$ \'etale over~$Y$ and every torsor $Q$
under $G_1$, $(Q/F_1, f_1)$ is cohomologically proper relative
to~$S$ in dimension $\leqslant -1$ (\resp in dimension $\leqslant
0$). Then $(F,f)$ is cohomologically proper relative to~$S$
in dimension $\leqslant 0$ (\resp in dimension $\leqslant 1$).
\item[3)] Let $F \to G$ be a monomorphism of sheaves of groups over~$X$. Suppose that $(F,f)$ is cohomologically proper relative
to~$S$ in dimension $\leqslant 1$ and that $(G,f)$ is
cohomologically proper relative to~$S$ in dimension
$\leqslant 0$. Then, for every torsor $Q$ under $G$, $(Q/F, f)$ is
cohomologically proper relative to~$S$ in dimension
$\leqslant 0$.
\end{enumerate}
\end{proposition}

\subsubsection*{Proof}
1) Let $\Phi$ (\resp $\Phi_1$, \resp $\Phi_2$) be the stack in
discrete categories associated to $F$ (\resp $G$,
\resp $H$) and let $\Phi_3$ be the final stack over~$X$. One then has an
exact diagram
$$
\xymatrix@C=.5cm{ \Phi \ar[r] & \Phi_1 \ar@<2pt>[r] \ar@<-2pt>[r] &
\Phi_2 \ar@<4pt>[r] \ar[r] \ar@<-4pt>[r] & \Phi_3 }
\text{.}
$$
By hypothesis $(\Phi_1, f)$ is cohomologically proper
relative to~$S$ in dimension $\leqslant 1$ and $(\Phi_2, f)$ is
cohomologically proper relative to~$S$ in dimension
$\leqslant 0$ \eqref{XIII.1.2}. Since $(\Phi_3, f)$ is obviously
cohomologically proper relative to~$S$ in dimension
$\leqslant -1$, it follows from~\Ref{XIII.1.12} that $(\Phi, f)$ is
cohomologically proper relative to~$S$ in dimension
$\leqslant 1$, \ie that $(F,f)$ is cohomologically proper
relative to~$S$ in dimension $\leqslant 0$.

\smallskip
2) Let us first show that, if $(G,f)$ is cohomologically
proper relative to~$S$ in dimension $\leqslant 0$ and if the
$(Q/F_1, f_1)$ are cohomologically
% original p. 364
proper relative to~$S$ in dimension $\leqslant -1$, then
$(F,f)$ is cohomologically proper relative to~$S$ in
dimension $\leqslant 0$. By~\Ref{XIII.1.3.1} it suffices to
prove that, for every scheme $Y_1$ \'etale over~$Y$ and
every torsor $P$ over~$X_1$ with group $F_1$, $(\leftexp{P\mkern-4mu}{F}_1,
f_1)$ is cohomologically proper relative to~$S$ in dimension
$\leqslant 0$ when one considers $\leftexp{P\mkern-4mu}{F}_1$ as a
sheaf of sets, and that the canonical morphism
$$
d \colon g^*(\R^1 f_* F) \to \R^1 f'_* F'
$$
is injective. The first assertion follows at once from 1),
since, if $Q$ denotes the torsor deduced from~$P$ by the extension
$F_1 \to G_1$ of the structure group, one has an isomorphism
$\leftexp{Q}{G}_1 / \leftexp{P\mkern-4mu}{F}_1 \isomto Q/F_1$.

Let us show that $d$ is injective. It suffices to prove that, if $Y_1$ is
a scheme \'etale over~$Y$ and if $P$ and $\tilde{P}$ are
two torsors under $F_1$ whose inverse images $P'$ and~$\tilde{P}'$
over~$X'_1$ are isomorphic, then, after possibly making a surjective
\'etale extension of~$Y_1$, $P$ and~$\tilde{P}$ become
isomorphic. Choose an isomorphism $p' \colon P' \isomto
\tilde{P}'$. Let $Q$ (\resp~$\tilde{Q}$) be the torsor deduced from
$P$ (\resp $\tilde{P}$) by the extension of the structure group $F_1 \to
G_1$. The inverse images $Q'$ (\resp $\tilde{Q}'$) of~$Q$
(\resp $\tilde{Q}$) over~$X'_1$ are deduced from~$P'$
(\resp $\tilde{P}'$) by
the extension
of the structure group $F'_1 \to G'_1$; let $q' \colon Q' \isomto
\tilde{Q}'$ be the isomorphism that one obtains in the same way from
$p'$. Since $(G,f)$ is cohomologically proper relative to~$S$ in dimension $\leqslant 0$, one may assume, after possibly making
a surjective \'etale extension of~$Y_1$, that $q'$ is the image of an
isomorphism $q \colon Q \isomto \tilde{Q}$. To the torsor~$P$
(\resp~$\tilde{P}$) is associated a section $x$ of~$Q/F_1$
(\resp a section $\tilde{x}$ of~$\tilde{Q}/F_1$), and, for $P$
and $\tilde{P}$ to be isomorphic, it is necessary and sufficient that there be an
isomorphism $Q \isomto \tilde{Q}$ such that the isomorphism
$$
e \colon \H^0(X_1, Q/F_1) \to \H^0(X_1,\tilde{Q}/F_1)
$$
deduced from it transforms $x$ into $\tilde{x}$. One takes
the isomorphism $q$. The sections $e(x)$ and $\tilde{x}$ of
$\H^0(X_1,\tilde{Q}/F_1)$ have the same image in
$\H^0(X'_1,\tilde{Q}'/F'_1)$. Since $(\tilde{Q}/F_1,f_1)$ is
cohomologically proper relative to~$S$ in dimension
$\leqslant -1$,
% original p. 365
after possibly making a surjective \'etale extension of~$Y_1$, one does
have $e(x) = \tilde{x}$, which proves the injectivity of~$d$.

To complete the proof, it remains to prove that, if
$(G,f)$ is cohomologically proper relative to~$S$ in
dimension $\leqslant 1$, and if, for every scheme $Y_1$ \'etale
over~$Y$ and every torsor $Q$ over~$X_1$ with group $F_1$,
$(Q/F_1, f_1)$ is cohomologically proper relative to~$S$ in
dimension $\leqslant 0$, then the morphism $d$ is surjective. Let
$P'$ be a torsor over~$X'_1$ with group $F'_1$, and $Q'$ the torsor under
$G'_1$ obtained from~$P'$ by extension of the structure
group. Giving~$P'$ is equivalent to giving
$Q'$ and a section $x'$ of~$\H^0(X'_1,Q'/F'_1)$. It then
follows from the surjectivity of the morphism
$$
g^*(\R^1 f_* G) \to \R^1 f'_* G'
$$
that, after possibly making a surjective \'etale extension of~$Y_1$,
there exists a torsor $Q$ under~$G_1$ whose inverse image over~$X'_1$
is isomorphic to $Q'$. Using the fact that $(Q/F_1, f_1)$ is
cohomologically proper relative to~$S$ in dimension
$\leqslant 0$, one may likewise assume that there exists an
element $x$ of~$\H^0(X_1, Q/F_1)$ whose image in
$\H^0(X'_1, Q'/F'_1)$ is $x'$. The data of~$Q$ and of~$x$
determine a torsor $P$ under $F_1$ whose inverse image over~$X_1$ is isomorphic to $P'$, which proves the
surjectivity of~$d$.

\smallskip
3) Let us show that $(Q/F,f)$ is cohomologically proper
relative to~$S$ in dimension $\leqslant -1$, \ie that, for
every $S$-scheme $S'$ and every scheme $Y_1$ \'etale
over~$Y$, if $x$, $\tilde{x}$ are two elements of
$\H^0(X_1,Q_1/F_1)$ whose images $x'$, $\tilde{x}'$ in
$\H^0(X'_1, Q'_1/F'_1)$ are equal, then, after a surjective
extension of~$Y_1$, one has $x = \tilde{x}$. To~$x$ (\resp~$\tilde{x}$)
is associated a torsor $P$ (\resp~$\tilde{P}$) under $F_1$ such that
$Q_1$ is deduced from~$P$ (\resp~$\tilde{P}$) by the extension $F_1
\to G_1$ of the structure group. From the relation $x' = \tilde{x}'$
it follows that one has an isomorphism $u' \colon P' \isomto
\tilde{P}'$ such that the isomorphism induced on~$Q'_1$ by~$u'$ is
the identity. Since $(F,f)$ is cohomologically proper
relative to~$S$ in dimension $\leqslant 0$, one deduces
that,
% original p. 366
after a surjective \'etale extension of~$Y_1$, one has an
isomorphism $u \colon P \to \tilde{P}$ lifting $u'$; the fact that
$(G,f)$ is cohomologically proper relative to~$S$ in
dimension $\leqslant -1$ then implies that $x =
\tilde{x}$.

Let us show that $(Q/F,f)$ is cohomologically proper relative to~$S$ in dimension $\leqslant 0$. Let $Y''$ be a scheme \'etale
over~$Y'$ and $x''$ an element of~$\H^0(X'',Q''/F'')$. To~$x''$
is associated a torsor $P''$ over~$X''$ with group $F''$. Since
$(F,f)$ is cohomologically proper relative to~$S$ in
dimension $\leqslant 1$, one can find surjective \'etale
morphisms $Y''_1 \to Y''$ and $Y_1 \to Y$ such that one has a
morphism $Y''_1 \to Y'_1$, and a torsor $P$ over~$X_1$ with group
$F_1$ whose inverse image over~$X''_1$ is isomorphic to the inverse
image of~$P''$. It then follows from the fact that $(G,f)$ is
cohomologically proper relative to~$S$ in dimension
$\leqslant 0$ that one can even choose $Y''_1$ and $Y_1$ such that
the torsor deduced from~$P$ by extension of the structure group $F_1
\to G_1$ is isomorphic to $Q_1$; to~$P$ there corresponds an
element $x$ of~$\H^0(X_1,Q_1/F_1)$ whose image in
$\H^0(X''_1,Q''_1/F''_1)$ is isomorphic to the inverse image of
$x''$, which completes the proof.

\begin{proposition}
\label{XIII.1.14}
Let $f \colon X \to S$ be an $S$-scheme and $F$ a sheaf
of sets or of groups over~$X$ (\resp a sheaf of
ind-$\LL$-groups, where $\LL$ is a set of prime
numbers). Suppose $F$ locally constant, $(F,f)$
cohomologically proper in dimension $\leqslant 0$ (\resp in
dimension $\leqslant 1$), and $f$ locally $0$-acyclic
(\resp locally $1$-aspherical for $\LL$) \textup{(SGA~4 XV~1.11)}.
Then, for every specialization $\overline{s}_1 \to
\overline{s}_2$ of geometric points of~$S$, the
specialization morphism \textup{(SGA~4 VIII~7.1)}
$$
a_0 \colon (f_* F)_{\overline{s}_2} \to (f_*
F)_{\overline{s}_1}
$$
is an isomorphism, and, if $F$ is a sheaf of groups, the
morphism
$$
a_1 \colon (\R^1 f_* F)_{\overline{s}_2} \to (\R^1 f_*
F)_{\overline{s}_1}
$$
is injective (\resp the morphisms $a_0$ and $a_1$ are
isomorphisms).
\end{proposition}

The
% original p. 367
proof is obtained by copying word for word that of SGA~4
XVI~2.3, but replacing in it the expression ``proper'' by
the expression ``cohomologically proper''.

\begin{corollary}
\label{XIII.1.15}
Let $f \colon X \to S$ be a morphism, $\Phi$ a stack over~$X$, and $\LL$
a set of prime numbers. Suppose that, for every scheme
$X_1$ \'etale over~$X$ and every pair of objects $x$, $y$ of
$\Phi_{X_1}$, the sheaf $\SheafHom_{X_1}(x,y)$ is locally
constant, that the sheaf $\SheafAut_{X_1}(x)$ is a locally constant
ind-$\LL$-group, and that the sheaf of
maximal subgerbes $S\Phi$ of~$\Phi$ \textup{\cite[III 2.1.7]{XIII.1}} is locally
constant. Suppose that $(\Phi,f)$ is cohomologically proper in
dimension $\leqslant 1$ and that $f$ is locally $1$-aspherical
for $\LL$. Then, for every specialization $\overline{s}_1 \to
\overline{s}_2$ of geometric points of~$S$, the
specialization morphism
$$
a \colon (f_* \Phi)_{\overline{s}_2} \to (f_*
\Phi)_{\overline{s}_1}
$$
is an equivalence of categories.
\end{corollary}

Let $\overline{S}_1$ (\resp $\overline{S}_2$) be the strict
localization of~$S$ at $\overline{s}_1$ (\resp the strict localization of~$S$
at $\overline{s}_2$), $\overline{X}_2$, $\overline{\Phi}_2$
(\resp $\overline{X}_1$, $\overline{\Phi}_1$) the inverse images of
$X_2$, $\Phi_2$ over~$\overline{S}_2$ (\resp of~$X_1$, $\Phi_1$ over~$\overline{S}_1$), and consider the cartesian square
$$
\xymatrix{ \overline{X}_1 \ar[r]^h \ar[d]_{\overline{f}_1} &
\overline{X}_2 \ar[d]^{\overline{f}_2} \\ \overline{S}_1 \ar[r]^g &
\overline{S}_2 }
\text{.}
$$
One must show that the functor
$$
\varphi \colon \overline{\Phi}_2(\overline{X}_2) \to \overline{\Phi}_1
(\overline{X}_1)
$$
is an equivalence. The functor $\varphi$ is fully
faithful; indeed, let
% original p. 368
$x$, $y$ be two objects of~$(\overline{\Phi}_2)_{\overline{X}_2}$;
the canonical morphism
$$
\Hom_{\overline{X}_2}(x,y) \to
\Hom_{\overline{X}_1}(\varphi(x),\varphi(y))
$$
is identified with the canonical morphism
$$
\H^0(\overline{X}_2, \SheafHom_{\overline{X}_2}(x,y)) \to
\H^0(\overline{X}_1,h^*(\SheafHom_{\overline{X}_2}(x,y)) \text{.}
$$
This morphism is an isomorphism by~\Ref{XIII.1.14}.

Let us show that $\varphi$ is an equivalence. Let $x_1$ be an object
of~$\overline{\Phi}_1(\overline{X}_1)$ and $G_1$ the maximal subgerbe
of~$\overline{\Phi}_1$ generated by $x_1$. The morphism
$$
\H^0(\overline{X}_2, S\overline{\Phi}_2) \to \H^0(\overline{X}_1,
h^*(S \overline{\Phi}_2)) =
\H^0(\overline{X}_1,S\overline{\Phi}_1)
$$
is bijective, and there therefore exists a maximal subgerbe $G_2$ of
$\overline{\Phi}_2$ such that $h^* G_2$ is isomorphic to
$G_1$. It then suffices to prove that the functor
$$
G_2 \to h_* h^* G_2
$$
is an equivalence. But, in this form, the question is
local for the \'etale topology on~$\overline{X}_2$. One may therefore
assume that $G_2$ is a gerbe of torsors under the group of
automorphisms of an object of~$G_2$, a case in which the assertion follows
from~\Ref{XIII.1.14}.

\begin{corollary}
\label{XIII.1.16}
The hypotheses are those of~\Ref{XIII.1.14}. If one assumes
moreover that $F$ is a sheaf of sets (\resp of ind-$\LL$-groups)
and that $f_* F$ (\resp $\R^1 f_* F$) is constructible, then
$f_* F$ (\resp $\R^1 f_* F$) is locally constant.
\end{corollary}

The corollary follows from~\Ref{XIII.1.14} by means of SGA~4
IX~2.11.

\begin{remark}
\label{XIII.1.17}
Let us recall that the condition $f$ locally $0$-acyclic is satisfied
if $f$ is flat with separable fibers, $X$ and $Y$ locally
noetherian (SGA~4 XV~4.1), and that the condition $f$ locally
$1$-aspherical
% original p. 369
for $\LL$ is satisfied if $f$ is smooth, $\LL$ being the set
of prime numbers distinct from the residue
characteristics of~$S$ (SGA~4 XV~2.1).
\end{remark}

\section[A particular case of cohomological properness]{A particular case of cohomological properness: divisors
with relative normal crossings}
\label{XIII.2}

\setcounter{subsection}{-1}

\subsection{}
\label{XIII.2.0}
Let $R$ be a discrete valuation ring with field of fractions
$K$ and $L$ an étale $K$-algebra; $L$ is then a direct product
of a finite number of fields $L_i$, where $L_i$ is an étale
extension of~$K$. If $L'_i$ denotes the Galois extension
generated by $L_i$ in an algebraic closure of~$L_i$,
one says that $L$ is \emph{tamely ramified}
over~$R$ if the $L'_i$ are tamely
ramified extensions in the sense of~X~\Ref{X.3}, \ie if an inertia group~$I_i$ of~$L'_i|K$ is of order prime to the residue
characteristic $p$ of~$R$.

We know that $I_i$ is in any case an extension of a cyclic group
of order prime to $p$ by a $p$-group. (This follows from \cite[ch\ptbl IV prop\ptbl 7 cor\ptbl 4]{XIII.5} when the residue extension of~$R$ is
separable. The proof given in
\loccit extends to the general case in the following
way. Let us take up again the hypotheses and the notation of
\loccit but without assuming the residue extension
separable. Let $H_i$ be the subgroup of the inertia group $G_0$,
the set of the elements $s$ of~$G_0$ such that one has $s \pi /
\pi \in U^i$ for every uniformizer $\pi$ of~$A_L$. One then checks
that $G_0/H_1$ is a group of order prime to $p$ and that,
for $i \geqslant 1$, the $H^i/H^{i+1}$ are $p$-groups, from which
one deduces the announced result.)

\subsubsection{}
\label{XIII.2.0.1}
In the case where $R$ is strictly local, one has the following
simple characterization: the $K$-algebra $L$ is
tamely ramified over~$R$ if and only if the
$[L_i:K]$ are prime to $p$. Moreover, if $L$ is
tamely ramified over~$R$, the $L_i$ are cyclic extensions
of~$K$. Indeed, when $R$ is strictly local, $I_i$
is equal to the Galois group
% original p. 370
of~$L'_i$ over~$K$. As we have just recalled, $I_i$ is an extension
of a cyclic group of order prime to $p$ by a $p$-group. If
one assumes $L'_i$ tamely ramified over~$R$, $I_i$ is
then a cyclic group of order prime to $p$. It follows
that $[L_i:K]$ is prime to $p$ and that one has $L_i =
L'_i$. Conversely, if $[L_i:K]$ is prime to $p$, $I_i$ cannot
contain a nontrivial normal $p$-subgroup; $I_i$ is therefore
a cyclic group of order prime to $p$, which proves that $L_i$
is tamely ramified over~$R$.

\subsubsection{}
\label{XIII.2.0.2}
Let $R$ be a discrete valuation ring with field of fractions
$K$, $L$ an étale $K$-algebra, and let $\overline{R}$ be a
strict localization of~$R$, $\overline{K}$ its field of fractions,
$\overline{L} = L \otimes_K \overline{K}$. Then, for $L$ to be
tamely ramified over~$R$, it is necessary and sufficient that
$\overline{L}$ be tamely ramified over~$\overline{R}$. Indeed one reduces to the case where $L$ is a
field. Let then $\overline{L} = \prod_i \overline{L}_i$, where
the $\overline{L}_i$ are fields, extensions of~$\overline{K}$; if
$L'$ is the Galois extension generated by $L$, and, if
$\overline{L'} = L' \otimes_K \overline{K}$,
one likewise has a decomposition of~$\overline{L'}$ into a product
of fields,
$\overline{L'} = \prod_j \overline{L'_j}$
and each $\overline{L}_i$ is a subextension of at least one of the
$\overline{L'_j}$. Since $L'$ is a Galois extension of~$K$,
the $\overline{L'_j}$ are Galois extensions of
$\overline{K}$. Suppose $L$ tamely ramified over~$R$;
since the Galois group of~$\overline{L'_j}|K$ is isomorphic to the
inertia group of~$L'|K$, the $L'_j$ are also tamely
ramified over~$\overline{R}$, and hence so are the
$\overline{L}_i$. Conversely, suppose $\overline{L}$
tamely ramified over~$\overline{R}$. For each $j$,
let $v_j$ be the discrete valuation of~$\overline{L'_j}$ which extends
the valuation of~$\overline{K}$, and let us still denote by $v_j$ the valuation
induced on~$L'$. As $j$ varies, $v_j$ runs through the set of
valuations of~$L$ which extend the valuation of~$K$. Let $G =
\Gal(L'|K)$, $H=\Gal(L'|L)$, $I_j$ the inertia group of~$L'|K$ at
$v_j$, $J_j$ the inertia group of~$L'|L$ at $v_j$. The group $I_i$
is an extension of a cyclic group of order prime to $p$ by a
$p$-group $P_j$. Since the $\overline{L}_i$ are tamely
ramified over~$\overline{R}$, $I_j/J_j$ is of order prime to
$p$, hence one has $P_j \subset J_j$. Consequently the group $H$ contains
all the $P_j$, hence also the group $P$ generated by
% original p. 371
the $P_j$ for variable $j$. But the group~$P$ is invariant in $G$,
since an inner automorphism of~$G$ transforms the~$I_j$ into one
another, hence also the $P_j$ into one another. It follows that $P$ is a
subgroup of~$H$ normal in~$G$, hence, since $L'$ is
the Galois extension generated by $L$, that one has $P=1$, which
proves that $L$ is tamely ramified over~$R$.

More generally, let $R \to R'$ be a morphism of discrete
valuation rings such that the image of a uniformizer $\pi$ of
$R$ is a uniformizer $\pi'$ of~$R'$ and such that the residue
extension $\kres(R')$ is a separable extension of
$\kres(R)$. Let $K$ be the field of fractions of~$R$, $K'$ the field of
fractions of~$R'$, $L$ an étale $K$-algebra, $L' = L
\otimes_K K'$. Then, for $L$ to be tamely
ramified over~$R$, it is necessary and sufficient that $L'$ be
tamely ramified over~$R'$. Indeed we may assume
$R$ and $R'$ strictly local. By~\Ref{XIII.2.0.1} it suffices
to prove that, when $L$ is a field, so is~$L'$. Let $\tilde{R}$ be the normalization of~$R$ in $L$,
$\tilde{\pi}$ a uniformizer of~$\tilde{R}$, $\tilde{R}' =
\tilde{R} \otimes_R R'$. The extension $\kres(\tilde{R})|\kres(R)$ being
radicial and the extension $\kres(R')|\kres(R)$ étale,
$\kres(\tilde{R})\otimes_{\kres(R)}\kres(R')$ is a field [EGA IV 4.3.2 and
4.3.5]. This proves that $R'$ is a local ring, and, since $\pi$ has
image $\pi'$ in $R'$, one has $\kres(\tilde{R}')= R'/(\tilde{\pi})$;
consequently $R'$ is a discrete valuation ring \cite[ch\ptbl I \S 2
prop\ptbl 2]{XIII.5}, hence $L'$ is a field.

\subsubsection{}
\label{XIII.2.0.3}
By reduction to the strictly local case, one sees that a
subalgebra of a tamely ramified algebra is
tamely ramified, that the tensor product of two
tamely ramified algebras is tamely
ramified, that a tamely ramified algebra remains
so after extension of the discrete valuation ring,
that an algebra which becomes tamely ramified
after a tamely ramified extension is
tamely ramified.

\subsection{}
\label{XIII.2.1}
Let
% original p. 372
$X$ be an $S$-scheme, $D$ a divisor $\geq 0$ on~$X$.
Recall (SGA~5 II~4.2) that one says that $D$ has strictly
\emph{normal crossings relative to}~$S$
if there exists a finite family $(f_i)_{i\in I}$ of elements of
$\Gamma(X,\cal{O}_X)$ such that one has $D=\sum_{i\in
I}\divisor(f_i)$ and such that the following condition is satisfied:

\setcounter{subsubsection}{-1}
\subsubsection{}
\label{XIII.2.1.0}
For every point $x$ of~$\Supp D$, $X$ is smooth over~$S$ at $x$, and, if
one denotes by $I(x)$ the set of the $i\in I$ such that $f_i(x)=0$, the
subscheme $V((f_i)_{i\in I(x)})$ is smooth over~$S$ of
codimension
$\card I(x)$
in~$X$.

The divisor $D$ is said to have \emph{normal crossings relative
to} $S$ if, locally on~$X$ for the étale topology, it has
strictly normal crossings.

Let $D$ be a divisor with normal crossings relative
to~$S$. We set $Y=\Supp D$, $U=X-Y$, and we denote by $i\colon U\to X$
the canonical immersion. For every geometric point
$\overline{s}$ of~$S$ and for every maximal point $y$ of the geometric
fiber $Y_{\overline{s}}$, the ring
$R=\cal{O}_{X_{\overline{s}},y}$ is a discrete valuation
ring.

In the rest of this no., we shall use the following
technical definition:

\begin{subdefinition}
\label{XIII.2.1.1}
Let $F$ be a sheaf of sets on~$U$. One says that $F$ is
\emph{tamely ramified over}~$X$
(\emph{along}~$D$) \emph{relative to}~$S$ if, for every
geometric point $\overline{s}$ of~$S$, the following condition
is satisfied:

For every maximal point $y$ of~$Y_{\overline{s}}$, the restriction of
$F$ to the field of fractions $K$ of~$\cal{O}_{X_{\overline{s}}}$ is
representable by the spectrum of an étale $K$-algebra $L$,
tamely ramified over~$\cal{O}_{X_{\overline{s},y}}$.
\end{subdefinition}

Most often, when no confusion can result from it, we
shall omit the mention of~$D$ in the terminology.

\begin{subdefinition}
\label{XIII.2.1.2}
If
% original p. 373
$F$ is a sheaf of groups on~$U$, tamely
ramified over~$X$ relative to~$S$, one denotes by
$$
\H^1_{\tame}(U,F)
$$
the subset of~$\H^1(U,F)$ formed by the classes of torsors under
$F$ which are tamely ramified over~$X$ relative
to~$S$.
\end{subdefinition}

Let
$$
\xymatrix{ U \ar[r]^i \ar[d]_g & X \ar[dl]^f \\ T}
$$
be a commutative diagram of~$S$-schemes, with $i$ as
in~\Ref{XIII.2.1}; one denotes by
$$
\R^1_{\tame} g_* F
$$
the sheaf on~$T$ associated with the presheaf
$T'\mto\H^1_{\tame}(U',F)$, where $T'$ runs through the schemes
étale over~$T$ and where $U'=U\times_T T'$; $\R^1_{\tame} g_*F$
is a subsheaf of~$\R^1g_*F$.

Note that, if $g$ is coherent, if $\overline{t}$ is a geometric
point of~$T$, $\overline{T}$ the strict localization of
$T$ at $\overline{t}$, $\overline{U}=U\times_T\overline{T}$, one has an
isomorphism
\begin{equation*}
\label{eq:XIII.2.1.2.1}
\tag{\thesubsubsection.1} (\R^1_{\tame} g_*F)_{\overline{t}}\simeq
\H^1_{\tame}(\overline{U},\overline{F})
\end{equation*}

\subsubsection{}
\label{XIII.2.1.3}
Let $C_t((U,X)/S)$, or simply $C_t$,
be the category of étale coverings of~$U$ which are
tamely ramified over~$X$ relative
to~$S$. Suppose $U$ connected and let $a$ be a
geometric point of~$U$. Let $\Gamma_t$ be the functor which, to an
étale covering~$U'$ of~$U$ tamely ramified
over~$X$ relative to~$S$, associates the set of the geometric
points of~$U'$ over~$a$. From~\Ref{XIII.2.0}
it follows that the pair $(C_t,\Gamma_t)$ satisfies axioms
($G_1$) to ($G_6$) of V~\Ref{V.4}. Consequently $\Gamma_t$ is
representable by a pro-object which is called the
\emph{tamely ramified universal covering
% original p. 374
of~$(U,X)$ relative to~$S$ pointed at~$a$}.
The group opposite to the group of $U$-automorphisms of the
tamely ramified universal covering is called
the \emph{tamely ramified fundamental group}
and denoted
$$
\pi_1^{\tame}((U,X)/S,a)\text{ or simply }\pi_1^{\tame}(U,a)\text{ or
even }\pi_1^{\tame}(U).
$$
It is obviously a quotient of the fundamental group $\pi_1(U,a)$
(V~\Ref{V.6.9}).

\subsubsection{}
\label{XIII.2.1.4}
Let $F$ be a sheaf of groups on~$U$, $P$ a right torsor
with group $F$, $Q$ a left torsor with group $F$, and suppose
$P$ and $Q$ tamely ramified over~$X$ relative to~$S$. Then so is $P\stackrel{F}{\wedge}Q$. Indeed one
reduces to showing that, if $R$ is a discrete
valuation ring with field of fractions $K$ and if~$F$ is a
finite étale group scheme over~$K$, and if $P$ and $Q$ are
two torsors under~$F$ tamely ramified over~$R$, then
so is $P\stackrel{F}{\wedge} Q$. Now
$T=P\stackrel{F}{\wedge} Q$ is a quotient of~$P\times_K Q$. If $L$,
$M$, $N$ denote the $K$-algebras representing
$T$, $P$, $Q$ respectively, then $L$ is a subalgebra of
$M\otimes_K N$, and it follows from~\Ref{XIII.2.0.3} that $L$ is
tamely ramified over~$R$.

We deduce from the preceding that, if $F$ is a sheaf of
groups on~$U$ and if there exists a torsor with group $F$
tamely ramified over~$X$ relative to~$S$, then
$F$ is tamely ramified over~$X$ relative to~$S$. Indeed the torsor $P^{\circ}$ opposite to~$P$ is
tamely ramified over~$X$ relative to~$S$,
since it is isomorphic to $P$ as a sheaf of sets. If
$\leftexp{P\mkern-4mu}{F}$ is the group obtained by twisting $F$ by $P$, one has an
isomorphism
$$
F\simeq P^{\circ}\stackrel{\leftexp{P\mkern-4mu}{F}}{\wedge}P
$$
and consequently $F$ is tamely ramified over~$X$
relative to~$S$.

One sees as before that, if $F\to F'$ is a morphism
of sheaves of groups on~$U$, tamely ramified over~$X$ relative to~$S$, and if $P$ is a torsor under $F$
tamely ramified over~$X$ relative
% original p. 375
to~$S$, then the torsor $P'$ deduced from~$P$ by the extension of the
structure group $F\to F'$ is tamely ramified over~$X$
relative to~$S$.

In particular the canonical morphism
$$
\H^1(U,F)\to \H^1(U,F')
$$
gives by restriction to $\H^1_{\tame}(U,F)$ a canonical morphism
$$
\H^1_{\tame}(U,F)\to\H^1_{\tame}(U,F')
$$

\subsubsection{}
\label{XIII.2.1.5}
Let $S'\to S$ be a morphism and let us denote by $U'$ (\resp $X'$, etc.)
the inverse image of~$U$ (\resp $X$, etc.) over~$S'$. If $F$ is a
sheaf of sets on~$U$ tamely ramified over~$X$
relative to~$S$, it follows from
definition~\Ref{XIII.2.1.1} and from~\Ref{XIII.2.0.3} that $F'$ is
tamely ramified over~$X'$ relative to~$S'$.

If now $F$ is a sheaf of groups on~$U$, the inverse image
over~$S'$ of a torsor under $F$ tamely ramified over~$X$
relative to~$S$ is a torsor under $F'$ tamely
ramified over~$X'$ relative to~$S'$. In particular one has a
canonical functor
\begin{equation*}
\label{eq:XIII.2.1.5.1}
\tag{\thesubsubsection.1} C_t((U,X)/S)\to C_t((U',X')/S').
\end{equation*}
Suppose $U$ and $U'$ connected and let $a$ be a
geometric point of~$U$, $a'$ a geometric point of~$U'$
over~$a$; one deduces from the preceding a canonical
morphism
\begin{equation*}
\label{eq:XIII.2.1.5.2}
\tag{\thesubsubsection.2} \pi_1^{\tame}(U',a')\to\pi_1^{\tame}(U,a).
\end{equation*}

If $S'\to S$ is a morphism and $h\colon T'\to T$ the canonical
projection, the morphism
$$
h^*(\R^1g_*F)\to\R^1g_*'F'
$$
gives by restriction a canonical morphism
\begin{equation*}
\label{eq:XIII.2.1.5.3}
\tag{\thesubsubsection.3} h^*(\R^1_{\tame}g_*F)\to\R^1_{\tame}g_*'F'
\end{equation*}

\subsubsection{}
\label{XIII.2.1.6}
Let
% original p. 376
$F$ be a sheaf of groups on~$U$, tamely
ramified over~$X$ relative to~$S$. The notation being
that of~\Ref{XIII.2.1.2}, one has canonical exact sequences:
\begin{equation*}
\label{eq:XIII.2.1.6.1}
\tag{\thesubsubsection.1}
\begin{array}{ccccccc}
1 & \longrightarrow & \H^1(X,i_*F)& \longrightarrow & \H^1_{\tame}(U,F) &
\longrightarrow & \H^0(X,\R^1_{\tame} i_*F)\\[10pt] 1 &
\longrightarrow & \R^1f_*(i_*F) & \longrightarrow & \R^1_{\tame}g_*F &
\longrightarrow & f_*(\R^1_{\tame}i_*F).\end{array}
\end{equation*}

The first is obtained from the exact sequence
(SGA~4 III~3.2):
$$
1\to \H^1(X,i_*F)\to \H^1(U,F)\to
\H^0(X,\R^1i_*F).
$$
Indeed it suffices to show that the image of~$\H^1(X,i_*F)$ in
$\H^1(U,F)$ is in fact contained in $\H^1_{\tame}(U,F)$ and that the image of
$\H^1_{\tame}(U,F)$ in $\H^0(X,\R^1i_*F)$ is in fact contained in
$\H^0(X,\R^1_{\tame}i_*F)$. Now the inverse image over~$U$ of a torsor under
$i_*F$ is a torsor under $i^*i_*F$ which is obviously
tamely ramified over~$X$ relative to~$S$, hence the same
holds after the extension of the structure group
$i^*i_*F\to F$, which proves the existence of the arrow
$\H^1(X,i_*F)\to\H^1_{\tame}(U,F)$. The fact that the image of~$\H^1_{\tame}(U,F)$
in $\H^0(X,\R^1i_*F)$ is contained in $\H^0(X,\R^1_{\tame}i_*F)$
follows immediately from the definition of~$\R^1_{\tame}i_*F$. This
proves the existence of the first exact sequence, and the second
is deduced from it by localization.

\subsection{}
\label{XIII.2.2}
We keep the notation of~\Ref{XIII.2.1}. We are going to define
a notion of tamely ramified object of a stack $\Phi$
on~$U$ when the latter is given, locally \emph{on} $X$ and
on~$S$ for the étale topology, as \emph{the inverse image of a
stack} $\Psi$ \emph{on}~$S$.

First let $G$ be a gerbe on~$U$ and suppose given an
étale surjective morphism $S_1\to S$, an étale surjective morphism
$X_2\to X\times_S S_1$, a trivial gerbe $H$ on~$S_1$ and an
isomorphism
$$
G|{U_2}\to H|{U_2},
$$
where
% original p. 377
$U_1=U\times_X X_1$, $U_2=U\times_X X_2$. When one chooses a
trivialization of~$H|{X_2}$, the isomorphism above identifies
$G|{U_2}$ with the stack of torsors under a sheaf of groups $F$. One
says that an element $x$ of~$G_U$ is tamely
ramified over~$X$ relative
to~$S$
if the restriction of~$x$
to $U_2$ is a torsor tamely ramified over~$X$
relative to~$S$. By~\Ref{XIII.2.1.4} this notion does not
depend on the way in which one has trivialized~$H|X_2$.

Now let $\Phi$ be a stack on~$U$ and suppose given an
étale surjective morphism $S_1\to S$, an étale surjective
morphism $X_2\to X\times_S S_1$, a stack $\Psi$ on~$S_1$ and an
isomorphism
$$
i\colon \Phi|U_2\to \Psi|U_2.
$$
Let $x$ be an element of~$\Phi_U$, $G_x$ the maximal subgerbe
of~$\Phi$ generated by $x$
\cite[III~2.1.7]{XIII.1}, $S\Phi$ the sheaf of
maximal subgerbes of~$\Phi$. The isomorphism $i$ induces an
isomorphism
$$
S\Phi|U_2\to S\Psi|U_2.
$$
It follows from~\Ref{XIII.5.7} that, after replacing $S_1$ if necessary by
an étale surjective extension, one has a unique maximal
subgerbe $H$ of~$\Psi$, which one may assume trivial, such that $i$
defines an isomorphism
$$
G_x|U_2\to H|U_2.
$$
One says that the element $x$ is tamely ramified
over~$X$ relative to~$S$ if it is so as an element
of~$G_x$ endowed with the isomorphism above.

\subsubsection{}
\label{XIII.2.2.1}
Let $\Phi$ be a stack on~$U$ given, locally on~$X$ and $S$,
as the inverse image of a stack on~$S$, and let
$$
\xymatrix{ U \ar[r]^i
\ar[d]_g & X \ar[dl]^f \\ T }
$$
be a diagram as in~\Ref{XIII.2.1.2}. For every scheme $T'$
étale over~$T$, if $U'=U\times_T T'$, one considers the
subset $(g_*^{\tame}\Phi)_{T'}$
of~$(g_*\Phi)_{T'}=\Phi_{U'}$
formed
% original p. 378
by the elements of~$\Phi_{U'}$ which are tamely
ramified
over~$X$ relative to~$S$. One calls tamely ramified
direct image
of~$\Phi$ by~$g$, and one denotes by
$$
g_*^{\tame}\Phi
$$
the full subcategory of~$g_*\Phi$ whose objects over
a scheme $T'$ étale over~$T$ are the elements of
$(g_*^{\tame}\Phi)_{T'}$. It is clear that $g_*^{\tame}\Phi$ is a substack of
$g_*\Phi$.

\subsubsection{}
\label{XIII.2.2.2}
By reduction to the case of a stack of torsors, one sees that, if
$\Phi$ is a stack on~$U$ which is, locally for the
étale topology of~$S$ and $X$, the inverse image of a stack on~$S$, the canonical
morphism
$$
h^*(g_*\Phi)\to g_*'\Phi'
$$
gives by restriction a canonical morphism
$$
h^*(g_*^{\tame}\Phi)\to g_*^{\prime\tame}\Phi'
$$

\begin{remarks}
\label{XIII.2.3}

a) If $F$ is a locally constant constructible sheaf of sets
on~$U$, for $F$ to be tamely
ramified over~$X$ relative to~$S$, it suffices that the condition
of~\Ref{XIII.2.1.1} be satisfied for the geometric
points of~$S$ over the maximal points of~$S$. To
see this we may assume that the divisor $D$ has strictly normal
crossings. The sheaf $F$ is representable by an étale
covering $V$ of~$U$. If $\overline{s}$ is a geometric
point of~$S$, $y$ a maximal point of
$Y_{\overline{s}}$, we denote by $\overline{S}$ the strict localization of
$S$ at $\overline{s}$, $\overline{X}$ the strict localization of~$X$ at
$\overline{y}$, $\overline{U}=U\times_X \overline{X}$,
$\overline{V}=V\times_X\overline{X}$. If the condition
of~\Ref{XIII.2.1.1} is satisfied at the geometric points
over the maximal points of~$S$, it follows from~\Ref{XIII.5.5}
below that $\overline{V}$ is a covering of~$\overline{U}$
tamely ramified over~$\overline{X}$ relative to
$\overline{S}$; consequently $V$ is an étale covering of~$U$
tamely ramified over~$X$ relative to~$S$.

b)
% original p. 379
Let $F$ be a sheaf of groups on~$U$ tamely
ramified over~$X$ relative to~$S$. If $\overline{s}$ is a
geometric point of~$S$, $y$ a maximal point of
$Y_{\overline{s}}$, we denote by $K$ the field of fractions of
$\cal{O}_{X_{\overline{s}},y}$. Suppose that, for every point
$\overline{s}$ and for every point $y$, the $K$-algebra $L$ whose
spectrum represents $F|K$ is of rank prime to the
residue characteristic $p$ of
$\cal{O}_{X_{\overline{s}}}$. We shall sometimes say, by abuse of language,
that $F$ is prime to the residue characteristics of~$S$.
When this is so, every torsor $P$ under $F$ is
tamely ramified over~$X$ relative to~$S$.

Indeed, let $\overline{R}$ be the strict localization of
$\cal{O}_{X_{\overline{s}},y}$ at $\overline{y}$, $\overline{K}$ its
field of fractions, $\overline{F}$ the inverse image of~$F$ on~$\overline{K}$. Let us show that we may assume $F$ constant. Since
$F$ is tamely ramified over~$X$ relative to~$S$,
$\overline{F}$ is representable by the spectrum of a
$\overline{K}$-algebra $L=\prod L_i$, where the $L_i$ are
extensions of~$\overline{K}$ of degree prime to $p$. One can
therefore find an extension $K'$ of~$\overline{K}$ of degree prime
to $p$ such that $\overline{F}|K'$ is a constant
sheaf. By~\Ref{XIII.2.0.3}, to prove that
$P|\overline{K}$ is tamely ramified over~$\overline{R}$, it suffices to see that $P|K'$ is tamely
ramified over the integral closure of~$\overline{R}$ in
$K'$, whence the reduction to the case where $\overline{F}$ is
constant. Suppose henceforth $\overline{F}$ constant. The
$\overline{K}$-algebra $H$ which represents $P|\overline{K}$ is
then a product of extensions $H_i$ of~$\overline{K}$ isomorphic to
one another. Since the rank of~$H$ is prime to $p$, so is
$[H_1:\overline{K}]$, which proves that $H$ is
tamely ramified over~$X$ relative to~$S$.

c) Let $X$ be a regular scheme, $D$ a divisor with normal
crossings of~$X$ (SGA~5 I~3.1.5), $U=X-\Supp D$, $F$ a
sheaf of sets on~$U$. If $y$ is a maximal point of~$\Supp
D$, we denote by $K$ the field of fractions of
$\cal{O}_{X,y}$. One says that $F$ is \emph{tamely
ramified relative to}~$D$
if, for every maximal point $y$ of~$\Supp D$, $F|K$ is
representable by a $K$-algebra tamely
ramified over~$\cal{O}_{X,y}$.
\end{remarks}

\begin{theorem}
\label{XIII.2.4}
Let
% original p. 380
$f\colon X\to S$ be an $S$-scheme, $D$ a divisor on~$X$
with normal crossings relative to~$S$ \eqref{XIII.2.1},
$Y=\Supp D$, $U=X-Y$, $i\colon U\to X$ the canonical immersion. Let $F$
be a sheaf of sets (\resp of groups) on~$U$, satisfying
one of the following conditions:
\begin{enumerate}
\item[a)] $F$ is, locally for the \'etale topology \emph{on}
$X$ and on~$S$, the inverse image of a sheaf of sets (\resp of a
constructible sheaf of groups) on~$S$.
\item[b)] $F$ is locally constant constructible on~$U$ and
tamely ramified over~$X$ relative to~$S$.
\end{enumerate}
Then we have the following conclusions:
\begin{enumerate}
\item[1)]$(F,i)$ is cohomologically proper relative to~$S$
in dimension $\leq 0$ (\resp for every morphism $h\colon S'\to S$, if
$i'\colon U'\to X'$ is the inverse image of~$i$ over~$S'$, if $F'=F|U'$
and if $k=h_{(X)}$, the canonical morphism
$$
\Psi\colon k^*(\R^1_{\tame}i_*F)\to\R^1_{\tame}i_*'F'
$$
is an isomorphism).

If $F$ is a sheaf of groups prime to the residue characteristics
of~$S$ (\Ref{XIII.2.3}.b)) (tamely
ramified over~$X$ relative to~$S$), then $(F,i)$ is
cohomologically proper relative to~$S$ in dimension $\leq 1$.
\item[2)]If $F$ is a constructible sheaf of sets (\resp of groups),
$i_*F$ (\resp $\R^1_{\tame}i_*F$) is constructible.
\end{enumerate}
\end{theorem}

\subsubsection*{Proof}
For every $S$-scheme $S'$, we consider the following diagram,
all of whose squares are cartesian:
$$
\xymatrix{ U \ar[dd]_-g \ar[dr]^-i & & & U'\ar[lll] \ar[dd]_{\hbox{\raise10mm\hbox{$g'$}}}
\ar[dr]^-{i'}& \\ & X \ar[dl]_{f\mkern-8mu} & & & X'\ar[lll]_{k\hspace*{1cm}} \ar[dl]_{f'\mkern-10mu} \\ S
& & & S'\ar[lll]_-h & & }
$$
Since
% original p. 381
the question is local on~$X$ for the \'etale topology, we may
assume that $D$ is a divisor with strictly normal crossings
relative to~$S$ \eqref{XIII.2.1}; moreover, after
restricting $X$ to a neighborhood of~$Y$, we may assume $X$ smooth
over~$S$.

\Subsubsection*{Proof of~\Ref{XIII.2.4}~\textup{1)}}
\subsubsection{}
\label{XIII.2.4.1}
\emph{Case of a sheaf of sets satisfying} a). We may
assume that we have $F=g^*G$, where $G$ is a sheaf on~$S$. It
then follows from (SGA~4 XVI~3.2) that the canonical morphism
\begin{equation*}
\label{eq:XIII.2.4.1.1}
\tag{\thesubsubsection.1} f^*G\to i_*F
\end{equation*}
is an isomorphism. For every $S$-scheme $h\colon S'\to S$, we have
likewise an isomorphism $f'^*G'\to i_*'F'$; consequently the canonical
morphism
$$
\varphi\colon k^*(i_*F)\to i_*'F'
$$
is identified with the natural isomorphism
$$
k^*f^*G\simeq f'^*G'
$$

\subsubsection{}
\label{XIII.2.4.2}
\emph{Case of a sheaf of sets satisfying} b). We must
show that $\varphi$ is an isomorphism, and for this it suffices to
see that this is so at each geometric point
$\overline{x}'$ of~$X'$. Let $\overline{S}$ (\resp $\overline{X}$,
\resp $\overline{S'}$, \resp $\overline{X'}$) be the strict localization
of~$S$ (\resp~$X$, \resp~$S'$, \resp $X'$) at $\overline{x}'$ and let us put
$\overline{U}=U_{(\overline{X})}$,
$\overline{U'}=U_{(\overline{X'})}$, etc. The morphism
$\varphi_{\overline{x}}$ is identified with the canonical morphism
$$
\overline{\varphi}\colon \H^0(\overline{U},\overline{F})\to
\H^0(\overline{U'}, \overline{F'})
$$
We can find a principal covering $V$ of~$\overline{U}$, of the
type appearing in~\Ref{XIII.5.4}, such that the inverse image of
$\overline{F}$ on~$V$ is a constant sheaf with value $C$. If
$\Pi$ is the Galois group of~$V$ over~$\overline{U}$, $\Pi$
operates on~$\overline{F}|V$, and we have
\begin{equation*}
\label{eq:XIII.2.4.2.1}
\tag{\thesubsubsection.1}
\H^0(\overline{U},\overline{F}))\simeq \H^0(V,C_V)^{\Pi}
\end{equation*}
where the right-hand side denotes the set of the
elements of~$\H^0(V,C_V)$ invariant
% original p. 382
under~$\Pi$. Since $V'=V\times_{\overline{U}}\overline{U'}$ is a
principal covering of~$\overline{U'}$ with Galois group
$\Pi'\simeq \Pi$, we see that the morphism $\overline{\varphi}$
is obtained, by taking the invariants under $\Pi$, from the
canonical morphism
$$
\H^0(V,C_V)\to \H^0(V',C_{V'}).
$$
Since $V$ and $V'$ are connected \eqref{XIII.5.4}, this morphism, hence
also $\overline{\varphi}$, is an isomorphism.

Let us note that if moreover $F$ is a sheaf of groups and if $P$ is a
torsor on~$\overline{U}$ with group $\overline{F}$,
tamely ramified relative to $D$, it follows
from the preceding proof and from~\Ref{XIII.2.2} that
$(\leftexp{P}{\overline{F}},\overline{i})$ is cohomologically
proper relative to~$S$ in dimension $\leq 0$.

\subsubsection{}
\label{XIII.2.4.3}
\emph{Case of a sheaf of groups.} To show that $\Psi$ is an
isomorphism, it suffices to prove that, for every geometric
point $\overline{y}'$ of~$Y'$, the morphism
$$
\Psi_{\overline{y}'}\colon
(k^*(\R^1_{\tame}i_*F))_{\overline{y}'}\to (\R^1_{\tame}
i_*'F')_{\overline{y}'}
$$
is an isomorphism. Now, by~\Ref{XIII.2.1.2},
$\Psi_{\overline{y}'}$ is identified with the canonical morphism
$$
\overline{\Psi}\colon
\H^1_{\tame}(\overline{U},\overline{F})\to
\H^1_{\tame}(\overline{U'},\overline{F'}).
$$
Let $\widetilde{U}$ be the universal tamely
ramified covering of~$\overline{U}$ \eqref{XIII.2.1.3} and $\widetilde{F}$
the inverse image of~$\overline{F}$ on~$\widetilde{U}$. It follows
from~\Ref{XIII.5.7} in case a) and from~\Ref{XIII.5.5} in case b)
that $\H^1_{\tame}(\overline{U},\overline{F})$ is identified with the
subset $\H^1(\overline{U},\overline{F})$ formed by the classes
of~$\overline{F}$-torsors whose inverse image on~$\widetilde{U}$
is trivial. On the other hand a
classical argument (\cf IX~\Ref{IX.5}, p\ptbl\pageref{page-300})
shows that
the set of the elements of~$\H^1(\overline{U},\overline{F})$
whose inverse image on~$\widetilde{U}$ is trivial is identified
with
$\H^1(\pi_1^{\tame}(\overline{U}),\H^0(\widetilde{U},\widetilde{F})).$
We thus obtain a canonical isomorphism
\begin{equation*}
\label{eq:XIII.2.4.3.1} \tag{\thesubsubsection.1} {}
\H^1_{\tame}(\overline{U},\overline{F})\isomto
\H^1(\pi_1^{\tame}(\overline{U}),\H^0(\widetilde{U},\widetilde{F})).
\end{equation*}

Consequently the morphism $\overline{\Psi}$ is identified with the
canonical morphism
$$
\H^1(\pi_1^{\tame}(\overline{U}),\H^0(\widetilde{U},\widetilde{F}))
\to
\H^1(\pi_1^{\tame}(\overline{U'}),\H^0(\widetilde{U}',\widetilde{F}')).
$$
Let us show
% original p. 383
that this morphism is an isomorphism. The morphism
$\pi_1^{\tame}(\overline{U'})\to\pi_1^{\tame}(\overline{U})$
is an isomorphism by~\Ref{XIII.5.6}, and the same
holds for the morphism $\H^0(\widetilde{U},\widetilde{F})\to\allowbreak
\H^0(\widetilde{U}',\widetilde{F}')$. Indeed, this is obvious
in case b) since $\widetilde{F}$ is constant and $\widetilde{U}$ and
$\widetilde{U}'$ are connected. In case a), let $\overline{G}$ be a
constructible sheaf of groups on~$\overline{S}$ such that we have
$\overline{F}=\overline{g}^*\overline{G};$ since the morphisms
$\widetilde{U}\to\overline{S}$ and
$\widetilde{U}'\to\overline{S'}$ are $0$-acyclic
\eqref{XIII.5.7}, we have
$$
\H^0(\widetilde{U},\widetilde{F})\simeq
\H^0(\overline{S},\overline{G}) \simeq
\H^0(\overline{S'},\overline{G'})\simeq
\H^0(\widetilde{U}',\widetilde{F}'),
$$
which implies that $\Psi$ is an isomorphism. The last
assertion of \Ref{XIII.2.4}~1) follows from what precedes,
taking into account \Ref{XIII.2.3}~b).

\subsubsection*{Proof of~\Ref{XIII.2.4}~\textup{2)}}

The case of a constructible sheaf of sets satisfying a)
follows at once from \eqref{eq:XIII.2.4.1.1}. Let $F=g^*G$ be a
sheaf of groups satisfying a), where $G$ is a constructible
sheaf; we may assume $S$ affine; let $(S_j)_{j\in J}$ be
a finite family of reduced closed subschemes of~$S$
whose union covers $S$, such that the inverse image of~$G$ on~$S_j$ is a locally constant sheaf. Taking into account
\Ref{XIII.2.4}~1), it suffices, in order to establish that $\R^1_{\tame}i_*F$
is constructible, to see that this is so after the change
of base $S_j\to S$, for each $j\in J$. We are thus
reduced to case~b) where $F$ is locally constant.

We assume from now on that $F$ is a sheaf of sets or of
groups satisfying b). Since the question is local for the
\'etale topology on~$X$, we may assume $X$ of finite
presentation over~$S$, and, by passage to the limit, we may assume $X$
and $S$ noetherian.

Let $D=\sum_{1\leq i\leq r}\divisor f_i$, where, for each point
$x$ of~$\Supp D$, if $I(x)$ is the set of the $i$ such that
$f_i(x)=0$, the subscheme $V((f_i)_{i\in I(x)})$ is smooth over~$S$ of codimension $\card I(x)$ in~$X$. Let $\mathcal{P}$ be
the set of subsets of~$[ 1,r ]$ and, for each $I\in\mathcal{P}$,
let us put
$$
X_I=\textstyle\Big(\bigcap_{i\in I}V(f_i)\Big)\cap\Big(\bigcap_{i\not\in I}X_{f_i}\Big).
$$
Let
% original p. 384
$z$ be a point of~$X_I$. After first restricting to an
\'etale neighborhood of~$z$, we can find an open subset $W$ of~$X$
containing $z$ and a principal covering $V$ of~$U\cap W$,
tamely ramified over~$W$ relative to~$S$, of the type
considered in~\Ref{XIII.5.6.1}, such that the inverse
image of~$F$ on~$V$ is a constant sheaf with
value~$C$. Let $\pi$ be the Galois group of the covering~$V$. For
every geometric point $\overline{x}$ of~$X_I$, we then have,
by \eqref{eq:XIII.2.4.2.1},
$$
\H^0(\overline{U},\overline{F})\simeq \H^0(V,C_V)^{\pi}.
$$
It follows that $i_*F|X_I\cap W$ is locally constant (SGA~4
IX~2.13) and consequently $i_*F$ is constructible.

Let us finally show that if $F$ is a locally constant sheaf of groups,
$\R^1_{\tame}i_*F|_{X_I}$ is constructible. If
$\overline{x}$ is a geometric point of~$X$, we have obtained
in \eqref{eq:XIII.2.4.3.1} the expression
$$
\H^1_{\tame}(\overline{U},\overline{F})\isomto
\H^1(\pi^{\tame}_1(\overline{U}),\H^0(\widetilde{U},\widetilde{F})).
$$
If $p$ is the residue characteristic of~$\overline{X}$, we
have, by~\Ref{XIII.5.6},
$\pi_1^{\tame}(\overline{U})=\prod_{l\neq p }\ZZ_l(1)^{\card I}$. Let us denote by $\LL$ the set of prime numbers which
divide the order of the finite group $\H^0(\widetilde{U},\widetilde{F})$ and
let $K=\prod_{l\in\LL -\{p\}\cap\LL }\ZZ_l(1)^{\card I}$. It
follows from \cite[I \S5 ex.2]{XIII.4} that we have
$$
\H^1_{\tame}(\overline{U},\overline{F})\simeq
\H^1(K,\H^0(\widetilde{U},\widetilde{F})).
$$
Since $K$ is topologically of finite type and
$\H^0(\widetilde{U},\widetilde{F})$ finite, we first of all
deduce from this that the fibers of the sheaf $\R^1_{\tame}i_*F|_{X_I}$ are
finite. On the other hand, the set $\LL$ does not depend on the point
$\overline{x}$. For every $q\in\LL$, let $X_{I,q}$ be the closed subset of
$X_I$ with equation $q=0$ and let $X_{I'}$ be the open subset of~$X_I$
complementary to the union of the $X_{I,q}$. Then
$\R^1_{\tame}i_*F|_{X_{I,q}}$ and $\R^1_{\tame}i_*F|_{X_I'}$ are
locally constant; indeed a specialization arrow of
geometric points of~$X_{I,q}$ (\resp of~$X_{I'}$) induces an
isomorphism on the groups $K$ \eqref{XIII.5.6.1} hence also on the
sets $\H^1_{\tame}(\overline{U},\overline{F})$, and one can apply
SGA~4 IX~2.13.

\begin{corollary}
\label{XIII.2.5}
Let
% original p. 385
$f\colon X\to S$ be a morphism, $D$ a divisor on~$X$
with normal crossings relative to~$S$ \eqref{XIII.2.1},
$Y=\Supp D$, $U=X- Y$, $i\colon U\to X$ the canonical
immersion. Let $\Phi$ be a stack on~$U$ and suppose given
surjective \'etale morphisms $S_1\to S$ and $X_2\to
X\times_{S}S_1$, a stack $\Psi$ on~$S_1$ and an isomorphism
$\Phi|U_2\simeq\Psi|U_2$ (\cf \Ref{XIII.2.2}).

Then, for every morphism $h\colon S'\to S$, if $k=h_{(X)}$,
the canonical functor
$$
\varphi\colon k^*i_*\Phi\to i_*'\Phi'
$$
is fully faithful. If $\Psi$ is constructible, the canonical
functor
$$
\psi\colon k^*i_*^{\tame}\Phi\to i_*^{\prime\tame}\Phi'
$$
is an equivalence of categories.

Moreover, if the stack $\Psi$ is constructible (\resp if $\Psi$ is
$1$-constructible \eqref{XIII.0} and $S$ locally noetherian),
$i_*\Phi$ is constructible (\resp $i_*^t\Phi$ is
$1$-constructible).
\end{corollary}

Let us show that $\varphi$ is fully faithful. It suffices to see
that, for every geometric point $\overline{x}'$ of~$X'$, this
is so for the functor
$$
\overline{\varphi}\colon
\Phi(\overline{U})\to\Phi'(\overline{U'})
$$
(we have taken up again the notation of~\Ref{XIII.2.4.2}). Let $a$, $b$ be two
elements of~$\Phi(\overline{U})$, $a',\ b'$ their images under
$\overline{\varphi}$. Since the morphism
$\overline{U}\to\overline{S}$ is locally $0$-acyclic
\eqref{XIII.5.7}, we have an isomorphism
$$
\H^0(\overline{U},\overline{S\Phi})
\overset{\sim}{\to}\H^0(S,\overline{S\Psi}).
$$
Consequently $a$ and $b$ come by inverse image from elements
of~$\Psi$, and hence so does
$F=\SheafHom_{\overline{U}}(a,b)$. Since the sheaf
$F'=\SheafHom_{\overline{U'}}(a',b')$ is the inverse image on~$\overline{U'}$ of~$F$, it follows from~\Ref{XIII.2.4.1} that the
canonical morphism
$$
\H^0(\overline{U},F)\to \H^0(\overline{U'},F')
$$
is an isomorphism, which proves that $\overline{\varphi}$ is
fully faithful.

Let us show
% original p. 386
that, for every geometric point $\overline{x}'$ of~$X'$, the
functor
$$
\overline{\psi}\colon i^{\tame}_*\Phi(\overline{X'})\to
i_*^{\prime\tame}\Phi'(\overline{X'})
$$
is an equivalence. By what precedes,
$\overline{\psi}$ is fully faithful. Let us show that
$\overline{\psi}$ is essentially surjective. Let $a'$ be an
element of~$\Phi'(\overline{U'})$ tamely
ramified over~$X'$ relative to~$S'$ \eqref{XIII.2.2} and
let us show that it is the image of a tamely
ramified element of~$\Phi(\overline{U})$. It follows from~\Ref{XIII.2.4}
1) that the canonical morphism
$$
\H^0(\overline{U},\overline{S}\overline{\Phi})\to
\H^0(\overline{U'},\overline{S}\overline{\Phi}')
$$
is an isomorphism. Let $G'$ be the maximal subgerbe of~$\Phi'$
generated by $a'$; there then exists a maximal subgerbe $G$ of
$\overline{\Phi}$, inverse image of a gerbe on~$\overline{S}$, such
that we have
$$
\overline{m}^*G\simeq G',
$$
where $m$ is the morphism $\overline{U'}\to
\overline{U}$. The canonical functor
\begin{equation*}
\label{eq:XIII.2.5.*}
\tag{$*$} {\overline{k}^*\overline{i}_*^{\tame}G\to
\overline{i}_*^{\prime\tame}G'}
\end{equation*}
is an equivalence, since $G$ is identified with a gerbe of
torsors under a constructible sheaf of groups coming from
$\overline{S}$, and one can apply \Ref{XIII.2.4}~1). It
then follows from \eqref{eq:XIII.2.5.*} that there exists an
element $a$ of~$G(\overline{U})$, tamely
ramified over~$X$ relative to~$S$, whose inverse image on~$\overline{U'}$ is $a'$, which proves that $\psi$ is an
equivalence.

If $\Psi$ is constructible, so is~$i_*\Phi$; indeed
an object $x$ of~$i_*\Phi$ is, locally for the
\'etale topology of~$S$, the inverse image of an object $y$ of~$\Psi$; it
therefore follows from~\Ref{eq:XIII.2.4.1.1} that $\SheafAut (x)$ is
the inverse image of~$\SheafAut (y)$, hence is constructible. Finally, if
$\Psi$ is $1$-constructible, so is
$i^{\tame}_*\Phi$ by~\Ref{XIII.6.3} below.

\begin{corollary}
\label{XIII.2.6}
The
% original p. 387
notation is that of~\Ref{XIII.2.4}. Assume that $S$ is of
characteristic zero at every point $s$ such that we have $Y_s\neq
\emptyset.$ Then, if $\cal{ F}$ is a locally
constant constructible sheaf of groups on~$U$ (\resp a constructible stack on~$U$ which
is, locally on~$X$ and $S$, the inverse image of a constructible stack
on~$S$), $(\cal{ F},i)$ is cohomologically proper relative
to~$S$ in dimension $\leq 1.$
\end{corollary}

Since every constructible sheaf of sets on~$U$ is
tamely ramified over~$X$ relative to~$S,$ the
corollary follows from~\Ref{XIII.2.4} (\resp \Ref{XIII.2.5}).

\begin{corollary}
\label{XIII.2.7}
The notation is that of~\Ref{XIII.2.4}, but we are given in addition
an $S$-scheme $T$, a proper morphism $p\colon X\to T,$
and we assume $X$ and $T$ of finite presentation over~$S$; let
$q=pi$. Let $\cal{ F}$ be a constructible sheaf of sets on~$U$
satisfying one of the conditions \textup{a)} or \textup{b)} of~\Ref{XIII.2.4}
(\resp a sheaf of groups satisfying one of the conditions
\textup{a)}, \textup{b)} of~\Ref{XIII.2.4}, \resp a stack on~$U$ which is, locally on~$X$ and $S$, the inverse image of a constructible stack $G$ on~$S$). Then
we have the following conclusions:
\begin{enumerate}
\item[1)] $(\cal{ F},q)$ is cohomologically proper relative to~$S$ in dimension $\leq 0$ (\resp for every morphism $h\colon
S^{'}\to S$, if $m=h_{(T)}$, the canonical morphism:
$$
\Theta\colon m^*(\R^1_{\tame}q_*\cal{ F})\to
\R^1_{\tame}q_*^{'}\cal{ F}^{'}
$$
is an isomorphism, \resp for every morphism $S^{'}\to S$,
the canonical morphism
$$
\xi \colon m^*(q_*^{\tame}\cal{ F})\to
{q'}_*^{{\tame}}\cal{F}^{'}
$$
is an equivalence).
\item[2)] the sheaf $q_*\cal{ F}$ (\resp the sheaf
$\R_{\tame}^1q_*\cal{ F}$, \resp the stack $q_*^{\tame}\cal{ F}$)
is constructible. In the last case, if we assume $S$ locally
noetherian and $G$ $1$-constructible, the same is true
of $q_*^{\tame}\cal{ F}.$
\end{enumerate}
\end{corollary}

The
% original p. 388
first part follows at once from~\Ref{XIII.2.4},
\Ref{XIII.2.5} and from the proof
of~\Ref{XIII.1.8}. Let us prove~2). If $\cal{ F}$ is a constructible sheaf
of sets on~$U$ satisfying \Ref{XIII.2.4}~a)
or \Ref{XIII.2.4}~b), it follows from \Ref{XIII.2.4}~2) that
$i_*\cal{ F}$ is constructible; hence so is
$q_*\cal{ F}=p_*(i_*\cal{ F})$ (SGA~4 XIV~1.1).

Let $\cal{F}$ be a constructible sheaf of groups on~$U$
satisfying \Ref{XIII.2.4}~a) or \Ref{XIII.2.4}~b) and let us prove
that $\R_{\tame}^1q_*\cal{ F}$ is constructible. By passage to
the limit (EGA~IV 8.10.5 and 17.7.8) and using 1), we may
assume $S$ noetherian. Let then $\Phi$ be the stack on~$X$ whose
fiber at every scheme $X^{'}$ \'etale over~$X$ consists
of the torsors on~$U^{'}=U\times_{X}X^{'}$, with group $\cal{F}|U$, which
are tamely ramified over~$X$ relative to~$S$; we
therefore have
$$
S(i_*^{\tame}\Phi)\simeq \R_{\tame}^1i_*\cal{ F}
$$
and this sheaf is constructible by \Ref{XIII.2.4}~2).
It therefore follows from~\Ref{XIII.6.3} below that $S(p_*\Phi)$
is constructible, \ie that $\R^1_{\tame}q_*\cal{ F}$ is
constructible.

Finally, if $\cal{ F}$ is a stack on~$U$ which is, locally on~$X$ and
$S$, the inverse image of a constructible stack on~$S,$
$i_*^{\tame}\cal{ F}$ is constructible, and hence so
is~$q_*^{\tame}\cal{ F}=p_*(i_*^{\tame}\cal{ F})$. If
moreover $S$ is locally noetherian and $SG$ constructible,
$S(i_*^{\tame}\cal{ F})$ is constructible by 6.3; hence
so is~$S(q_*i_*^{\tame}\cal{ F})$, \ie
$S(q_*^{\tame}\cal{ F})$~\Ref{XIII.6.2}.

\begin{corollary}
\label{XIII.2.8}
Let
$$
\xymatrix{ U \ar[r]^{i} \ar[d]_{g} & X \ar[ld]^{f} \\ S }
$$
be a commutative diagram of schemes, in which $U$ is the open set
complementary
% original p. 389
in $X$ to a divisor with normal crossings relative to~$S$, $f$ a \emph {proper} morphism of finite presentation. Let
$\LL$ be a set of prime numbers. We assume~$g$ locally
$0$-acyclic (\resp locally $1$-aspherical for $\LL$). Then, if
$\cal{ F}$ is a sheaf of sets on~$U$ (\resp a sheaf of
$\LL$-groups) on~$U,$ locally constant constructible,
tamely ramified over~$X$ relative to~$S$,
$f_*\cal{ F}$ (\resp $\R^1_{\tame}f_*\cal{ F}$) is locally
constant constructible and $(\cal{ F},f)$ is cohomologically proper
relative to~$S$ in dimension $\leq 0$ (\resp the formation of
$\R^1_{\tame}f_*\cal{ F}$ commutes with every change of base
$S^{'}\to S$). In the case without ``resp.'', if $\cal{ F}$ is a
sheaf of groups, for every specialization $\overline{s}_1
\to \overline{s}_2$ of geometric points of~$S$, the
specialization morphism
$$
(\R^1_{\tame}f_*\cal{ F})_{\overline{s}_2}\to
(\R^1_{\tame}f_*\cal{ F})_{\overline{s}_1}
$$
is injective.
\end{corollary}
The corollary follows at once from~\Ref{XIII.2.6} and
from~\Ref{XIII.1.14} (\resp from the analogue of~\Ref{XIII.1.14} for
$\R^1_{\tame}f_*\cal{ F}$, which is proved as in \loccit).

\begin{corollary} \label{XIII.2.9}
Let
$$
\xymatrix{ U \ar[r]^{i} \ar[d]_{g} & X \ar[ld]^{f} \\ S }
$$
be a commutative diagram of schemes, in which $U$ is the open set
complementary in~$X$ to a divisor with normal crossings
relative to~$S$, $f$~a proper smooth morphism of
finite presentation. Let~$\LL $ be the set of prime numbers
distinct from the residue characteristics of~$S$. Let~$\cal{F}$ be a locally constant constructible sheaf of~$\LL $-groups
on~$U$, tamely ramified over~$X$ relative
to~$S$. Then $\R^1 f_*\cal{ F}$ is locally constant
constructible and $(\cal{ F},f)$ is cohomologically proper
relative
% original p. 390
to~$S$ in dimension $\leq 1$.
\end{corollary}
The corollary follows from~\Ref{XIII.2.8} and from the fact that we have
$\R_{\tame}^1f_*\cal{ F}=\R^1f_*\cal{ F}$ (\Ref{XIII.2.3}~b)).

\subsection{}
\label{XIII.2.10}
If $U$ is a connected scheme, $a$ a geometric point
of~$U$, $\LL$~a set of prime numbers, we denote by
\begin{equation*}
\label{eq:XIII.2.10.0}
\tag{2.10.0} \pi_1^{\LL}(U,a)
\end{equation*}
the projective limit of the finite quotients of~$\pi_1(U,a)$ whose
orders have all their prime factors in~$\LL$.

We shall define specialization morphisms for the
fundamental group, generalizing X.\Ref{X.2}.

Let $g\colon U\to S$ be a coherent morphism with geometrically
connected fibers (\resp a morphism of the form $g=
fi$, where $f\colon X \to S$ is a proper morphism of
finite presentation and where $i\colon U\to X$ is an
open immersion such that $U$ is the complement in $X$
of a divisor with normal crossings relative to~$S$
(\cf \Ref{XIII.2.8})). Let ${\LL}$ be a set of prime numbers and
suppose, in the non-resp.\ case, that, for every finite constant
${\LL}$-group $C$, $(C_U,g)$ is cohomologically proper
relative to~$S$ in dimension $\leq 1$. Let $\overline{s}_1
\to \overline{s}_2$ be a specialization morphism of
geometric points of~$S$, $\overline{S}$ the strict
localization of~$S$ at $\overline{s}_2$,
$\overline{U}=U\times_S\overline{S}$. We have a commutative diagram
$$
\xymatrix{ U_{\bar{s}_1} \ar[r]^{h_1} \ar[d] & \bar{U}
\ar[d]^{\bar{g}} & U_{\bar{s}_1} \ar[l]_{h_2} \ar[d] \\ \bar{s}_1
\ar[r] & \bar{S} & \,\bar{s}_2\,. \ar[l] }
$$
If $a_1$ is a geometric point of~$U_{\bar{s}_1},\ a_2$ a
geometric point of~$U_{\bar{s}_2}$, the morphisms $h_1$ and
$h_2$ define canonical morphisms
% original p. 391
\begin{align*}
\pi_1\colon \pi_1^{\LL}(U_{\bar{s}_1},a_1)&\to
\pi_1^{\LL}(\bar{U},a_1) & \pi_2\colon
\pi_1^{\LL}(U_{\bar{s}_2},a_2)&\to \pi_1^{\LL}(\bar{U},a_2)\\
(\text{\resp~} \pi_1\colon \pi_1^\tame (U_{\bar{s}_1},a_1)&\to
\pi_1^\tame (\bar{U},a_1) & \pi_2\colon
\pi_1^\tame(U_{\bar{s}_2},a_2)&\to \pi_1^\tame(\bar{U},a_2))
\end{align*}
(V~\Ref{V.7} and~\Ref{eq:XIII.2.1.5.2}). The hypotheses of
cohomological properness (\resp \Ref{XIII.2.8}) prove that $\pi_2$
is an isomorphism. If one chooses a class of paths from~$a_1$
to $a_2$, one obtains an isomorphism
$$
\pi_{12}\colon \pi_1^{\LL}(\bar{U},a_1)\isomto
\pi_1^{\LL}(\bar{U},a_2) \quad (\text{\resp~} \pi_{12}\colon
\pi_1^\tame (\bar{U},a_1)\isomto \pi_1^\tame(\bar{U},a_2))
$$
whence a morphism $\pi=\pi_2^{-1}\pi_{12}\pi_1$
$$
\pi\colon \pi_1^{\LL}(U_{\bar{s}_1},a_1)\to
\pi_1^{\LL}(U_{\bar{s}_2},a_2) \quad (\text{\resp~}\pi\colon
\pi_1^\tame (U_{\bar{s}_1},a_1)\to \pi_1^\tame
(U_{\bar{s}_2},a_2)).
$$
Changing the class of paths from~$a_1$ to $a_2$ amounts to
modifying $\pi$ by an inner automorphism of
$\pi_1^{\LL}(X_{\bar{s}_2},a_2)$ (\resp of
$\pi_1^\tame(X_{\bar{s}_2},a_2)$). We call \emph{specialization
morphism for the fundamental group}
associated with the morphism $\bar{s}_1\to \bar{s}_2$, and we denote
simply by
$$
\pi\colon \pi_1^{\LL}(X_{\bar{s}_1})\to
\pi_1^{\LL}(X_{\bar{s}_2}) \quad \text{(\resp~} \pi\colon
\pi_1^{\tame}(X_{\bar{s}_1})\to \pi_1^{\tame}(X_{\bar{s}_2}))
$$
any one of the morphisms defined above.

\begin{lemma}
\label{XIII.2.11}
Let $f\colon X\to S$ be a proper morphism of
finite presentation, $D$ a divisor on~$X$ with normal
crossings relative to~$S$, $Y=\Supp D$, $U=X-Y$, $i\colon
U\to X$ the canonical morphism, $\bar{s}_1\to
\bar{s}_2$ a specialization morphism of geometric points
of~$S$, $y_1$ a geometric point of
$Y_{\bar{s}_1}$, $y_2$ a geometric point of
$Y_{\bar{s}_2}$, such that the projection $z_1$ of~$y_1$ on~$X$ is a
generization of the projection $z_2$ of~$y_2$. Let $I_{y_1}$
be an inertia subgroup of~$\pi_1^{\tame}(\bar{U}_{\bar{s}_1})$ at
$y_1.$ Then the image of~$I_{y_{1}}$ under the
specialization morphism
$$
\pi\colon \pi_1^{\tame}(U_{\bar{s}_1})\to
\pi_1^{\tame}(U_{\bar{s}_2})
$$
is
% original p. 392
an inertia subgroup of~$\pi_1^{\tame}(U_{\overline{s}_2})$
at~$y_2$.
\end{lemma}

Indeed, let $\overline{X}$ (\resp $\widetilde{X}$) be the strict
localization of~$X$ at $y_2$ (\resp at $y_1$), $\overline{U}=U \times_X
\overline{X}$ (\resp $\widetilde{U}=U\times_X \widetilde{X}$). We have a
canonical morphism $\widetilde{U}\to \overline{U}$, and it
follows from~\Ref{XIII.1.10} that we have a commutative diagram
$$
\xymatrix{ \pi_1^{\tame} \left(\widetilde{U}_{\overline{s}_1} \right)
\ar[r]^{\pi '} \ar[d] & \pi_1^{\tame} \left(
\overline{U}_{\overline{s}_2} \right) \ar[d] \\ \pi_1^{\tame} \left(
U_{\overline{s}_1} \right) \ar[r]^{\pi } & \pi_1^{\tame} \left(
U_{\overline{s}_2} \right) }
$$
where $\pi'$ is the composite of the canonical morphism
$\pi_1^{\tame}(\widetilde{U}_ {\overline{s}_1})\to
\pi_1^{\tame}(\overline{U}_{\overline{s}_1})$ and of the
specialization morphism. Since
$\pi_1^{\tame}(\widetilde{U}_{\overline{s}_1})$
(\resp $\pi_1^{\tame}(\overline{U}_{\overline{s}_1})$)
is an inertia group of
$\pi_1^{\tame}(U_{\overline{s}_1})$ at~$y_1$ and (\resp of~$\pi_1^{\tame}(U_{\overline{s}_2})$ at~$y_2$), it suffices to
prove that $\pi'$ is surjective. But this follows from
the expression obtained in~\Ref{XIII.5.6}.

\begin{corollary}
\label{XIII.2.12}
Let $X$ be a proper and smooth connected curve of genus $g$ over a
separably closed field $k$ of characteristic $p\geq
0$. Let $U$ be the open set obtained by removing from $X$ $n$ distinct
closed points $a_1,\dots,a_n$. Then the tamely ramified
fundamental group $\pi_1^{\tame}(U)$ \eqref{XIII.2.1.3}
can be generated by $2g+n$ elements $x_i,
y_i,\sigma_j$, with $1\leq i\leq g$, $1\leq j\leq n$, such that
$\sigma_j$ is a generator of an inertia group
corresponding to $a_j$, and such that we have the relation
\begin{equation*}
\label{eq:XIII.2.12.*}
\tag{$*$} \prod_{1\leq i\leq
g}(x_iy_ix^{-1}_iy^{-1}_i)\cdot\prod_{1\leq j\leq n}\sigma_j=1.
\end{equation*}
For every finite group $G$ \emph{of order prime to} $p$,
generated by elements $\overline{x}_i,
\overline{y}_i,\overline{\sigma}_j$ satisfying the relation
\eqref{eq:XIII.2.12.*}, there exists an \'etale covering of~$U$,
with group $G$, corresponding to a homomorphism
$\pi_1^{\tame}(U)\to G$ which
% original p. 393
sends $x_i, y_i,\sigma_j$ to~$\overline{x}_i,
\overline{y}_i,\overline{\sigma}_j$ respectively. In other
words, if $p'$ denotes the set of prime numbers distinct
from~$p$, $\pi_1^{p'}(U)$ is the pro-$p'$-group generated by the
generators $x_i, y_i,\sigma_j$ bound by the single relation~\eqref{eq:XIII.2.12.*}.
\end{corollary}

\subsubsection*{Proof}
We may assume $k$ algebraically closed. Suppose first $k$ of
characteristic zero. There then exists an algebraically closed
subextension $k'$ of~$k$, of finite transcendence degree
over~$\QQ$, such that $X$ comes by extension of scalars
from a proper and smooth curve $X'$ defined over~$k'$, and we may
assume that the points $a_1,\dots,a_n$ come from rational points
$a'_1,\dots,a'_n$ of~$X'$. Since $k'$ is of finite transcendence
degree over~$\QQ$, we can find an embedding of~$k'$ into the field
of complex numbers $\CC$; let $\widetilde{U}=U'\times_{k'}\CC.$
Let $k''$ be an algebraically closed extension of~$k'$ such that
we have $k'$-morphisms of~$k$ and of~$\CC$ into $k''$. If
$g'\colon U'\to k'$ is the structural morphism, and if
$\cal{F}$ is a finite constant sheaf of groups on~$U''$, it
follows from~\Ref{XIII.2.9} that the specialization morphisms
$$
(\R^1g'_*\cal{F}')_{\CC}\to(\R^1g'_*\cal{F}')_{k''}
\from(\R^1g'_*\cal{F}')_k
$$
are isomorphisms. In terms of fundamental groups, this shows
that we have an isomorphism, defined up to inner
automorphism
$$
\pi_1(U)\to\pi_1(\widetilde{U}),
$$
and it is clear that this isomorphism transforms an inertia group
relative to a point of~$X'-U'$ into an inertia group relative to the
same point. We may therefore assume that $k=\CC$. In this
last case it follows from Riemann's existence theorem
(XII~\Ref{XII.5.2}) that the fundamental group $\pi_1(U)$ is none other
than the completion, for the topology of subgroups of finite
index, of the fundamental group of the analytic space associated with~$U$. Now the latter can be computed by transcendental means \cite[ch\ptbl 7
\S 47]{XIII.3};
% original p. 394
it can be generated by $2g+n$ elements $x_i,
y_i,\sigma_j$ such that $\sigma_j$ is the image of a generator
of the local fundamental group $\pi_1(D_j)$ of a small disk centered
at $a_j$, \ie a generator of an inertia group
corresponding to the point $a_j$, these elements satisfying
the single relation \eqref{eq:XIII.2.12.*}.

If now $k$ is of
characteristic $p>0$, we can find a complete discrete valuation
ring $A$, with residue field $k$, with field of
fractions $K$ of characteristic zero, and a connected
scheme $X_1$, proper and smooth over~$S=\Spec A$, such that we have
$X_1\times_S \Spec k\simeq X$ (III~\Ref{III.7.4}). The points $a_j$ then
lift to sections $s_j$ of~$X_1$ over~$S$;
let $Y_{1j}$ be the reduced closed subscheme with
underlying space $s_j(S)$, $Y_1$ the union of the $Y_{1j}$,
$U_1=X_1-Y_1$, $g_1\colon U_1\to S$ the structural
morphism. Let $\overline{K}$ be an algebraically closed extension
of~$K$, $\overline{U}= U_1\times_S\overline{K}$. If $C$ is a
finite constant group, it follows from~\Ref{XIII.2.8} that the
specialization morphism
$$
(\R^1g_{1*}C_{U_1})_k\to(\R^1g_{1*}C_{U_1})_{\overline{K}}
$$
is injective and even bijective if $C$ is of order prime to
$p$. Now this means, in terms of fundamental groups, that the
specialization morphism \eqref{XIII.1.10}
$$
\pi \colon \pi_1(\overline{U})\to\pi^{\tame}_1(U)
$$
is surjective, and that the specialization morphism
$$
\pi^{p'}_1(\overline{U})\to\pi^{p'}_1(U)
$$
is bijective. Finally, if $x_i, y_i,\sigma_j$ are
generators of~$\pi_1(\overline{U})$ such that $\sigma_j$ is
a generator of an inertia group corresponding to the point
$b_j=Y_{1j}(\overline{K})$ of~$\overline{X}$, then,
by~\Ref{XIII.1.11}, $\pi(\sigma_j)$ is a generator
of an inertia group corresponding to $a_j$, which completes the
proof.

\begin{remarkMR}
\label{XIII.2.13}
Let $k$ be an algebraically closed field of characteristic $p>0$, $X$~a proper smooth algebraic curve over~$k$, connected, of genus $g$, $U$~an affine open set of~$X$, the complement of~$r\geq1$ rational points
of~$X$. We have at our disposal the fundamental group $\pi_1(U)$, its quotient
$\pi_1^{\tame}(U)$ which classifies the finite \'etale coverings of~$U$
tamely ramified at the points of~$X-U$, and the quotient
$\pi_1^{p'}(U)$ of~$\pi_1^{\tame}(U)$ which classifies the Galois
\'etale coverings of~$U$ with Galois group of order prime to~$p$.

In (Coverings of algebraic curves, Amer.\ J.~Math.\ \textbf{79}
(1957), p\ptbl825--856) S\ptbl Abhyankar formulated a certain number of conjectures on the structure of the finite groups that are Galois groups of finite
\'etale coverings of~$U$.

The conjectures concerning the finite groups that are Galois groups
of connected \'etale coverings of~$U$ of order prime to~$p$
are proved and made more precise in Corollary \Ref{XIII.2.12}. Before tackling the finite quotients of~$\pi_1(U)$, let us begin by giving some
indications about the ``size'' of this group. Let $X = \Spec(A)$. It is
known that $\Hom_{\mathrm{cont}}(\pi_1(U), \ZZ/p\ZZ)$ is described by
Artin-Schreier theory (cor\ptbl XI~\Ref{XI.6.9}). This group is
isomorphic to $A/\wp(A)$ where $\wp$ is the map from~$A$ to $A$ which
sends $a$ to~$a^p-a$. Suppose first that $U$ is the affine
line $\AA^1$, with ring $k[T]$. Let~$E$ be the set of
elements of~$k[T]$ of the form $\sum_i a_iT^i$, with $a_i = 0$
when $p$ divides~$i$. It is immediate that the composite map
$E\to k[T]\to k[T]/\wp k[T]$ is bijective. Consequently the coefficients
$a_i$, $(i,p)=1$, behave like coordinates of a space which
parametrizes the cyclic coverings of degree
$p$ of the affine line. In particular, one deduces that $\pi_1(\AA^1)$ is not
topologically of finite type and that, if one makes an extension
$k\to k'$ of algebraically closed fields, the natural morphism $\pi_1(\AA^1\times_kk')\to\pi_1(\AA^1)$, which is surjective, is not bijective,
contrary to the proper case. In the general case, the curve $U$ is
realized as a scheme finite over the affine line and one concludes that the
same phenomena occur for $\pi_1(U)$ (and moreover, more
generally, for $\pi_1(V)$ for every connected affine $k$-scheme $V$
of finite type, of dimension $\geq1$).

This being so, if $G$ is a finite group, let us denote by $G^{(p')}$ the largest
quotient group of~$G$ of order prime to~$p$. For a finite group
$G$ to be a topological quotient of~$\pi_1(U)$, it is necessary that $G^{(p')}$
be a topological quotient of~$\pi_1^{p'}(U)$, a condition which
one knows in principle how to answer thanks to Corollary \Ref{XIII.2.12}. In the
aforementioned article, Abhyankar conjectures that this necessary condition is
also sufficient. This conjecture was proved by M\ptbl Raynaud
in the case of the affine line and by D\ptbl Harbater in the general case
(Invent. Math.\ \textbf{116} (1994), p\ptbl425--462 and \textbf{117}, p\ptbl1--25). For example,
in the case of the affine line $\pi_1^{p'}(\AA^1)=1$ and one
concludes that a finite group $G$ is the Galois group of a connected
\'etale covering of~$\AA^1$ if and only if $G^{(p')} = 1$,
that is, if and only if $G$ is generated by its
Sylow $p$-subgroups. Thus every finite simple group of order a
multiple of~$p$ will do.
\end{remarkMR}


\section{Cohomological properness and generic local acyclicity}
\label{XIII.3}
% original p. 395

\begin{theorem}
\label{XIII.3.1}
Let $S$ be an irreducible scheme with generic
point~$s$, $X$ and $Y$ two $S$-schemes of
finite presentation, $f\colon X \to Y$ an $S$-morphism. For every
$S$-scheme~$S'$, we denote by $Y'$, $X'$, etc.\ the inverse image of~$Y$,
$X$, etc.\ under the morphism $S' \to S$. We have the following
properties:

\begin{enumerate}
\item[1)]
{\upshape a)} One can find a nonempty open set~$S'$ of~$S$ such that, for
every finite constant sheaf of sets $F'$ on~$X'$, $f'_*F'$ is
constructible, and $(F',f')$ is cohomologically proper
relative to~$S'$ in dimension $\le 0$.

{\upshape b)} Let $F$ be a constructible sheaf of sets on~$X$. Then
one can find a nonempty open set~$S'$ of~$S$ (depending on~$F$)
such that $f'_*F'$ is constructible and $(F',f')$ is
cohomologically proper relative to~$S'$ in dimension $\le 0$.

\item[2)] Suppose that the schemes of finite type of dimension $\le
\dim X_s$ over an algebraic closure $\overline{k}$ of~$\kres(s)$
are strongly desingularizable \textup{(SGA~5 I~3.1.5)}. Then we have
moreover the following properties:

{\upshape a)} One can find a nonempty open set~$S'$ of~$S$ such that, for
every finite constant sheaf of groups $F'$ on~$X'$, of order prime
to the residue characteristics of~$S$, if $\Phi'$ is the
stack of torsors under $F'$, $f'_*\Phi'$ is $1$-constructible, and
$(F',f')$ is cohomologically proper relative to~$S'$ in
dimension $\le 1$.

{\upshape b)} Let $\LL$ be the set of prime numbers distinct from the
residue characteristics of~$S$ and $\Phi$ a
$1$-constructible ind-$\LL$-stack on~$X$ \eqref{XIII.0}, such that,
for every scheme~$X_1$ \'etale over~$X$ and for every pair
of objects $x, x_1$ of~$\Phi_{X_1}$, the sheaf $\SheafHom
_{X_1}(x,x_1)$ is constructible. Suppose moreover $S$ locally
noetherian. Then one can find a nonempty open set~$S'$ of~$S$
such that $f'_*\Phi'$ is $1$-constructible, that, for every pair
of objects $y, y_1$ of a fiber $(f'_*\Phi')_{Y_1}$,
$\SheafHom_{Y_1}(y,y_1)$ is constructible,
% original p. 396
and that $(\Phi',f')$ is cohomologically proper relative to~$S'$ in dimension $\le 1$.
\end{enumerate}
\end{theorem}

\subsubsection*{Proof}
We may assume $S$ affine; by SGA~4 VIII~1.1, we may
assume $S$ integral; finally, by passage to the limit, we may
assume that $S$ is the spectrum of an algebra of finite type over~$\ZZ$; in particular $S$ is then noetherian. Since the question
is local on~$Y$, we may assume $Y$ affine. Moreover, to prove the
theorem, it suffices to do so after a
finite extension $S'\to S$, where $S'$ is an integral scheme
and where $S'\to S$ is a composite of \'etale morphisms and of
finite radicial surjective morphisms.

\Subsection*{1) Case of constant sheaves of sets}

1) 1. Reduction to the case where $X$ is normal over~$S$. Let
$X_{1\overline{s}}$ be the normalization of~$(X_{\overline{s}})_{\red}$;
after restricting $S$ to a nonempty open set and making
a radicial extension of~$S$ if necessary, we may assume that
$X_{1\overline{s}}$ comes from a scheme $X_1$ normal over~$S$, and
that the morphism $X_{1\overline{s}}\to X_{\overline{s}}$ comes from a
finite surjective morphism $p\colon X_1\to X$ (EGA~IV 8.8.2 and
9.6.1). Suppose the theorem proved for
$fp$. After restricting $S$ to an open set if necessary, we may assume
that, for every constant sheaf of sets $F$ on~$X$, $(p^*F,fp)$
is cohomologically proper relative to~$S$ in dimension $\le
0$ and that $f_*p_*(p^*F)$ is
constructible. By~\Ref{XIII.1.9}, $(p_*p^*F,f)$ is then
cohomologically proper relative to~$S$ in dimension $\le
0$. The morphism
$$
F \to p_* p^* F=G
$$
is a monomorphism. It already follows from this, $f_*F$ being
a subsheaf of~$f_*G$, that $f_*F$ is constructible (SGA~4
IX~2.9~(ii)) and that $(F,f)$ is cohomologically proper relative
to~$S$ in dimension $\le -1$.

Let $X_2=X_1 \times_{X} X_1$, $q \colon X_2 \to X$ the
canonical morphism. By \Ref{XIII.1.11}~1),
% original p. 397
we have an exact sequence
$$
\xymatrix@C=.5cm{F \ar[r] & G \ar@<2pt>[r]\ar@<-2pt>[r] & q_*q^*F }.
$$

By what has just been proved, applied to
$fq$ instead of~$f$, we may assume, after restricting $S$
to a nonempty open set if necessary, that, for every constant sheaf of sets
$F$ on~$X$, $(q^*F,fq)$ is cohomologically proper relative
to~$S$ in dimension $\le -1$, hence that $(q_*q^*F,f)$ is
cohomologically proper relative to~$S$ in dimension $\le
-1$. It then follows from \Ref{XIII.1.13}~1) that $(F,f)$ is
cohomologically proper relative to~$S$ in dimension $\le 0$.

1) 2. Reduction to the case where $X$ is normal affine over~$S$.

Let $U_s$ be an affine open set of~$X_s$ dense in $X_s$. After
restricting $S$ to a nonempty open set if necessary, we may assume that $U_s
\to X_s$ lifts to an open immersion $i\colon U\to X$,
schematically dominant relative to~$S$ (EGA~IV
8.9.1). Since the morphism $X \to S$ is normal, we have by
SGA~2 XIV~1.18:
$$
\prof \et_{S- U} (X) \ge 2 \quoi ;
$$
consequently, for every constant sheaf $F$ on~$X$, the canonical
morphism
$$
F \to i_*i^*F
$$
is an isomorphism. It follows that, if we suppose the
theorem proved for $fi$ and $i$, then, after
restriction of~$S$ to a nonempty open set, $(i^*F,i)$ and $(i^*F,fi)$
are cohomologically proper relative to~$S$ in dimension
$\le 0$. The same therefore holds for~$(F,f)$ (\Ref{XIII.1.6}~2)). Since moreover $f_* F=(fi)_*(i^*F)$ is constructible, this
completes the reduction.

1) 3. End of the proof.

We may assume $S$ normal (EGA~IV 7.8.3). We can find a
compactification of~$X_s$:
% original p. 398
$$
\xymatrix{ X_s \ar[r]^{j_s} \ar[d]_{f_s} & P_s \ar[ld]^{g_s} \\ Y_s }
$$
where $j_s$ is a dominant open immersion and $g_s$ a
proper morphism; after making a radicial extension of~$\kres(s)$ and
replacing $P_s$ by its normalization, which does not change $X_s$, we
may assume $P_s$ geometrically normal. After
restricting $S$ to a nonempty open set and making a surjective
radicial extension if necessary, we may assume that the diagram above
comes from a diagram
$$
\xymatrix{ X \ar[r]^{j} \ar[d]_{f} & P \ar[ld]^{g} \\ Y }
$$
where $P$ is a scheme normal over~$S$, $j$ an open immersion
schematically dominant relative to~$S$ and $g$ a
proper morphism (EGA~IV 6.9.1, 9.9.4 and 9.6.1). For every finite
constant sheaf of sets $F$ on~$X$ with value $C$, $j_*F$ is the
constant sheaf with value $C$ (SGA~4 2.14.1), and the same holds
after any change of base $S' \to S$; it
follows that $(F,j)$ is cohomologically proper relative
to~$S$ in dimension $\le 0$. The same therefore holds for~$(F,f)$,
since $g$ is proper \eqref{XIII.1.8}. Since $g$ is proper
$f_*F=g_*C_P$ is constructible, which completes the
proof of 1)\kern2pt a).

\Subsection*{2) Case of a constructible sheaf of sets}

Let $F$ be a constructible sheaf of sets on~$X$. By
SGA~4 IX 2.14~(ii), we can find a finite family of morphisms
$p_i \colon Z_i \to X$, and, on each $Z_i$, a finite constant sheaf of sets
$C_i$, such that we have a monomorphism
% original p. 399
$$
j \colon F \to \prod_{i} {p_i}_* C_i =G \quoi.
$$
By~1)\kern2pt a),
we may assume, after restricting $S$ to a nonempty
open set if necessary, that the $(C_i,fp_i)$ are cohomologically proper relative
to~$S$ in dimension $\le 0$, and that the $f_*{p_i}_*C_i$ are
constructible. We already conclude from this that $(G,f)$ is
cohomologically proper relative to~$S$ in dimension $\le 0$
\eqref{XIII.1.9}, hence that $(F,f)$ is cohomologically proper
relative to~$S$ in dimension $\le -1$, and that $f_*F$ is
constructible. Let~$K$ be the amalgamated sum $K= G \coprod_{F} G$;
since $F$ and $G$ are constructible, so is~$K$. We
therefore conclude from the preceding that, after restricting
$S$ to a nonempty open set if necessary, we may assume that $(K,f)$ is
cohomologically proper relative to~$S$ in dimension $\le
-1$. It then follows from \Ref{XIII.1.13}~1) that $(F,f)$ is
cohomologically proper relative to~$S$ in dimension $\le 0$.

\Subsection*{3) Case of constant sheaves of groups}

If $F$ is a constant sheaf of groups on~$X$, we denote by $\Phi$ the
stack of torsors under $F$.

3) 1. Let us first show that, after restricting $S$ to a
nonempty open set if necessary, for every constant sheaf of groups $F$ on~$X$,
of order prime to the residue characteristics of~$S$,
$(F,f)$ is cohomologically proper relative to~$S$ in
dimension $\le 0$, and that $f_*\Phi$ is constructible.

For this we reduce to the case where $X$ is smooth over~$S$. After
making a finite extension of~$\kres(s)$ if necessary, which is permissible since one
can regard it as the composite of an \'etale extension and
of a radicial extension, we can find a proper
surjective morphism $p_s \colon X_{1s}\to X_s$, where $X_{1s}$ is a
scheme smooth over~$S$ of the same dimension as $X_s$, and, after
restricting $S$ to a nonempty open set if necessary, we may assume that
$p_s$ comes from a proper surjective morphism $p \colon X_1 \to X$,
where $X_1$ is a scheme smooth over~$S$
% original p. 400
(EGA~IV 9.6.1 and 12.1.6). Let $X_2=X_1 \times_X X_1$, $q \colon
X_2 \to X$ the canonical morphism. We have an exact diagram of stacks
on~$X$
$$
\xymatrix@C=.5cm{\Phi \ar[r] & p_*p^* \Phi \ar@<2pt>[r]\ar@<-2pt>[r] &
q_*q^*\Phi }
$$
(\Ref{XIII.1.11}~2)). The theorem being assumed
proved in the smooth case, we see first of all that we may
assume $f_*p_*p^*\Phi$ constructible; hence so is
$f_*\Phi$ (\Ref{XIII.3.1.1} below). Moreover, by \Ref{XIII.1.6}~2),
we may assume $(p_*p^*\Phi,f)$ cohomologically proper
relative to~$S$ in dimension $\le 0$; it follows that
$(\Phi,f)$ is cohomologically proper relative to~$S$ in
dimension $\le -1$. We may therefore assume that $(q_*q^*\Phi,f)$ is
cohomologically proper relative to~$S$ in dimension $\le -1$,
and this implies that $(\Phi,f)$ is cohomologically proper
relative to~$S$ in dimension $\le 0$ (\Ref{XIII.1.12}~1)).

We then reduce as in 1)\kern2pt2 to the case where $X$ is smooth
and affine over~$S$. Let then
$$
\xymatrix{ X \ar[r]^{i} \ar[d]_{f} & P \ar[ld]^{q} \\ Y }
$$
be a compactification of~$X$, where $i$ is a dominant open
immersion and $q$ a proper morphism. Since we have $\dim P_s=\dim X_s$,
we can apply the hypothesis of resolution of
singularities to $P_{\overline s}$. After making an
\'etale extension and a radicial extension of~$S$ if necessary, we can
find a proper morphism $r \colon Z \to P$, where $Z$ is smooth
over~$S$, $r^{-1}(X)\simeq X$, and where $r^{-1}(X)$ is the
complement in $Z$ of a divisor with normal crossings
relative to~$S$. Every torsor under $F$ is then
tamely ramified on~$Z$ relative to~$S$
(\Ref{XIII.2.3}~b)). It therefore follows from~\Ref{XIII.2.7} that $(F,f)$
is cohomologically proper relative to~$S$ in dimension $\le
0$, which proves our assertion.

3) 2. Reduction to the case where $X$ is smooth over~$S$.
% original p. 401
\par
After making a finite extension of~$\kres(s)$ if necessary, we can find a
proper surjective morphism $p_s\colon X_{1s}\to X$, where~$X_{1s}$
is a scheme smooth over~$s$, and we may assume that $p_s$
comes from a proper surjective morphism $p\colon X_1 \to X$, where
$X_1$ is smooth over~$S$. Suppose the theorem
proved for $fp$ and let us prove it for $f$. Let $F$ be a
finite constant sheaf of groups on~$X$, of order prime to the
residue characteristics of~$S$ and $\Phi$ the stack of
torsors under $F$. Let $X_2=X_1 \times_X X_1$, $X_3=X_1 \times_X
X_1 \times_X X_1$, $q\colon X_2 \to X$, $r\colon X_3 \to X$ the
canonical morphisms. By \Ref{XIII.1.11}~2), we have an
exact diagram of stacks
$$
\xymatrix@C=.5cm{\Phi \ar[r] & p_*p^*\Phi \ar@<2pt>[r]\ar@<-2pt>[r] &
q_*q^*\Phi \ar@<4pt>[r] \ar[r] \ar@<-4pt>[r] & r_*r^*\Phi}\quoi.
$$
To prove that $(\Phi,f)$ is cohomologically proper relative
to~$S$ in dimension $\le 1$, it suffices to show that the same
holds for~$(p_*p^*\Phi,f)$, that $(q_*q^*\Phi, f )$ is
cohomologically proper relative to~$S$ in dimension $\le 0$
and that $(r_*r^*\Phi,f)$ is cohomologically proper relative
to~$S$ in dimension $\le -1$ (\Ref{XIII.1.12}~2)). By~3)\kern.3em 1
above we may assume that, for every finite constant sheaf of
groups~$F$, $(q^*\Phi,fq)$, $(q^*\Phi,q)$, $(r^*\Phi,fr)$, $(r^*\Phi,r)$
are cohomologically proper relative to~$S$ in dimension
$\le 0$. It then follows from \Ref{XIII.1.6}~2) that
$(q_*q^*\Phi,f)$ and $(r_*r^*\Phi,f)$ are cohomologically proper
relative to~$S$ in dimension $\le 0$. The theorem
being assumed proved in the smooth case, $(p^*\Phi,
fp)$ and $(p^*\Phi,p)$, hence also $(p_*p^*\Phi,f)$ (\Ref{XIII.1.6}~2)),
are cohomologically proper relative to~$S$ in dimension
$\le 1$. This indeed shows that $(\Phi,f)$ is cohomologically proper
relative to~$S$ in dimension $\le 1$.

Moreover $f_*p_*p^*\Phi$ is $1$-constructible by hypothesis;
by~\Ref{XIII.3.1} we may assume that $f_*q_*q^*\Phi$ is
constructible; it therefore follows from~\Ref{XIII.3.1.1} below that
$f_*\Phi$ is $1$-constructible.

3) 3. Reduction to the case where $X$ is smooth affine over~$S$.
% original p. 402

By 3)\kern2pt2, we may assume $X$ smooth over~$S$. Let $(U_i)_{i\in
I}$ be a finite covering of~$X$ by affine open sets, and let $X_1$
be the direct sum of the $U_i$, $p\colon X_1 \to X$ the canonical
morphism. Since $p$ is an effective descent morphism for the
category of \'etale sheaves of finite type over
variable schemes, we see as in 3)\kern2pt2 that, if we suppose the
theorem proved for the $U_i$, \ie for $X_1$, it
is also true for~$X$.

3) 4. Case where $X$ is smooth affine over~$S$.

We see as in 3)1 that, after restricting $S$ to a
nonempty open set and making a surjective \'etale extension and a
surjective radicial extension of~$S$ if necessary, we can find a
commutative diagram
$$
\xymatrix{ X \ar[r]^{i} \ar[d]_{f} & P \ar[ld]^{q} \\ Y }
$$
where $P$ is a scheme smooth over~$S$, $X$ the complement
in $P$ of a divisor with normal crossings relative to~$S$ and $g$ a proper
morphism. If $F$ is a constant sheaf
of order prime to the residue characteristics of~$S$, every
torsor under $F$ is tamely ramified on~$P$
relative to~$S$ (\Ref{XIII.2.3}~b)). The fact that $(F,f)$ is
cohomologically proper relative to~$S$ in dimension $\le 1$
and that $f_*\Phi$ is constructible then follows
from~\Ref{XIII.2.7}.

\Subsection*{4) Proof of 2)\kern2pt b)}

4) 1. Case where $\Phi$ is a gerbe.

One can find a surjective \'etale morphism of finite type $p
\colon X_1 \to X$ such that $p^*\Phi$ is a trivial gerbe. By
descent, as in 3)\kern2pt2, one sees that it suffices to prove the
theorem for $X_1$, $X_1\times_X X_1$ and
% original p. 403
$X_1\times_X X_1\times_X X_1$. One is therefore reduced to the case where
$\Phi$ is the gerbe of torsors under a constructible sheaf of groups
$F$ whose fibers are of order prime to the
residue characteristics of~$S$.

By SGA~4 IX~2.14, one can find a finite family of
finite morphisms $p_i \colon Z_i \to X$ and, for each $i$, a
finite constant sheaf of groups $C_i$, of order prime to the
residue characteristics of~$S$, such that one has a
monomorphism
$$
j\colon F \to \prod_{i} p_{i*}C_i=G \quoi.
$$
Let $\Phi_i$ be the stack of torsors under $C_i$ and $\Psi$ the stack of
torsors under $G$. It follows from 2)\kern2pt a) that, at the cost of
restricting $S$ to a nonempty open subset, one may assume that the
$(C_i,fp_i)$ are cohomologically proper relative to~$S$ in
dimension $\le 1$, and that the stacks $f_*p_{i*}\Phi_i$ are
$1$-constructible. It then follows from~\Ref{XIII.1.9} that the
$(p_{i*}C_i,f)$ are cohomologically proper relative to~$S$
in dimension $\le 1$; hence the same holds for~$(G,f)$. Moreover,
since $p_{i*}\Phi_i$ is equivalent to the stack of torsors under the
group $p_{i*}C_i$ (SGA~4 VIII~5.8), one sees that $f_*\Psi$ is
$1$-constructible.

Since $\R^1f_*G$ is constructible, one can find a sheaf
representable by an \'etale $Y$-scheme of finite type $T$ and
an epimorphism
$$
a \colon T \to \R^1f_*G
$$
(SGA~4 IX~2.7); moreover one may assume that the image of the identity
section of~$T(T)$ is defined by a torsor $Q$ on~$X\times_Y
T=X_T$, with group $G|X_T$. Let $f_T \colon X_T \to T$ be the canonical
morphism, $F_T=F|X_T$, etc. By 1)\kern2pt b) one may assume,
at the cost of restricting $S$ to a nonempty open subset, that
$(Q/F_T,f_T)$ is cohomologically proper relative to~$S$ in
dimension $\le 0$ and that $f_{T*}(Q/F_T)$ is constructible. It
then follows from~\Ref{XIII.3.1.2} that $f_*\Psi$ is constructible.

Let us show that $(F,f)$ is cohomologically proper relative to
% original p. 404
$S$ in dimension $\le 1$. By \Ref{XIII.1.13}~2) it suffices to
prove that, for every scheme $Y_1$ \'etale over~$Y$ and every
torsor $Q_1$ on~$X_1=X\times_Y Y_1$, if $f_1 \colon X_1 \to Y_1$ is
the canonical morphism, then $(Q_1/F_1,f_1)$ is cohomologically
proper relative to~$S$ in dimension $\le 0$. Now, by
definition of~$T$, $Q_1$ is, locally for the \'etale topology
of~$Y_1$, an inverse image of~$Q$, which proves our
reduction.

4) 2. General case.\par
Using lemma~\Ref{XIII.6.1.1}, 4)\kern2pt1 and 1)\kern2pt a), one sees that,
at the cost of restricting $S$ to a nonempty open subset, one may
assume $S(f_*\Phi)$ constructible and $(S\Phi,f)$ cohomologically
proper relative to~$S$ in dimension $\le 0$. One can then
find a sheaf representable by an \'etale $Y$-scheme
of finite type $T$ and an epimorphism
$$
a \colon T \to S(f_*\Phi)\quoi.
$$
One takes up again the notation of 4)\kern2pt1, and one sets moreover $Z=T\times_Y T$,
$X_Z=X\times_Y Z$, and one denotes by $f_Z \colon X_Z \to Z$ the canonical
morphism. One may assume $T$ chosen so that the image $q$ of the
identity section of~$T(T)$ under $a$ is defined by an object~$p$
of~$(f_*\Phi)_T=\Phi_{X_T}$. Let $p_1$ and $p_2$ (\resp~$q_1$ and
$q_2$) be the inverse images of~$p$ (\resp~$q$) under the two projections
of~$Z$ into $T$. One may assume, at the cost of restricting $S$ to
a nonempty open subset, that $(\SheafAut_{X_T}(p),f_T)$ is
cohomologically proper relative to~$S$ in dimension $\le 1$,
that $(\SheafHom_{X_Z}(p_1,p_2),f_Z)$ is cohomologically proper
relative to~$S$ in dimension $\le 0$, and that
$f_{T*}(\SheafAut_{X_T}(p))$, $f_{Z*}(\SheafHom_{X_Z}(p_1,p_2))$ are
constructible ( a)\kern2pt1) and 4)\kern2pt1)).

From this one first deduces that, for every scheme $Y_1$ \'etale over~$Y$ and every pair of objects $y$, $y_1$ of~$(f_*\Phi)_{Y_1}$, the
sheaf $\SheafAut_{Y_1}(y)$ (\resp $\SheafHom_{Y_1}(y,y_1)$) is
constructible, such a sheaf being, locally for the \'etale
topology of~$Y_1$, an inverse image of~$f_{T*}(\SheafAut_{X_T}(p))$
(\resp of~$f_{Z*}(\SheafHom_{X_Z}(p_1,p_2))$).

It remains to prove that $(\Phi,f)$ is cohomologically proper
relative
% original p. 405
to~$S$ in dimension $\le 1$. For this it suffices to show that,
for every $S'$-scheme $S$ and every geometric point
$\overline{y}'$ of~$Y$, if one denotes by $\overline{Y}$ the strict
localization of~$Y$ at $y'$, $\overline{X}=X \times_Y \overline{Y}$, etc.,
then the canonical functor
$$
\overline{\varphi} \colon \Phi(\overline{X}) \to
\Phi'(\overline{X'})
$$
is an equivalence of categories.

Let us show that $\overline{\varphi}$ is fully faithful. Let
$x$, $y \in \Phi(\overline{X})$, let $x'$, $y'$ be their images in
$\Phi'(\overline{X'})$, and let us show that the canonical morphism
\begin{equation*}
\label{eq:XIII.4.*} \tag{$*$} {\Hom_{\overline{X}}(x,y)
\to \Hom_{\overline{X'}}(x',y') }
\end{equation*}
is bijective. By definition of~$T$ there exist two morphisms from
$\overline{Y}$ to $T$ such that $x$ and~$y$ are the inverse images
of~$p$ under these two morphisms. This amounts to saying that there is a
morphism $\overline {Y} \to Z$, whence a morphism $h\colon
\overline{X} \to X_Z$ such that one has
$$
h^*(p_1)=x \qquad h^*(p_2)=y.
$$
Consequently one has a canonical isomorphism
$$
\SheafHom_{\overline{X}}(x,y)=h^*(\SheafHom_{X_Z}(p_1,p_2) )\quoi.
$$
but, taking into account the fact that $(\SheafHom_{X_Z}(p_1,p_2),f_Z)$ is
cohomologically proper relative to~$S$ in dimension $\le 0$,
one sees that the same holds for
$(\SheafHom_{\overline{X}}(x,y),\overline{f})$, which proves that the
morphism \eqref{eq:XIII.4.*} is bijective.

Let us show that $\overline{\varphi}$ is essentially surjective. Let
$x'\in \Phi'(\overline{X'})$. Since $(S\Phi,f)$ is cohomologically
proper relative to~$S$ in dimension $\le 0$, the canonical
morphism
$$
\H^0(\overline{X},\overline{S\Phi}) \to
\H^0(\overline{X'},\overline{S\Phi}')
$$
is bijective. Let $G'$ be the maximal subgerbe of~$\Phi'$
generated by $x'$; there then exists a maximal subgerbe $G$ of
$\overline{\Phi}$ whose inverse image on~$\overline{X'}$ is
$G'$. By 4)\kern2pt1, $(G,\overline{f})$ is cohomologically proper
relative to~$S$ in dimension $\le 1$; consequently the canonical
functor
$$
G(\overline{X}) \to G'(\overline{X'})
$$
% original p. 406
is an equivalence of categories, which proves the existence
of an element $x$ of~$\Phi(\overline{X})$ whose image in
$\Phi'(X')$ is isomorphic to $x'$, and completes the
proof of the theorem.

\begin{sublemma}
\label{XIII.3.1.1}
Let $S$ be a locally noetherian scheme and
$$
\xymatrix@C=.5cm{ \Phi \ar[r]^p & \Phi_1 \ar@<2pt>[rr]^{p_1,p_2}
\ar@<-2pt>[rr] && \Phi_2 }
$$
an exact diagram of stacks on~$S$ \eqref{XIII.1.10.1}. If $\Phi_1$
is constructible, so is $\Phi$. If, for every
scheme $S'$ \'etale over~$S$ and every pair
of objects
$x_1,~y_1$ of~$(\Phi_1)_{S'}$, the sheaf $\SheafHom_{S'}(x_1,y_1)$
is constructible, then, for every pair of objects $x,~y$ of
$\Phi_{S'}$, the same holds for~$\SheafHom_{S'}(x,y)$. Suppose
that $\Phi_1$ is $1$-constructible (\Ref{XIII.0}) and that $\Phi_2$ is
constructible; then $\Phi$ is $1$-constructible.
\end{sublemma}

For every scheme $S'$ \'etale over~$S$ and every object $x$ of
$\Phi_{S'}$, one has a monomorphism
$$
\SheafAut_{S'}(x)\to\SheafAut_{S'}(p(x)).
$$
It therefore follows from SGA~4 IX~2.9 that, if $\Phi_1$ is
constructible, so is $\Phi$, and the second
assertion of the lemma is proved in the same way.

Suppose now $\Phi_1$ $1$-constructible and $\Phi_2$
constructible. The morphism $p$ induces on the sheaves of
maximal subgerbes a morphism
$$
\varphi\colon S\Phi\to S\Phi_1.
$$
Let $G$ be the image of~$S\Phi$ under $\varphi$; by SGA~4 IX~2.9,
$G$ is a constructible sheaf. One can therefore find a sheaf
representable by an \'etale $S$-scheme of finite type $T$ and
an epimorphism
$$
a\colon T\to G
$$
(SGA~4 IX~2.7).
% original p. 407
Moreover one can choose $T$ so that the image $y$ of the identity
section of~$T(T)$ under $a$ is defined by an object $x_1$ of
$(\Phi_1)_T$ of the form $x_1=p(x)$, where $x\in\Phi_T$.

It suffices to show that, for every point $s$ of~$S$, there exists a
nonempty open subset~$U$ of~$\overline{\{ s\} }$ such that $S\Phi|U$ is
locally constant constructible. Let $s\in S$, let $\overline{s}$ be a
geometric point over~$s$ and $\overline{y}_1,\dots,\overline{y}_n$ the elements of~$G_{\overline{s}}$. By
definition of~$T$ there exist morphisms
$h_i\colon\overline{s}\to T$ such that one has
$h_i^*(y)=\overline{y}_i$. Let $S'$ be the fiber product over~$S$ of
$n$ schemes isomorphic to $T$, $h\colon\overline{s}\to S'$ the
fiber product of the $h_i$, $y_i$ (\resp $x_i$) the inverse image of
$y$ (\resp $x$) under the $i$th projection of~$S'$ onto~$T$. Let
$F_i$ be the subsheaf of~$S\Phi |S'$ that is the inverse image of~$y_i$,
and let us show that the $F_i$ are constructible. The sheaf $F_i$ is
a quotient of the sheaf $F'_i$ such that, for every scheme $S''$
\'etale over~$S'$, one has
\begin{multline*}
F'_i(S'') = \{\text{classes mod.\ isomorphism of objects $z$ of
$\Phi_{S''}$ equipped}\\
\text{ with an isomorphism $i\colon p(z)\simeq p(x_i|S'')$}\}.
\end{multline*}
It suffices to show that the $F_i'$ are constructible. Now, if one
sets $z_i=p_1p(x_i)$, $z_i'=p_2p(x_i)$, one has a monomorphism
$$
\Psi_i\colon F'_i\to\SheafIsom_{S'}(z_i,z_i'),
$$
obtained by associating, with every scheme $S''$ \'etale over~$S'$ and
every object $z$ of~$\Phi_{S''}$ such that one has an isomorphism
$i\colon p(z)\simeq p(x_i|S'')$, the isomorphism from~$z_i$ to $z_i'$
defined by the condition that the diagram
$$
\xymatrix@C=1.2cm{
p_1p(z) \ar[r]^-{p_1(i)}\ar[d]_-j& z_i\ar[d]\\ p_2p(z) \ar[r]^-{p_2(i)}& z'_i
}
$$
be commutative ($j$ is the canonical morphism associated with the
exact diagram).
% original p. 408

The morphism $\Psi_i$ is injective because to say that two objects $z,~z'$
such that one has isomorphisms $i\colon p(z)\isomto
p(x_i|S''),~i'\colon p(z')\isomto p(x_i|S'')$ define the
same element of~$\SheafIsom_{S''}(z_i,z'_i)$ amounts to
saying that one has $p_1(i'^{-1}i)=p_2(i'^{-1}i)$, \ie that $i'^{-1}i$
comes from an isomorphism $z\to z'$. The sheaf $F'_i$, being a
subsheaf of~$\SheafIsom_{S'}(z_i,z'_i)$, is constructible.

Since one can find a nonempty open subset~$U$ of~$\overline{\{ s\}}$
such that $y_1|U,\dots,y_n|U$ generate $G$, it follows from
lemma~\Ref{XIII.6.1.2} below that $S\Phi |U$ is constructible,
hence, at the cost of shrinking $U$, $S\Phi |U$ is locally constant
constructible.

\begin{sublemma}
\label{XIII.3.1.2}
Let $S$ be a locally noetherian scheme, $f\colon X\to S$
a morphism, $F\to G$ a monomorphism of sheaves of groups on~$X$, $\Psi$ (\resp $\Psi_1$) the stack of torsors under $F$
(\resp under $G$), $\Phi=f_*\Psi$, $\Phi_1=f_*\Psi_1$. Suppose
given a sheaf on~$S$ representable by an \'etale $S$-scheme
of finite type $T$ and a surjective morphism
$$
a\colon T\to S\Phi_1\simeq\R^1f_*G,
$$
such that there exists a torsor $Q$ on~$X_T=X\times_ST$ with group
$G|X_T$ which defines in $\R^1f_*G(T)$ the image under $a$ of the
identity section of~$T(T)$. Let $f_T\colon X_T\to T$ be the canonical
morphism and set $F_T=F|X_T$. Suppose that $\Phi_1$ is
$1$-constructible and that $f_{T*}(Q/F_T)$ is constructible. Then
$\Phi$ is $1$-constructible.
\end{sublemma}

It suffices indeed to copy the proof
of~\Ref{XIII.3.1.1}, the fact that one has morphisms $\Psi_i$
being replaced by the fact that one has isomorphisms
$$
F'_i\isomto f_*(Q/F_T)|S'.
$$

\begin{subremark}
\label{XIII.3.1.3}
Suppose
% original p. 409
that $k=\kres(s)$ is of characteristic zero; then the
schemes of finite type over~$\overline{k}$ of dimension $\leq\dim X$
are strongly desingularizable, and the proof
of~\Ref{XIII.3.1} makes it possible to prove the following results:
\begin{enumerate}
\item[a)] There exists a nonempty open subset~$S_1$ of~$S$ such that, for every
scheme $S'$ over~$S_1$ whose maximal points are of
characteristic zero and every locally
constant constructible sheaf of sets $F$ on~$X'=X\times_SS'$, $(F,f')$ is
cohomologically proper relative to~$S'$ in dimension $\leq
0$.

\item[b)] If all the residue characteristics of~$S$ are
zero, there exists a nonempty open subset~$S_1$ of~$S$ such that, for every
scheme $S'$ over~$S_1$ and every locally constant
constructible sheaf of groups $F$ on~$X'$, $(F,f')$ is
cohomologically proper relative to~$S'$ in dimension $\leq1$.
\end{enumerate}

It suffices indeed to copy the proof of
\Ref{XIII.3.1}.2)\kern2pt a). Proposition~\Ref{XIII.2.7} used in
3)\kern2pt 4 applies to the case of a locally constant sheaf $F$, because,
taking up again the notation of 3)\kern2pt4, every torsor under $F$ is
tamely ramified on~$P$ relative to~$S$, since
all the residue characteristics of~$S$ are zero.
\end{subremark}

\begin{corollary}
\label{XIII.3.2}
Let $k$ be a field of characteristic $p\geq 0$, $p'$ the set
of prime numbers distinct from~$p$, $f\colon X\to k$ a coherent
morphism.
\begin{enumerate}
\item[1)] For every sheaf of sets $F$, $(F,f)$ is
cohomologically proper in dimension $\leq 0$.
\item[2)] Suppose that one of the following two conditions is satisfied:
\begin{enumerate}
\item[a)] $f$ is of finite type and the schemes of finite type of
dimension $\leq\dim X$ over an algebraic closure of~$k$ are
strongly de\-sin\-gu\-lar\-iz\-able.
\item[b)] The schemes of finite type over an algebraic
closure of~$k$ are
% original p. 410
strongly desingularizable.
\end{enumerate}
Then, for every sheaf of ind-$p'$-groups $F$, $(F,f)$ is
cohomologically proper in dimension $\leq 1$.
\end{enumerate}
\end{corollary}

Let $F$ be a sheaf of sets (\resp of
ind-$p'$-groups). By SGA~4 IX~2.7.2 one can write
$F$ as a filtered inductive limit
$$
F=\varinjlim F_i,
$$
where the $F_i$ are constructible sheaves of sets (\resp of
ind-$p'$-groups). Since $f$ is coherent, $f_*$
(\resp $\R^1f_*$) commutes with inductive limits (SGA~4 VII~3.3). If
one knows that the $(F_i,f)$ are cohomologically proper in
dimension $\leq 0$ (\resp that every sheaf of sets is
cohomologically proper in dimension $\leq 0$ and that the $(F_i,f)$
are cohomologically proper in dimension $\leq 1$), the same will
hold for~$(F,f)$. One may therefore assume $F$ constructible.

If one assumes $f$ of finite type (\resp satisfying a)), the
proposition follows from \Ref{XIII.3.1}~1)\kern2pt b) (\resp from
\Ref{XIII.3.1}~2)\kern2pt b)). Let us now prove~\Ref{XIII.3.2} when one no longer
assumes $f$ of finite type. For every scheme $S'$ over
$k$ and every geometric point $\overline{s}$ of~$S'$, one
denotes by $\overline{k}$ (\resp $\overline{S'}$) the strict localization of
$k$ at $\overline{s}$ (\resp of~$S'$ at $\overline{s}$),
by $\overline{X}$ the inverse image of~$X$ over~$\overline{k}$, and one
considers the cartesian square
$$
\xymatrix@C=1.2cm{
\overline{X}\ar[d]_-{\overline{f}}&\ar[l]_g \overline{X'}\ar[d]^-{\overline{f}'}\\ \overline{k}&\ar[l]\overline{S'}
}
$$
It suffices to prove that, for every $S'$ and every $\overline{s}$, the
canonical morphism
$$
\H^0(\overline{X},\overline{F})\to
\H^0(\overline{X'},\overline{F'})\qquad
\text{(\resp $\H^1(\overline{X},\overline{F})\to
\H^1(\overline{X'},\overline{F'})$)}
$$
is an isomorphism. It suffices to show that one has the relations
% original p. 411
\begin{equation*}
\label{eq:XIII.3.*}
\tag{$*$} \overline{F}\simeq g_*g^*\overline{F}\qquad
\text{(\resp $\R^1g_*(g^*\overline{F})=0$).}
\end{equation*}
Now, in this form, the question is local on~$X$ for the \'etale
topology. One may therefore assume $X$ affine; by passage to the
limit one may assume $X$ of finite type over~$k$. One then knows that
$(F,f)$ is cohomologically proper in dimension $\leq 0$
(\resp $\leq 1$) and that the same holds when $X$ is replaced by
a scheme \'etale of finite type over~$X$, which proves
\eqref{eq:XIII.3.*}.

\begin{theorem}
\label{XIII.3.3}
Let $S$ be an irreducible scheme with generic point
$s$, $f\colon X\to S$ a morphism of finite
presentation. Suppose that the schemes of finite type of dimension
$\leqslant\dim X_s$ over an algebraic closure $\bar{k}$ of
$k$ are desingularizable \textup{(EGA~IV 7.9.1)}. Then, if $\LL$
denotes the set of prime numbers distinct from the
residue characteristics of~$S$, one can find a nonempty open
subset~$S_1$ of~$S$ such that the morphism $f|S_1$ is
universally locally $1$-aspherical for~$\LL$.
\end{theorem}

One may assume $S$ integral and $X$ reduced (SGA~4
VIII~1.1). By passage to the limit one may assume $S$
noetherian. Moreover, to prove the theorem, it
suffices to do so after a finite extension $S_1\to S$, where $S_1$
is an integral scheme and where $S_1\to S$ is a composite
of \'etale extensions and of surjective radicial extensions.

Let us first show that, at the cost of restricting $S$ to a nonempty open
subset, $f$ is universally locally $0$-acyclic. At the cost of
making a radicial extension of~$\kres(s)$, one may assume that the
morphism $(X_s)_\red\to s$ is separable; one may therefore assume,
at the cost of restricting $S$ to a nonempty open subset and of making
a surjective radicial extension of~$S$, that the morphism $f$ is
flat, with separable geometric fibers (EGA~IV
12.1.1), which implies that $f$ is universally
$0$-acyclic (SGA~4 XV~4.1).

Let us show that, at the cost of restricting $S$ to a nonempty open subset,
$f$ is universally locally $1$-aspherical for
$\LL$. At the cost of making a
% original p. 412
finite extension of~$\kres(s)$, which is permissible because one can
consider it as a composite of an \'etale extension and of a
radicial extension, one can find a surjective proper morphism
$p_s\colon Y_s\to X_s$, where $Y_s$ is a scheme smooth over~$S$,
of the same dimension as $X_s$, and one may assume that $p_s$
comes from a surjective proper morphism $p\colon Y\to X$, where $Y$
is a scheme smooth over~$S$ (EGA~IV 9.6.1 and 12.1.6). It suffices to
show that, at the cost of restricting $S$ to a nonempty open subset,
for every diagram with cartesian squares
$$
\xymatrix{ S''\ar[d]_i & X ''\ar[d]^j\ar[l]_{f''} \\ S'\ar[d] &
X'\ar[d]\ar[l]_{f'} \\ S & X, \ar[l]_f }
$$
where $i$ is \'etale of finite presentation, and for every
sheaf of ind-$\LL$-groups $F$ on~$S''$, if $\Phi$ is the stack
of torsors under $F$, then the canonical morphism
$$
{f'}^*i_*\Phi \to j_*{f''}^*\Phi
$$
is an equivalence. Let $Z=Y\times_{X}Y,
T=Y\times_{X}Y\times_{X}Y$. One has in a natural way a commutative
diagram
$$
\xymatrix{ S'' \ar[d]_i & X''\ar[l]_{f''}\ar[d]_j &
Y''\ar[l]_{p''}\ar[d] & Z''\ar[d]\ar@<-.4ex>[l]\ar@<.4ex>[l] &
T''\ar[d]\ar[l]\ar@<-.8ex>[l]\ar@<.8ex>[l] \\ S' \ar[d] & X'
\ar[l]_{f'}\ar[d] & Y'\ar[d]\ar[l]_{p'}&
Z'\ar[d]\ar@<-.4ex>[l]\ar@<.4ex>[l] &
T'\ar[d]\ar[l]\ar@<-.8ex>[l]\ar@<.8ex>[l] \\ S & X\ar[l]_f & Y\ar[l]_p
& Z'\ar@<-.4ex>[l]\ar@<.4ex>[l] & T'\ar[l]\ar@<-.8ex>[l]\ar@<.8ex>[l]
}
$$
Let
% original p. 413
$q\colon Z\to X$ and $r\colon T\to X$ be the canonical morphisms,
the notations $q', r', q'', r''$ having an obvious
meaning. By \Ref{XIII.1.11}~2), one has the following
essentially commutative diagram, whose rows are exact:
$$
\xymatrix{ f'^* i_*\Phi \ar[r]\ar[d]_a &
p'_*p'^*(f'^*i_*\Phi)\ar[d]_b\ar@<-.4ex>[r]\ar@<.4ex>[r] &
q'_*q'^*(f'^*i_*\Phi)\ar[d]_c\ar[r]\ar@<-.8ex>[r]\ar@<.8ex>[r]&
r'_*r'^*(f'^*i_*\Phi)\ar[d]_d \\ j_*f''^*\Phi \ar[r]&
j_*p''_*p''^*(f''^*\Phi)\ar@<-.4ex>[r]\ar@<.4ex>[r] &
j_*q''_*q''^*(f''^*\Phi)\ar[r]\ar@<-.8ex>[r]\ar@<.8ex>[r]&
j_*r''_*r''^*(f''^*\Phi) }
$$
Since $Y$ is smooth over~$S$, the morphism $fp$ is universally
locally $1$-aspherical for $\LL$ (SGA~XV~2.1),
and it follows from \cite[VII 2.1.7]{XIII.2} that $b$ is an equivalence of
categories. On the other hand, at the cost of restricting $S$ to a
nonempty open subset, one may assume that the morphisms $Z\to S$ and $T\to
S$ are universally locally $0$-acyclic. It follows
that the functors $c$ and $d$ are fully faithful, and the
diagram above then shows that $a$ is an equivalence, which
completes the proof.
\begin{corollary}
\label{XIII.3.4}
Let $k$ be a field of characteristic $p\geqslant0$, $p'$
the set of prime numbers distinct from~$p$, $f\colon X\to k$ a
coherent morphism. Suppose that one of the following two
conditions is satisfied:
\begin{enumerate}
\item[a)] $f$ is of finite type and the schemes of finite type of
dimension $\leqslant\dim X$ over an algebraic closure of~$k$
are desingularizable.
\item[b)] The schemes of finite type over an algebraic
closure of~$k$ are desingularizable.
\end{enumerate}
Then $f$ is universally locally $1$-aspherical for $p'$.
\end{corollary}

Case a) follows from~\Ref{XIII.3.3}. In case b), the question
being local on~$X$, one may assume $X$ affine; by passage
to the limit (SGA~4 XV~1.3), one reduces to the case where $X$ is
of finite type over~$k$.
\begin{corollary}
\label{XIII.3.5}
Let
% original p. 414
$S$ be an irreducible scheme with generic point
$s$, $f\colon X\to S$ a morphism of finite
presentation. Suppose that the schemes of finite type of dimension
$\leqslant\dim X_s$ over an algebraic closure $\overline{k}$
of~$\kres(s)$ are strongly desingularizable \textup{(SGA~5 I~3.1.5)}. If
$\LL$ denotes the set of prime numbers distinct from the
residue characteristics of~$S$, one can find a nonempty open
subset~$S_1$ of~$S$ such that, for every specialization
$\bar{s}_1\to\bar{s}_2$ of geometric points of~$S_1$, the
specialization morphism \eqref{XIII.2.10}
$$
\pi_1^{\LL}(X_{\bar{s}_1})\to\pi_1^{\LL}(X_{\bar{s}_2})
$$
is bijective.
\end{corollary}

By~\Ref{XIII.3.1} and~\Ref{XIII.3.3}, one may, at the cost of
restricting $S$ to a nonempty open subset, assume that $f$ is
locally $1$-aspherical for $\LL$, and that, for every finite constant sheaf
of~$\LL$-groups $F$ on~$X$, $(F,f)$ is
cohomologically proper in dimension $\leqslant 1$. It then follows
from~\Ref{XIII.1.14} that, for every specialization
$\bar{s}_1\to\bar{s}_2$ of geometric points of~$S_1$, the
specialization morphism
$$
\left(\R^1 f_* F\right)_{\bar{s}_2} \to \left(\R^1 f_*
F\right)_{\bar{s}_1}
$$
is bijective. The corollary is nothing but the translation of the
foregoing in terms of fundamental groups.

\section{Homotopy exact sequences}\label{XIII.4}\setcounter{subsection}{-1}
\subsection{}
\label{XIII.4.0}
Let $X$ and $S$ be two connected schemes, $f\colon X\to S$ a
morphism, $a$ a geometric point of~$X$, $\LL$ a set
of prime numbers. Let $K$ be the kernel of the canonical homomorphism
$\pi_1(X,a)\to\pi_1(S,a)$ and $N$ the smallest normal
pro-subgroup of~$K$ such that $K/N$ is a pro-$\LL$-group
$K^{\LL}$. Then $N$ is normal in $\pi_1(X,a)$ and we denote by
$$
\pi_1'(X,a)
$$
the
% original p. 415
quotient of~$\pi_1(X,a)$ by $N$. If $a$ is a geometric
point of a geometric fiber~$X_{\bar{s}}$, the
canonical morphisms
$$
\pi_1(X_{\bar{s}},a)\to \pi_1(X,a) \to
\pi_1(S,a)
$$
allow one to define canonical morphisms
$$
\pi_1^{\LL}(X_{\bar{s}},a)\lto{u} \pi_1'(X,a) \lto{v} \pi_1(S,a)
$$
We have $vu=0$.

\begin{proposition}
\label{XIII.4.1}
Let $S$ be a connected scheme, $f\colon X\to S$ a
locally $0$-acyclic morphism \textup{(SGA~4 XV~1.11)}; suppose moreover
that $f$ is $0$-acyclic (which, when $f$ is coherent, amounts to
saying that the geometric fibers of~$f$ are connected \textup{(SGA~4
XV~1.16)}). Let $\LL$ be a set of prime numbers. If $S'$ is a
scheme \'etale over~$S$, we denote by $X'$, $f'$ the inverse images
of~$X$, $f$ over~$S'$. Suppose that, for every \'etale covering
$S'$ of~$S$ and for every \'etale covering $E$ of
$X'$ that is a quotient of a Galois covering whose group is an
$\LL$-group, $(E,f')$ is cohomologically proper relative
to~$S'$ in dimension $\leqslant 0$ and $f'_* E$ is
constructible. Then, if $\bar{s}$ is a geometric point of
$S$ and $a$ a geometric point of the fiber $X_{\bar{s}}$, the
sequence of group homomorphisms
\begin{equation*}
\label{eq:XIII.4.1.1}
\tag{\thesubsection.1}
\pi_1^{\LL}(X_{\bar{s}},a)\lto{u} \pi_1'(X,a)\lto{v}\pi_1(S,a)\to1
\end{equation*}
is exact
\end{proposition}

This statement generalizes X.\Ref{X.1.4}, whose proof we shall
copy.

Let us first show that $v$ is surjective. It suffices to show that, for
every connected \'etale covering $S'$ of~$S$, $X'$ is also
connected (V~\Ref{V.6.9}). Let $C$ be a set having at least two
elements. It follows from the fact that $f$ is $0$-acyclic that
the canonical morphism
$$
\H^0(S',C_{S'})\to \H^0(X',C_{X'})
$$
is bijective, hence that $X'$ is connected, whence the surjectivity
of~$v$.

By
% original p. 416
definition of~$K^{\LL}$ \eqref{XIII.4.0}, we have the exact sequence
$$
1\to K^{\LL} \to \pi_1'(X,a)\to \pi_1(S,a)\to 1
$$
Let $\widetilde{S}$ be the universal covering of~$S$ and
$\widetilde{X}=\widetilde{S}\times_SX$; the group $K^{\LL}$ classifies the
Galois coverings $P$ whose group is an $\LL$-group, such that there
exist an \'etale covering $S'$ of~$S$ and a Galois covering
$Q$ of~$X'=X\times_SS'$ such that one has an isomorphism
$P\simeq Q\times_{X'}\widetilde{X}$. For the sequence
\eqref{eq:XIII.4.1.1} to be exact, it is necessary and sufficient that the
canonical morphism
$$
\pi_1^{\LL}(X_{\bar{s}},a)\to K^{\LL}
$$
be surjective. By the interpretation of~$K^{\LL}$ this
amounts to saying that, for every \'etale covering $S'$ of~$S$
and for every Galois covering $Q$ of~$X'$ whose group is an
$\LL$-group, such that $P=Q\times_{X'}\widetilde{X}$ is connected,
$Q|X_{\bar{s}}$ is then connected. Let us show that this last
condition is satisfied. Indeed, let $S'$ be an \'etale
covering of~$S$, $Q$ a Galois covering of~$X'$ whose group is an
$\LL$-group $F$, such that $Q|X_{\bar{s}}$ is disconnected, and let us show
that, after replacing $S'$ by an \'etale covering, $Q$
becomes disconnected. There exist a subgroup $G$ of~$F$ distinct from
$F$ and a torsor $R$ under $G|X_{\bar{s}}$ such that $Q|X_{\bar{s}}$
is obtained by extension of the structure group $G\to F$ from
$R$. The \'etale covering $E=Q/G$ of~$X'$ is such that
$E|X_{\bar{s}}$ has a section. By~\Ref{XIII.1.16} $f_*E$
is locally constant constructible and, after replacing $S'$
by an \'etale covering, one may even assume that $f_*E$
is constant. Since $(E,f')$ is cohomologically proper relative
to~$S'$ in dimension $\leqslant0$ and since $\H^0(X_{\bar{s}},
E|X_{\bar{s}})$ is nonempty, we see that $E$ has a section. But this
proves that $Q$ is disconnected, which completes the proof.

From~\Ref{XIII.4.1} we deduce the following lemma, which will be used
in~\Ref{XIII.4.6}.

\begin{lemma}
\label{XIII.4.2}
Let $S$ be a connected scheme, $f\colon X\to S$ a
locally $0$-acyclic and $0$-acyclic morphism, $\LL$ a set of
prime numbers. Suppose
% original p. 417
that, for every finite constant sheaf of~$\LL$-groups $F$ on~$X$,
$(F,f)$ is cohomologically proper relative to~$S$ in
dimension $\leqslant1$, and that, for every \'etale covering
$S'$ of~$S$ and for every \'etale covering $E$ of~$X'$
that is a quotient of a Galois covering whose group is an $\LL$-group,
$f'_*E$ is constructible. Then, if $\bar{s}$ is a
geometric point of~$S$ and $a$ a geometric point of the
fiber $X_{\bar{s}}$, the sequence of group homomorphisms
$$
\pi_1^{\LL}(X_{\bar{s}},a)\to\pi_1'(X,a)\to\pi_1(S,a)\to1
$$
is exact.
\end{lemma}

the hypotheses of~\Ref{XIII.4.1} are satisfied. It follows
indeed from~\Ref{XIII.1.13} 3) that, for every scheme $S'$
\'etale over~$S$ and for every \'etale covering $E$ of~$X'$
that is a quotient of a Galois covering whose group is an $\LL$-group,
$(E,f')$ is cohomologically proper relative to~$S'$ in
dimension $\leqslant 0$.
\begin{proposition}
\label{XIII.4.3}
Let $S$ be a connected scheme, $\LL$ a set of prime
numbers, $f\colon X\to S$ a $0$-acyclic morphism, locally
$1$-aspherical for $\LL$ \textup{(SGA~4 XV~1.11)}, $g\colon S\to X$ a
section of~$f$. Let $\bar{s}$ be a geometric point of~$S$,
$a$ a geometric point of the fiber $X_{\bar{s}}$. Suppose
that, for every constant sheaf of
$\LL$-groups~$F$, $(F,f)$ is cohomologically proper in dimension
$\leqslant1$, that the direct image under~$f$ of the stack of torsors under
$F$ is a $1$-constructible stack \eqref{XIII.0} and that, for every
\'etale covering $E$ of~$X'$ that is a quotient of a Galois
covering whose group is an $\LL$-group, $f'_*E$ is constructible. Then
the sequence of group homomorphisms
\begin{equation*}
\label{eq:XIII.4.3.1}
\tag{\thesubsection.1}
1\to\pi_1^{\LL}(X_{\bar{s}},a)\lto{u} \pi_1'(X,a)\lto{v}\pi_1(S,a)\to1
\end{equation*}
is exact.
\end{proposition}

Taking~\Ref{XIII.4.2} into account, it suffices to show the injectivity
of the morphism $u$, \ie to prove that, for every principal
covering $\overline{Z}$ of~$X_{\bar{s}}$
% original p. 418
whose group is an $\LL$-group~$C$, there exist an \'etale covering
$Z$ of~$X$ and a morphism from a connected component of~$Z|X_{\bar{s}}$
to~$\overline{Z}$ (V~\Ref{V.6.8}). So let $\overline{Z}$ be a
principal covering of~$X_{\bar{s}}$ whose group is an $\LL$-group~$C$, and~$\bar{z}$ its class in
$\H^1(X_{\bar{s}},C_{X_{\bar{s}}})$. By \Ref{XIII.1.5}~d), we
have a canonical isomorphism
$$
\left(\R^1f_*C_X\right)_{\bar{s}}\isomto\H^1(X_{\bar{s}},C_{X_{\bar{s}}}),
$$
and by~\Ref{XIII.1.16}, $\R^1f_*C_X$ is a locally constant
constructible sheaf. One can therefore find an
\'etale covering $S'$ of~$S$ such that $\R^1f_*C_X|S'$ is
constant. If $\bar{s}\to S'$ is a geometric point
above the geometric point $\bar{s}\to S$, there exists an
element~$z$ of~$\H^0(S', \R^1f_*C_X)$ whose image in
$\H^1(X_{\bar{s}},C_{X_{\bar{s}}})$ is $\bar{z}$. By
Lemma~\Ref{XIII.4.3.1} below, one can find an \'etale
covering $S'_1$ of~$S'$ and a torsor $P$ on~$X'_1$ with group~$C$
whose image in $\H^0(S'_1,\R^1f_*C)$ is equal to the
restriction of~$z$. The torsor~$P$ is representable by an
\'etale covering $Z$ of~$X'_1=X\times_SS'_1$ such that
$Z\times_{X'_1}X_{\bar{s}}$ is isomorphic to $\overline{Z}$. If
one regards $Z$ as an \'etale covering of~$X$, one
then has a morphism from~$Z\times_XX_{\bar{s}}$ to $\overline{Z}$, which
completes the proof.

\begin{sublemma}
\label{XIII.4.3.1}
Let $f\colon X\to S$ be a $0$-acyclic and locally
$0$-acyclic morphism, $g$ a section of~$f$. Let $C$ be a finite
constant group such that $(C_X,f)$ is cohomologically proper in dimension
$\leqslant0$ and such that the direct image under $f$ of the stack of torsors under
$C_X$ is constructible. Then, for every section $z$ of
$\H^0(S,\R^1f_*C_X)$, one can find an \'etale covering
$S_1$ of~$S$ and, if $X_1=X\times_SS_1$, an element of
$\H^1(X_1,C_{X_1})$ whose image under the canonical morphism
$$
\H^1(X_1,C_{X_1})\to\H^0(S_1,\R^1f_*C_{X_1})
$$
is equal to the restriction of~$z$ to
$\H^0(S_1,\R^1f_*C_{X_1})$.
\end{sublemma}

For every scheme $S'$ \'etale over~$S$, we set $X'=X\times_SS'$,
and we denote by $g'$ (\resp $F'$, etc.) the inverse image of~$g$ (\resp $F$,
etc.) under the morphism
% original p. 419
$S'\to S$. The presheaf $G$ on~$S$ defined by
\newlength{\tmpRacinet} \setlength{\tmpRacinet}{\textwidth}
\addtolength{\tmpRacinet}{-7em}
$$
G(S') = \left\{\text{\parbox{\tmpRacinet}{classes mod.\ isomorphism of
		torsors $P$ on~$X'$ with group $C_{X'}$, endowed with an
		isomorphism $g'^*P\isomto C_{S'}$ }} \right\}
$$
is then a sheaf. Indeed, this follows by descent from the fact
that an isomorphism of a torsor $P$ on~$X'$ is completely
determined by its restriction to $g'(S')$. Moreover, we have a
surjective morphism
$$
G\to \R^1f_*F.
$$
Let $z$ be an element of~$\H^0(S,\R^1f_*C_X)$ and $H$ the
subsheaf of~$G$ that is the inverse image of~$z$. It suffices to show
that $H$ is a locally constant constructible sheaf. Now, this
property being local on~$S$, we may assume that $z$
comes from an element of~$\H^1(X,C_X)$ represented by
a torsor $P$ such that $g^*P$ is isomorphic to $C_S$. Giving an
isomorphism $i\colon g^*P\isomto C_S$ amounts to giving a
global section of~$\SheafAut_{C_S}(g^*P)$, and two isomorphisms $i$
and $i'$ define the same element of~$G(X)$ if and
only if $ii'^{-1}$ is the image of an element of
$\Aut_{C_X}(P)$. If one considers the canonical injection
$$
f_*\SheafAut_{C_X}(P)\to\SheafAut_{C_S}(g^*P)\simeq C,
$$
$H$ is therefore identified with the quotient of~$\SheafAut_{C_S}(g^*P)$ by
$f_*\SheafAut_{C_X}(P)$. By~\Ref{XIII.1.16}
$f_*\SheafAut_{C_X}(P)$ is locally constant; hence the same
holds for~$H$, which completes the proof.

\begin{examples}
\label{XIII.4.4}
Let us note that, if $S$ is a connected scheme, the hypotheses
of~\Ref{XIII.4.1} are satisfied when the morphism $f$ is proper,
flat, of finite presentation, with separable connected geometric
fibers, $\LL$ being arbitrary
(\cf X~\Ref{X.1.3}). The hypotheses of~\Ref{XIII.4.3} are
satisfied if moreover $f$ is smooth and has a section, if one
denotes by $\LL$ the set of prime numbers distinct from the
residue characteristics of~$S$ (SGA~4 XV~2.1 and XVI~5.2).

The hypotheses of~\Ref{XIII.4.1} are also satisfied if $S$ is connected,
if one has a scheme~$Z$ proper of finite presentation, flat over~$S$, with
% original p. 420
separable connected geometric fibers, such that $X$ is
the complement in $Z$ of a divisor with normal crossings
relative to~$S$, $\LL$ being the set of prime
numbers distinct from the residue characteristics of~$S$
\eqref{XIII.2.9}. The hypotheses of~\Ref{XIII.4.3} are
satisfied if moreover $f$ is smooth and has a section.
\end{examples}

\subsection{}
\label{XIII.4.5}
Let us take up again the notation and hypotheses of~\Ref{XIII.4.3}. If
$\bar{s}$ is a geometric point of~$S$ and $a=g(\bar{s})$,
the section $g$ allows one to define a morphism
$$
w\colon \pi_1(S,a)\to \pi_1'(X,a),
$$
so that $\pi_1'(X,a)$ is identified with the semi-direct product of
$\pi_1(S,a)$ by $\pi_1(X_{\bar{s}},a)$. The profinite group
$\pi_1(S,a)$ therefore acts on~$\pi_1(X_{\bar{s}}, a)$. Since
$\pi_1^\LL(X_{\bar{s}},a)$ is a strict projective limit of groups
invariant under the action of~$\pi_1(S,a)$, the datum of
$\pi_1^\LL(X_{\bar{s}},a)$ endowed with this action is equivalent
to the datum of a strict projective system of finite \'etale
group schemes over~$S$, which we denote by
$$
\pi_1^\LL(X/S,g,\bar{s})\quad\text{or
simply}\quad\pi_1^\LL(X/S,g).
$$

We then have the following properties:
\setcounter{subsubsection}{0}

\subsubsection{}
\label{XIII.4.5.1}
For every finite \'etale group scheme $G$ over~$S$ whose
fibers are $\LL$-groups, the set $E$ of classes of torsors
$P$ under the inverse image $G_X$ of~$G$ on~$X$, endowed with an isomorphism
$g^*P\isomto G$, is canonically isomorphic to the set
$$
\Hom_S(\pi_1^\LL(X/S,g,\bar{s}),G) \quad\text{mod.\ inner
automorphisms of~$G$.}
$$

\subsubsection{}
\label{XIII.4.5.2}
For every finite \'etale group scheme $G$ over~$S$ whose
fibers are $\LL$-groups, the sheaf $\R^1f_*G_X$ is
canonically isomorphic to the sheaf associated with the presheaf
$$
S'\mto \Hom_{S'}(\pi_1^\LL(X/S,g,\bar{s}),G)\
\text{mod.\ inner automorphisms of~$G$.}
$$
($S'$ denotes a scheme \'etale over~$S$).

\subsubsection{}
\label{XIII.4.5.3}
Let
% original p. 421
$S'$ be a connected
$S$-scheme, $\bar{s}$ a geometric point of~$S'$, $X'$, $g'$
the respective inverse images of~$X,g$ over~$S'$. Then
$\pi_1^\LL(X'/S',g',\bar{s})$ is canonically isomorphic to
the inverse image of~$\pi_1^\LL(X/S,g,\bar{s})$ over~$S'$. For every
geometric point $\xi$ of~$S$, the fiber
$\pi_1^\LL(X/S,g,\bar{s})_\xi$ is isomorphic to $\pi_1^\LL(X_\xi)$.

Indeed, the datum of~$G$ is equivalent to the datum
of an abstract $\LL$-group $\mathbf{G}$
on which $\pi_1(S,a)$ acts, whence an action of
$\pi_1'(X,a)$ on~$\mathbf{G}$. The isomorphism defined
in~\Ref{XIII.4.5.1} is then obtained by restriction to the subset
$E$ from the canonical morphism
\[
\H^1(\pi_1'(X,a),\mathbf{G})\to
\H^1(\pi_1^\LL(X_{\bar{s}},a),\mathbf{G}) =
\Hom(\pi_1^\LL(X_{\bar{s}},a),\mathbf{G})/\text{inn.\ aut.\ }G,
\]
the set $E$ being mapped bijectively onto the subset of
morphisms from $\pi_1^\LL(X_{\bar{s}},a)$ to $\mathbf{G}$ that are
compatible with the action of
$\pi_1(S,a)$. Assertion~\Ref{XIII.4.5.3} then follows from the
definition of~$\pi_1^\LL(X/S,g,\bar{s})$ taking into account the homotopy
exact sequence \eqref{eq:XIII.4.3.1}, and~\Ref{XIII.4.5.2} is
deduced from \Ref{XIII.4.5.1} and~\Ref{XIII.4.5.3}.

\begin{proposition}[K\"unneth formula]
\label{XIII.4.6}
Let $k$ be a separably closed field of characteristic
$p\geq 0$, $X$ and $Y$ two connected $k$-schemes, $a$ a
geometric point of~$X$, $b$ a geometric point of~$Y$,
$c$ a geometric point of~$X\times_k Y$ above~$a$ and
$b$. Suppose that one of the following two conditions is satisfied:
\begin{enumerate}
\item[a)] $X$ is of finite type over~$k$ and the schemes of finite type
over an algebraic closure $\bar{k}$ of~$k$, of dimension
$\leq\dim X$, are strongly desingularizable. \textup{(SGA~5 I~3.1.5)}.
\item[b)] $X$ is quasi-compact and quasi-separated and every
scheme of finite type over~$\bar{k}$ is strongly
desingularizable.
\end{enumerate}

Then, if $p'$ is the set of prime numbers distinct from~$p$,
the morphism
\begin{equation*}
\label{eq:XIII.4.6.0} \tag{\thesubsection.0} {}
\pi_1^{p'}(X\times_k Y,c)\to\pi_1^{p'}(X,a)\times\pi_1^{p'}(Y,b),
\end{equation*}
deduced
% original p. 422
from the homomorphisms on the fundamental groups associated with the
projections
$$
X\times_k Y\to X \quad \text{and} \quad X\times_k Y\to Y,
$$
is an isomorphism.
\end{proposition}

We may assume $k$ algebraically closed and $X$ reduced
(SGA~4 VIII~1.1). Let $Z=X\times_k Y$, and let $g\colon X\to k$ and $f\colon
Z\to Y$ be the canonical morphisms. The morphism $g$, hence also $f$,
is universally locally $1$-aspherical for $p'$
by~\Ref{XIII.3.5}. Since $X$ is connected, $f$ is
$0$-acyclic (SGA~4 XV~1.16). On the other hand it follows
from~\Ref{XIII.3.2} that, for every finite $p'$-group $C$, $(C_X,g)$ is
cohomologically proper relative to $k$ in dimension $\leq
1$. It follows that $(C_Z,f)$ is cohomologically proper
relative to $Y$ in dimension $\leq 1$ (\Ref{XIII.1.5}~c)) and
that $f_*C_Z$ and $\R^1f_*C_Z$ are constant sheaves. Consequently
$g$ satisfies all the hypotheses of~\Ref{XIII.4.2}. We
therefore have the exact sequence
$$
\pi_1^{p'}(X_b,c)\to\pi_1^{p'}(Z,c)\to\pi_1^{p'}(Y,b)\to 1.
$$
Moreover the composite morphism
$$
\pi_1^{p'}(X_b,c)\to\pi_1^{p'}(Z,c)\to\pi_1^{p'}(Y,b)
$$
is an isomorphism and we therefore have the exact sequence
$$
1\to\pi_1^{p'}(X,a)\to\pi_1^{p'}(Z,c)\to\pi_1^{p'}(Y,b)\to 1.
$$
On the other hand the morphism \eqref{eq:XIII.4.6.0} allows one to define
a morphism from this exact sequence to the exact sequence
$$
1\to\pi_1^{p'}(X,a)\to
\pi_1^{p'}(X,a)\times\pi_1^{p'}(Y,b)\times\pi_1^{p'}(Y,b)\to 1,
$$
and it follows that the morphism \eqref{eq:XIII.4.6.0} is an
isomorphism.

\subsection{}
\label{XIII.4.7} Let
$$
\xymatrix@C=1.2cm{
X \ar[d]_-f&\ar[l] X_U \ar[d]_-{f_U}& \ar[l]X_{\bar{s}}\ar[d] \\
S&\ar[l] U &\ar[l]\bar{s}
}
$$
be a
% original p. 423
diagram whose squares are cartesian, where $S$ is an
arcwise connected scheme (SGA~4 IX~2.12), $U$ a connected open subset of
$S$, $\bar{s}$ a geometric point of~$U$. Let $a$ be a
geometric point of~$X_{\bar{s}}$, $\LL$ a set of
prime numbers. Let $g$ be a section of~$f$ and suppose that the
following conditions are satisfied:
\begin{enumerate}
\item[a)] The morphism $f$ is $0$-acyclic, locally $0$-acyclic,
and, for every \'etale covering $S'$ of~$S$ and every
\'etale covering $E$ of~$X\times_S S'$ that is a quotient of a
Galois covering whose group is an $\LL$-group, $(E,f_{(S')})$ is
cohomologically proper relative to~$S'$ in dimension $\leq
0$.
\item[b)] The morphism $f_U$ is locally $1$-aspherical for
$\LL$, and, for every finite constant sheaf of~$\LL$-groups $F$ on~$X_U$, $(F,f_U)$ is cohomologically proper in dimension $\leq 1$ and
the fibers of~$\R^1 f_{U*}F$ are finite.
\end{enumerate}

From~\Ref{XIII.4.1} and~\Ref{XIII.4.3} one then deduces the following
commutative diagram, whose rows are exact:
\begin{equation*}
\label{eq:XIII.4.7.0} \tag{\thesubsection.0} {}
\begin{array}{c}
\xymatrix{
1 \ar[r]& \pi_1^{\LL}(X_{\bar{s}},a) \ar[r]\ar@{=}[d]& \pi'_1(X_U,a) \ar[r]\ar[d]& \pi_1(U,a)\ar[r]\ar[d]& 1 \\
&\pi_1^{\LL}(X_{\bar{s}},a) \ar[r]&\pi'_1(X,a) \ar[r]& \pi_1(S,a) \ar[r]& 1
}
\end{array}
\end{equation*}

Thanks to the section $g$, we have morphisms
$$\pi_1(U,a)\to\pi'_1(X_U,a),\quad \pi_1(S,a)\to\pi'_1(X,a)\,\;$$
from which one
deduces a morphism from the amalgamated sum of~$\pi_1(S,a)$ and
$\pi'_1(X_U,a)$ over~$\pi_1(U,a)$ to $\pi'_1(X,a)$:
\begin{equation*}
\label{eq:XIII.4.7.1} \tag{\thesubsection.1} {}
\varphi\colon\pi=\pi_1(S,a)\coprod_{\pi_1(U,a)}\pi'_1(X_U,a)\to\pi'_1(X,a).
\end{equation*}
Suppose that the following condition is satisfied:
\begin{enumerate}
\item[c)] If $T=S-U$, we have $\prof \et_T(S)\geq 2$ (SGA~2 XIV~1.1).
\end{enumerate}
Then the functor which associates with an \'etale covering of~$S$
its restriction to $U$ is fully faithful
(SGA~2 XVI~1.4). It follows that the morphism
$\pi_1(U,a)\to\pi_1(S,a)$ is surjective (V~\Ref{V.6.9}) and one
deduces from the
% original p. 424
diagram \eqref{eq:XIII.4.7.0} that the same holds for the morphism
$\pi'_1(X_U,a)\to\pi'_1(X,a)$; a fortiori $\varphi$ is an
epimorphism. Let
\begin{equation*}
\label{eq:XIII.4.7.2} \tag{\thesubsection.2} {}
K=\Ker(\pi_1(U,a)\to\pi_1(S,a)).
\end{equation*}
The group $\pi$ of \eqref{eq:XIII.4.7.1} is identified with the quotient of
$\pi'_1(X_U,a)$ by the closed invariant subgroup generated
by the image $L$ of~$K$ in $\pi'_1(X_U,a)$. Let us regard
$\pi'_1(X_U,a)$ as the semi-direct product of~$\pi_1(U,a)$ by
$\pi_1^\LL (X_{\bar{s}},a)$. The group $K$ then acts by
inner automorphisms on~$\pi_1^\LL (X_{\bar{s}},a)$, and the
quotient $\pi=\pi'_1(X_U,a)/L$ is identified with the semi-direct product of
$\pi_1(U,a)/K=\pi_1(S,a)$ by the group $\pi_1^\LL(X_{\bar{s}},a)_K$
of coinvariants of~$\pi_1^{\LL}(X_{\bar{s}},a)$ under~$K$. We
finally have an epimorphism
\begin{equation*}
\label{eq:XIII.4.7.3} \tag{\thesubsection.3} {}
\varphi\colon\pi=\pi_1^\LL(X_{\bar{s}},a)_K\cdot\pi_1(S,a)\to\pi'_1(X,a).
\end{equation*}

The proposition which follows gives conditions under which the
morphism $\varphi$ is an isomorphism.

\setcounter{subsubsection}{3}

\begin{subproposition}
\label{XIII.4.7.4} The notation is that of~\Ref{XIII.4.7}. We
suppose that, in addition to conditions \textup{a), b), c)}, the following
conditions are satisfied:
\begin{enumerate}
\item[d)] For every point $t$ of~$T=S-U$, the morphism $f$ is
locally $1$-aspherical for $\LL$ at $g(t)$.
\item[e)] For every point $t\in T$, every irreducible component
of the fiber $X_t$ contains $g(t)$ and, for every point $x$ of
$X_t-\{g(t)\}$ which is not maximal, we have
$$
\prof\hop_x(X)\geq 3 \quad\textup{(SGA~2 XIV~1.2)}
$$
and the ring $\cal{O}_{X,x}$ is noetherian.
\end{enumerate}

Then the morphism \eqref{eq:XIII.4.7.3} is an isomorphism.
\end{subproposition}

As was said previously, the group $\pi$ is identified with the
quotient of~$\pi'_1(X_U,a)$ by the closed invariant subgroup $L$
generated by the image of~$K$ \eqref{eq:XIII.4.7.2} in
$\pi'_1(X_U,a)$. This amounts to saying that $\pi$ classifies the
% original p. 425
principal coverings $Z$ of~$X_U$ such that $g^{-1}_U(Z)$
extends to an \'etale covering of~$S$, and which induce on~$X_{\bar{s}}$ a covering obtained by extension of the structure
group from a principal covering whose group is an
$\LL$-group. To prove that $\varphi$ is an isomorphism, it
suffices to show that such a covering $Z$ extends to all
of~$X$.

Let us first show that $Z$ extends to an \'etale covering
of an open set containing $X_U$ and $g(S)$. Let $W$ be a scheme
\'etale over~$X$ whose image contains $X_U$ and $g(S)$, and let us set
$W_U=W\times_S U$. From the fact that the morphism $W\to S$ is $0$-acyclic
and from the relation $\mathrm{prof}.\,\et_T S\geq 2$, it follows that we have
$$
\mathrm{prof}.\,\et_{(W-W_U)} W\geq 2
$$
(SGA~2 XIV~1.13); consequently, if $Z|W_U$ extends to an
\'etale covering of~$W$, this extension is unique up to
unique isomorphism. It follows that the problem of
extending $Z$ to a neighborhood of~$g(S)\cup X_U$ is local for the
\'etale topology in the neighborhood of the points of~$g(T)$. If $t$ is a
point of~$T$, we set $x=g(t)$, and we denote by $\bar{X}$ (\resp $\bar{S}$)
the strict localization of~$X$ at $\bar{x}$ (\resp of~$S$ at
$\bar{t}$), $\bar{U}=U\times_S \bar{S}$, $\bar{X}_U=\bar{X}\times_X
X_U$, $\bar{g}\colon\bar{S}\to\bar{X}$ the morphism deduced from
$g$. It suffices to show that, for every point $t$ of~$T$, the inverse
image $\bar{Z}$ of~$Z$ on~$\bar{X}_U$ extends to $\bar{X}$
or, what amounts to the same, is trivial. Now, by definition
of~$Z$, the inverse image of~$\bar{Z}$ on~$\bar{U}$ is
trivial. To prove that $\bar{Z}$ is trivial, it suffices to show
that it is of the form $\bar{f}^*_U E$, where $E$ is a
principal covering of~$\bar{U}$; indeed we shall then have
$E\simeq\bar{g}^*_U\bar{f}^*_U E \simeq\bar{g}^*_U\bar{Z}$, whence
the result since $\bar{g}^*_U\bar{Z}$ is trivial. Now, the
morphism $\bar{f}_U$ being $0$-acyclic and locally $0$-acyclic,
it suffices, in order to prove that $\bar{Z}$ comes from~$\bar{U}$, to
show that, for every geometric point algebraic over a
point of~$\bar{U}$, which one may assume to be the point
$\bar{s}$, $\bar{Z}|X_{\bar{s}}$ is trivial (SGA~4 XV~1.15). But since
$\bar{Z}|X_{\bar{s}}$ is obtained by extension of the structure
group from a principal covering whose group is an
$\LL$-group, this follows from the fact that the morphism $\bar{f}$ is
$1$-aspherical for $\LL$.

We
% original p. 426
have therefore shown that there exists an open neighborhood $V$ of
$g(S)\cup X_U$ such that $Z$ extends to an \'etale covering
$Z_V$ of~$V$. Let us show that $Z_V$ extends to all of~$X$.
It suffices to see that, for every point $x$ of~$X-V$, we have
$$
\prof\hop_x X\geq 3.
$$
Now this follows from hypothesis e) and from the fact that a point $x$
of~$X-V$ cannot be maximal in its fiber $X_t$ since, every
irreducible component of~$X_t$ containing $g(t)$, every maximal
point of~$X_t$ belongs to $V$.

\begin{corollary} \label{XIII.4.8}
The hypotheses are those of~\Ref{XIII.4.7.4}, but we assume moreover
that $\pi_1(S,a)=1$. Then we have an isomorphism
\begin{equation*}
\label{eq:XIII.4.8.*} \tag{$*$}
{}\pi_1^\LL(X,a)\isomto\pi_1^\LL(X_{\bar{s}},a)_K.
\end{equation*}
In particular, if $\pi_1^\LL(X_{\bar{s}},a)$ is topologically of
finite presentation and if $K$ acts on~$\pi_1^\LL(X_{\bar{s}},a)$ through a group of finite
type, then $\pi_1^\LL(X,a)$ is topologically of finite
presentation.
\end{corollary}

The isomorphism \eqref{eq:XIII.4.8.*} was proved
in~\Ref{XIII.4.7.4}. Suppose that $\pi_1^\LL(X_{\bar{s}},a)$ is the quotient
of the free pro-$\LL$-group on $n$ generators
$L(x_1,\dots,x_n)$ by the closed invariant subgroup generated
by the elements $y_1,\dots,y_p$ of~$L(x_1,\dots,x_n)$, and
suppose that $K$ acts through a group
generated by elements $k_1,\dots,k_q$. If for every
$i\in[1,n]$, $j\in[1,q]$, we denote by $z_{ij}$ an element of
$L(x_1,\dots,x_n)$ lifting the element $(k_j\cdot
x_i)x_i^{-1}$, then $\pi_1^\LL(X_{\bar{s}},a)_K$ is the quotient of
$L(x_1,\dots,x_n)$ by the closed invariant subgroup generated
by the elements $(y_i)_{i\in[1,p]}$,
$(z_{ij})_{i\in[1,n],j\in[1,q]}$.

\begin{remarks}
\label{rem:XIII.4.6}
a) Conditions a) to e) of~\Ref{XIII.4.7} are
satisfied when $S$ is a connected normal scheme, $U$ a
dense retrocompact open subset of~$S$, $f$ a proper morphism of
finite presentation, with fibers that are geometrically connected
and irreducible at every point $t$ of~$T$, $f$ being moreover
separable,
% original p. 427
smooth at the points of~$X_U\cup g(T)$, $\LL$ being the set of
prime numbers distinct from the residue characteristics of
$S$, and $X$ being regular at every point of~$X_t$. Condition
a) follows indeed from SGA~4 XV~4.1 and 1.4; conditions
b) and d) follow from~\Ref{XIII.1.4} and SGA~4 XV~2.1 and
XVI~5.2. Finally e) follows from SGA~XIV~1.11.
\par\smallskip
b) Corollary~\Ref{XIII.4.5} applies to compute the
fundamental group $\pi_1^{p'}(X)$ of a proper and smooth surface $X$
over a separably closed field $k$ of characteristic $p$
($p'$~denoting the set of prime numbers distinct from
$p$). The method was communicated to us by
J.P\ptbl \textsc{Murre}; it consists in reducing, by blowing up
$X$, to the case where one has a fibration $X\to\PP^1_k$ and an open set $U$
of~$\PP^1_k$ satisfying the hypotheses of~\Ref{XIII.4.7} (see
SGA~7 for more details). The same method can be
used more generally (\loccit) to prove that, if
$X$ is a connected $k$-scheme of finite type, and if the schemes
of finite type of dimension $\leq\dim X$ over an algebraic
closure of~$k$ are strongly desingularizable
(SGA~5 I~3.1.5), then $\pi_1^{p'}(X)$ is topologically of
finite presentation.
\end{remarks}

\section{Appendix I: Variations on Abhyankar's lemma}
\label{XIII.5}
%
This appendix contains various variants of Abhyankar's lemma.

\begin{proposition}
\label{XIII.5.1} Let $X=\Spec A$ be a regular local scheme,
$D=\sum_{1\leq i\leq r} \divisor f_i$ a divisor with normal
crossings, where the $f_i$ are elements of the
maximal ideal of~$A$ forming part of a regular system of
parameters. Let $n_i$, $1\leq i\leq r$, be integers $\geq 0$
and set
$$
X'=X[T_1,\dots,T_r]/(T_1^{n_1}-f_1,\dots,T_r^{n_r}-f_r),
$$
$U'=U\times_X X'$. Then $X'$ is regular and $U'$ is the
complement in $X'$ of the divisor with normal crossings
$\sum_{1\leq i\leq r}\divisor T_i$. If the integers $n_i$ are prime
% original p. 428
to the residue characteristic $p$ of~$X$, $U'$ is a
connected \'etale covering of~$U$, tamely
ramified relative to $D$ \textup{(\Ref{XIII.2.3}~c))}.
\end{proposition}

Indeed $X'$ is the spectrum of a local ring $A'$ whose
maximal ideal is generated by $T_1,\dots,T_r$.
(We may assume $r=\dim(A)$.)
Since $A'$ is finite and flat over~$A$, hence of dimension $r$, $A'$ is
regular (EGA~$0_{\textup{IV}}$~17.1.1) and the $T_i$ form a regular
system of parameters of~$A'$. Suppose the $n_i$ prime
to $p$. Since all the $f_i$ are invertible on~$U$, the fact that
$U'$ is \'etale over~$U$ follows from I~\Ref{I.7.4}. Moreover $U'$
is tamely ramified relative to $D$; indeed, let
$x_i$ be the generic point of~$V(f_i)$, $\bar{R}$ the
strict localization of~$R=\cal{O}_{X,x_i}$, $\bar{K}$ the field of
fractions of~$\bar{R}$. Then the $\bar{K}$-algebra that
represents $U'|\bar{K}$ is obtained from the field
$\bar{K}[T_i]/(T_i^{n_i}-f_i)$ by making an unramified
extension; it is therefore tamely ramified over~$\bar{R}$.

\begin{proposition}[Absolute Abhyankar lemma]
\label{XIII.5.2}
%
Let $X$ be a regular local scheme,
$$
D=\sum_{1\leq i\leq r} \divisor f_i
$$
a divisor with normal crossings as in~\Ref{XIII.5.1},
$Y=\Supp D$, $U=X-Y$. Let~$V$ be an \'etale covering of~$U$
tamely ramified relative to $D$. If $x_i$ is the
generic point of the closed set $V(f_i)$, $\cal{O}_{X,x_i}$ is
a discrete valuation ring with field of fractions $K_i$, and
we have $V|K_i=\Spec(\prod_{j\in J_i}L_j)$, where the $L_j$ are
finite separable extensions of~$K_i$; denote by $n_j$ the order of the
inertia group of a Galois extension generated by $L_j$
and by $n_i$ the l.c.m.\ of the $n_j$ as $j$ runs through $J_i$. If we set
$$
X'=X[T_1,\dots,T_r]/(T_1^{n_1}-f_1,\dots,T_r^{n_r}-f_r),
$$
$U'=U_{(X')}$, $V'=V_{(X')}$, etc., the \'etale covering $V'$
of~$U'$ extends in a unique way, up to unique isomorphism,
to an \'etale covering of~$X'$, and the $n_i$ are
prime to the residue characteristic $p$ of~$X$.
\end{proposition}

Uniqueness follows from the fact that $X'$ is normal
\eqref{XIII.5.1}; indeed
% original p. 429
an \'etale covering of~$X'$ that extends $V'$ is isomorphic to the
normalization of~$X'$ in the fiber of~$V'$ at the generic point
of~$X'$ \eqref{I.10.2}. If $\bar{x}'$ is a geometric point
of~$Y'$, we denote by $\bar{X}'$ the strict localization of~$X'$ at
$\bar{x}'$, $\bar{V}'=V'_{(\bar{X}')}$, etc. By descent, taking
uniqueness into account, it suffices to show that, for every
geometric point $\bar{x}'$ of~$Y'$, the \'etale covering
$\bar{V}'$ of~$\bar{U}'$ extends to $\bar{X}'$. Given
that an \'etale covering of an open subset of the regular scheme~$X'$
which contains all the points $x'$ such that $\dim
\cal{O}_{X',x'}\leq 1$ extends to all of $X$
(SGA~2 XIV~1.11), we may even restrict ourselves to the points $\bar{x}'$ which
project onto a maximal point of~$Y'$. Now, at such a point
$\bar{x}'$, the fact that $\bar{V}'$ extends to an \'etale
covering of~$\bar{X}'$ follows from~\Ref{X.3.6}.

Let us show that the $n_i$ are prime to $p$. Indeed, if this
were not so, we would have for example $p|n_1$. Replacing
$X$ by
$$X[T_1,\dots,T_r]/(T_1^{n_1/p}-f_1,T_2^{n_2}-f_2,\dots,T_r^{n_r}-f_r),$$
we reduce to the case where $X'=X[T_1]/(T_1^p-f_1)$. It suffices
to show that $V$ extends to an \'etale covering of~$X$,
because we will then have $n_1=1$, contrary to the hypothesis. For this we may
assume that $X$ is strictly local. Let $Z$ be the
closed subscheme of~$X$ with equation $p=0$ and $Z_1=Z\cap
X_{f_1}$; $Z_1$ is a nonempty open subset of~$Z$. By what
precedes, the \'etale covering $V'$ of~$U'$ extends
to an \'etale covering $W'$ of~$X'$. Let $W''_1$ and $W''_2$
be the inverse images of~$W'$ under the two projections $X''=X'\times_X
X'\rightrightarrows X'$, and let us show that the descent isomorphism
$u\colon W''_1|U''\to W''_2|U''$ extends to an $X''$-morphism
$W''_1\to W''_2$, which will necessarily be a descent
datum on~$W'$ relative to $X'\to X$. Let $Z''$
(\resp $Z''_1$) be the inverse image of~$Z$ (\resp $Z_1$) in $X''$. Since
the morphism $Z''\to Z$ is radicial, there exists an isomorphism
$v\colon W''_1|Z''\to W''_2|Z''$ which extends the isomorphism
$u|Z''_1$. But, since $X$ is henselian, we have a bijection
$$
\Hom_{X''}(W''_1,W''_2)\simeq\Hom_{Z''}(W''_1|Z'',W_2''|Z''),
$$
whence a morphism $w\colon W''_1\to W''_2$ lifting $v$. The
subscheme, at once open
% original p. 430
and closed, of~$X''$ over which $u$ and $w$ coincide
contains $Z''_1$, hence is equal to~$X''$, whence the fact that $V$
extends to~$X$.

\setcounter{subsection}{3} \setcounter{subsubsection}{-1}

\subsubsection{}
\label{XIII.5.3.0} Let us take up again the hypotheses and the notation
of~\Ref{XIII.5.2}, assuming moreover
$X$
strictly local. Then
it follows from \loccit that every connected \'etale covering
of~$U$ tamely ramified relative to $D$ is a
quotient of a tamely ramified covering of the
form
$$
U'=U[T_1,\dots,T_r]/(T_1^{n_1}-f_1,\dots,T_r^{n_r}-f_r),
$$
where the $n_i$ are integers prime to $p$. Let
$\bbmu_n$ be the group of $n$th roots of unity of
$U$. The group of $U$-automorphisms of~$U'$ is none other than the
group $\bbmu_{n_1}\times\dots\times\bbmu_{n_r}$, an
$n_i$th root of unity $\xi_i$ acting on~$U'$ by
transforming $T_i$ into $\xi_i T_i$. We therefore have the following statement:

\setcounter{subsection}{2}

\begin{corollary}
\label{XIII.5.3} Let $X$ be a regular strictly local scheme
of residue characteristic $p\geq 0$, $D=\sum_{1\leq i\leq
n} \divisor f_i$ a divisor with normal crossings on~$X$,
$U=X-\Supp D$. Set
$$
\tilde U=\varprojlim_{(n_i)}
U[T_1,\dots,T_r]/(T_1^{n_1}-f_1,\dots,T_r^{n_r}-f_r),
$$
the projective limit being taken over the filtered ordered
set (for the divisibility relation) of families of integers
$n_i>0$, prime to $p$. Then $\tilde U$ is a tamely ramified
universal covering of~$U$. Consequently the tamely ramified
fundamental group of~$U$ is
$$
\pi_1^{\tame}(U)\simeq \prod_{\ell\neq p} \ZZ_\ell [1]^r
\qquad \text{(canonical isomorphism)}
$$
where we have set
$\ZZ_\ell[1]=\varprojlim_{n>0}\bbmu_{\ell^n}$. The group
$\pi_1^{\tame}(U)$ is non-canonically isomorphic to $\prod_{\ell\neq p}
\ZZ_\ell^r$.
%
\end{corollary}

\begin{proposition}
\label{XIII.5.4} Let $f\colon X\to S$ be a morphism of
schemes, $D=\sum_{1\leq i\leq r} \divisor f_i$ a divisor with
normal crossings relative to~$S$ \eqref{XIII.2.1}, where,
% original p. 431
for each point $x$ of~$Y=\Supp D$, if $I(x)\subset[1,r]$ is
the set of $i$ such that $f_i(x)=0$, the subscheme
$V((f_i)_{i\in I(x)})$ is smooth over~$S$ of codimension
$\card I(x)$
in $X$. Let $U=X-Y$. Let $x$ be a point of
$Y$, $X_1=\Spec\cal{O}_{X,x}$, $U_1=U\times_X X_1$, $n_i$, $i\in I(x)$,
integers, and
$$
X'=X[T_i]_{i\in I(x)}/(T_i^{n_i}-f_i).
$$
Then, if $x'$ is the point of~$X'$ over~$x$, $X'$ is smooth
over~$S$ at $x'$. If the integers~$n_i$ are prime to the
characteristic $p$ of~$\kres(x)$, $U'_1=U_1\times_X X'$ is a
connected \'etale covering of~$U_1$ tamely
ramified over~$X_1$ relative to~$S_1$ \eqref{XIII.2.1.1}.
\end{proposition}

If $s=f(x)$, the geometric fiber $X'_{\bar{s}}$ is
regular at the point $x'$ \eqref{XIII.5.1}; since $X'$ is flat
over~$S$ in a neighborhood of~$Y$, this proves that $X'$ is smooth over~$S$ at
$x'$ (EGA~IV~12.1.6). If the integers $n_i$ are prime to $p$,
$U'_1$ is an \'etale covering of~$U_1$ \eqref{I.7.4}; it is
tamely ramified over~$X_1$ relative to~$S$ because
this is so on the geometric fibers at each point of~$S$
\eqref{XIII.5.1}. Finally the fact that $U'_1$ is connected follows from
(SGA~4 XVI~3.2).

\begin{proposition}[Relative Abhyankar lemma]
\label{XIII.5.5}
Let $X$ be an $S$-scheme, $D$ a
divisor with normal crossings relative to~$S$, as
in~\eqref{XIII.5.4}. Let $Y=X-\Supp D$,
$U=X- Y$, $x$ a point of~$Y$, $X_{1}$ the strict
localization of~$X$ at a geometric point over~$x$,
$U_{1}=U\times_{X}X_{1}$, $V_{1}$ an \'etale covering of
$U_{1}$. Assume that, for every maximal point $s$ of~$S$,
$V_{1\overline{s}}$ is tamely ramified over~$X_{1\overline{s}}$ relative to $\overline{s}$. Then one can
find integers $n_{i}$ prime to the characteristic $p$
of~$\kres(x)$, with $i\in I(x)$, such that, if we set
$$
X'_{1}=X_{1}[T_{i}]_{i\in I(x)}\big/ (T_{i}^{n_{i}}-f_{i}) \quoi,
$$
$U'_{1}=U_{1}\times_{X_{1}}X'_{1}$, etc., the \'etale covering
$V'_{1}$ of~$U'_{1}$ extends in a unique way, up to
unique isomorphism, to an \'etale covering of
$X'_{1}$. In particular
% original p. 432
$V_{1}$ is tamely ramified over~$X_{1}$ relative
to~$S$.
\end{proposition}

We may assume $S$ local noetherian with closed point~$f(x)$.
For each maximal point $s$ of~$S$ and for each $i\in I(x)$, let
$x_{i}$ be the generic point of the closed set $V(f_{i})$ of the fiber
$X_{1\overline{s}}$. The local ring $(\cal{O}_{X_{1},x_{i}})_{\red}$
is a discrete valuation ring with field of fractions $K_{i}$
and we have
$V_{1}|K_{i}=\Spec(\prod_{j\in I(x_{i})}L_{j})$,
where $L_{j}$ is a finite separable extension of~$K$; we denote by
$n_{j}$ the order of the inertia group of a Galois extension
generated by $L_{j}$ and by $n_{i}$ the l.c.m.\ of the $n_{j}$ as
$s$ runs through the maximal points of~$S$ and
$j\in I(x_{i})$.

The $n_{i}$ being thus chosen, we shall show that $V'_{1}$
extends in a unique way to an \'etale covering of
$X'_{1}$. Uniqueness follows from the fact that, $X'$ being smooth
over~$S$ at the points of~$Y$, we have $\prof\et_{Y'_{1}}(X'_{1})\geqslant 2$
(SGA~4 XVI~3.2 or SGA~2 XIV~1.19). Let $x'_{1}$ be a point of
$Y'_{1}$, $\overline{x}'_{1}$ a geometric point over~$x'_{1}$,
and denote by $\overline{X'_{1}}$ the strict localization of
$X'_{1}$ at $\overline{x}'_{1}$,
$\overline{U'_{1}}=U'_{1(\overline{X'_{1}})}$,~etc. By descent,
taking uniqueness into account, it suffices to show that
$\overline{V}'_{1}$ extends to $\overline{X'_{1}}$. Moreover we
may restrict ourselves to taking for $x'_{1}$ the maximal points of
$Y'_{1}$; indeed we will then have an extension of~$V'_{1}$ over an
open subset $W'_{1}$ of~$X'_{1}$ containing the maximal points of~$Y'_{1}$;
now, if $Z'_{1}=X'_{1}- W'_{1}$, we have
$\codim(Z'_{1s},X'_{1s})\geqslant 2$ if $s$ is a maximal point of
$S$ and $\codim(Z'_{1s},X'_{1s})\geqslant 1$, $\prof\et_{s}(S)\geqslant
1$ if $s$ is a point of~$S$ which is not maximal; the fact that
$V'_{1}$ extends to all of $X'_{1}$ then follows
from~(SGA~2 XIV~1.20). But, at a geometric point
$\overline{x}'_{1}$ over a maximal point of~$Y'_{1}$,
$X'_{1\red}$ is the spectrum of a discrete valuation ring, and
the fact that $\overline{V'_{1}}$ extends to $\overline{X'_{1}}$
then follows from~X~\Ref{X.3.6}.

Let us show that the $n_{i}$ are prime to $p$. Indeed, if this
were not so, there would be an index $i_{0}\in I(x)$ such that $p$
divides $n_{i_{0}}$. Replacing $X$ by
$X_{1}[T_{i_{0}},T_{i}]_{i\in I(x)}\big/
(T_{i}^{n_{i_{0}}/p}-f_{i_{0}},T_{i}^{n_{i}}-f_{i})$, we reduce
to the case where $X'_{1}=X_{1}[T]/T^{p}-f_{i_{0}}$. By what
precedes, the \'etale
% original p. 433
covering $V'_{1}$ of~$U'_{1}$ extends to an \'etale
covering $E'_{1}$ of~$X'_{1}$. Let $\eta$ be the closed point of
$S$; since the morphism $X'_{1\eta}\to X_{1\eta}$ is radicial,
$V_{1\eta}$ extends to an \'etale covering $E_{1\eta}$ of
$X_{1\eta}$. One then deduces, as in~\Ref{XIII.5.2}, that
$E'_{1}$ is endowed with a descent datum relative to the
morphism $X'_{1}\to X_{1}$, which extends the natural descent
datum that we have on~$E'_{1}|U'_{1}$. It follows that
$V_{1}$ extends to $X_{1}$; but this implies that we have
$n_{i_{0}}=1$, contrary to the hypothesis $n_{i_{0}}=p$.

\begin{corollary}
\label{XIII.5.6}
Let $X$ be an $S$-scheme, $D=\sum_{1\leqslant i\leqslant
r}\divisor f_{i}$ a divisor with normal crossings relative
to~$S$, as in~\Ref{XIII.5.4}. Let $\overline{x}$ be a geometric
point of~$X$, $\overline{X}$ the strict localization of
$X$ at $\overline{x}$, $\overline{Y}=Y_{(\overline{X})}$,
$\overline{U}=\overline{X}-\overline{Y}$ and
$$
\widetilde{U}= \varprojlim_{(n_{i})}
\overline{U}[T_{i}]_{i\in I(x)}\big/ (T_{i}^{n_{i}}-f_{i}) \quoi,
$$
the projective limit being taken over the filtered set of
families of integers $n_{i}>0$, prime to the characteristic
$p$ of~$\kres(x)$. Then $\widetilde{U}$ is a tamely ramified
universal covering of~$\overline{U}$ relative to~$S$. Consequently
the tamely ramified fundamental group
of~$\overline{U}$ is
$$
\pi_{1}^{\tame}(\overline{U})\simeq \prod_{\ell\neq
p}\ZZ_{\ell}[1]^{I(x)} \qquad \textrm{(canonical isomorphism)\quoi.}
$$
The group $\pi_{1}^{\tame}(\overline{U})$ is non-canonically
isomorphic to $\prod_{\ell\neq p}\ZZ_{\ell}^{I(x)}$.
\end{corollary}

\begin{subremark}
\label{XIII.5.6.1}
Let $X$ be an $S$-scheme, $D=\sum_{1\leqslant i\leqslant
r}\divisor f_{i}$ a divisor with normal crossings relative
to~$S$, as in~\Ref{XIII.5.4}, $U=X-\Supp D$.
For every subset $I\subset[1,r]$ let
$$
X_{I}= \Big(\tbigcap_{i\in I}V(f_{i})\Big) \cap \Big(\tbigcap_{i\in\complement
I}X_{f_{i}}\Big)\quoi .
$$
Let $p$ be an integer that is prime or zero and let $Z$ be a subset of
$X_{I}$ all of whose points are of characteristic~$p$. Let
$$
\widetilde{U}_{I}= \varprojlim_{(n_{i})}
U[T_{i}]_{i\in I}
\big/ (T_{i}^{n_{i}}-f_{i}) \quoi,
$$
where
% original p. 434
the projective limit is taken over the filtered set of
families of integers $n_{i}>0$, prime to $p$. Then, for every
geometric point $\overline{x}$ of~$Z$, the inverse image of
$\widetilde{U}_{I}$ on~$\overline{U}$ is identified with the tamely
ramified universal covering of~$\overline{U}$.
\end{subremark}

\begin{corollary}
\label{XIII.5.7}
The notation is that of~\Ref{XIII.5.6}. Let $\overline{S}$ be the
strict localization of~$S$ at~$\overline{x}$, and
$$
\overline{g}\colon\overline{U}\to\overline{S} \quad \text{and} \quad
\tilde{g}\colon\widetilde{U}\to\widetilde{S}
$$
the canonical morphisms. Then the morphisms $\overline{g}$ and
$\tilde g$ are $0$-acyclic \textup{(SGA~4 XV~1.3)}. Let $G$ be a
constructible sheaf of groups on~$\overline{S}$, $F=\overline{g}^*G$, $P$
a torsor under $F$. Then, for $P$ to be tamely
ramified over~$\overline{X}$ relative to $\overline{S}$, it
is necessary and sufficient that its inverse image $\widetilde{P}$ on~$\widetilde{U}$ be trivial.
\end{corollary}

Indeed, for every scheme
$\overline{X'}=\overline{X}[T_{i}]_{i\in I(x)}\big/
(T_{i}^{n_{i}}-f_{i})$, where the $n_{i}$ are integers~$>0$
prime to $p$, the morphism
$\overline{f}'\colon\overline{X'}\to\overline{S}$ is $0$-acyclic.
The geometric fibers of~$\overline{f}'$ at the various
points of~$\overline{S}$ are therefore connected and even
irreducible. The same therefore holds for the geometric fibers
of the morphisms
$\overline{g}'\colon\overline{U'}\to\overline{S}$, which proves that
the $\overline{g}'$, hence also $\tilde g$, are
$0$-acyclic (SGA~4 XV~1.16).

It is clear that a torsor $P$ on~$\overline{U}$ with group $F$, whose
inverse image on~$\widetilde{U}$ is trivial, is
tamely ramified over~$\overline{X}$ relative to
$\overline{S}$. Let us show that conversely, if $P$ is
tamely ramified over~$\overline{X}$ relative to
$\overline{S}$, its inverse image on~$\widetilde{U}$ is trivial.

It follows from (SGA~4 IX~2.14~(ii)) that one can find a
finite morphism $n\colon S_{1}\to\overline{S}$, a constant sheaf of groups
$C$ on~$S_{1}$, a monomorphism $G\to n_*C$.
Consider the following commutative diagram made up of
cartesian squares:
$$
\xymatrix{ \widetilde{U}_{1} \ar[r] \ar[d]_{r} &U_{1} \ar[r]^{g_{1}}
\ar[d]_{q} &S_{1} \ar[d]_{n} \cr \widetilde{U} \ar[r] &\overline{U}
\ar[r]^{\overline{g}} &\overline{S} \rlap{\quoi.}\cr }
$$
Let
% original p. 435
$C_{1}$ (\resp $\widetilde{C}_{1}$) be the inverse image of~$C$ on~$U_{1}$
(\resp on~$\widetilde{U}_{1}$). We have a commutative diagram, in
which $i$ and $j$ are isomorphisms (SGA~4 VIII~5.8):
\begin{equation*}
\label{eq:XIII.5.7.*}
\tag{$*$}
\begin{array}{c}
\xymatrix{ \H^{1}(\overline{U},q_*C_{1}) \ar[r]^{i}
\ar[d] &\H^{1}(U_{1},C_{1}) \ar[d] \cr
\H^{1}(\widetilde{U},r_*\widetilde{C}_{1}) \ar[r]^{j}
&\H^{1}(\widetilde{U}_{1},\widetilde{C}_{1}) \rlap{\quoi.}\cr }
\end{array}
\end{equation*}
Let $Q$ be the torsor under $q_*C_{1}$ deduced from~$P$ by
the extension of the structure group $F\to q_*C_{1}$.
By~\Ref{XIII.2.1.4}, $Q$ is tamely ramified
over~$\overline{X}$ relative to $\overline{S}$. To the torsor $Q$
there corresponds, by means of $i$, a torsor $Q_{1}$ under $C_{1}$, and
it is clear that $Q_{1}$ is tamely ramified over~$X_{1}=\overline{X}\times_{\overline{S}}S_{1}$ relative to~$S_{1}$. It therefore follows from~\Ref{XIII.5.6} that the inverse image
$\widetilde{Q}_{1}$ of~$Q_{1}$ on~$\widetilde{U}_{1}$ is trivial,
and diagram~\eqref{eq:XIII.5.7.*} then shows that the inverse image
$\widetilde{Q}$ of~$Q$ on~$\widetilde{U}$ is trivial.

Consider the following commutative diagram, whose second
row is exact (SGA~4 XII~3.1):
\begin{equation*}
\label{eq:XIII.5.7.**}
\tag{$**$}
\begin{array}{c}
\xymatrix{ \H^{0}(\overline{S},n_*C/G) \ar[r]
\ar[d]_{k} &\H^{1}(\overline{S},G)\rlap{$~=~1$} \ar[d] \cr
\H^{0}(\widetilde{U},r_*\widetilde{C}_{1}/\widetilde{F}) \ar[r]
&\H^{1}(\widetilde{U},\widetilde{F}) \ar[r]
&\H^{1}(\widetilde{U},r_*\widetilde{C}_{1}) \rlap{\quoi.} \cr }
\end{array}
\end{equation*}
Since the morphism $\widetilde{U}\to\overline{S}$ is $0$-acyclic,
$k$ is an isomorphism. The fact that $\widetilde{P}$ is trivial
then follows from~\eqref{eq:XIII.5.7.**}.

\section[Appendix II: finiteness for direct
images of stacks]{Appendix II: finiteness theorem for direct
images of stacks}
\label{XIII.6}

\begin{proposition}
\label{XIII.6.1}
Let $S$ be a locally noetherian scheme, $f\colon X\to S$
a morphism. If $S'$ is an $S$-scheme, we denote by $X'$
(\resp $f'$, etc.) the inverse image of~$X$ (\resp $f$, etc.) under the
morphism $S'\to S$. Suppose that, for every
% original p. 436
scheme $S'$ \'etale over~$S$ and every constructible sheaf of sets
$F$ on~$X'$, $f'_*F$ is constructible, and that, for
every constructible sheaf of groups $F$ on~$X'$, $\R^{1}f'_*F$
is constructible. Let $\Phi$ be a $1$-constructible stack on~$X$~\eqref{XIII.0}. Then $f_*\Phi$ is $1$-constructible.
\end{proposition}

For every scheme $S'$ \'etale over~$S$ and for every object $x$ of
$(f_*\Phi)_{S'}$, we have an isomorphism
$$
\SheafAut_{S'}(x)\simeq f'_*\SheafAut_{X'}(x) \quoi,
$$
where, in the right-hand side of the equality, $x$ is
regarded as an object of~$\Phi_{X'}$. The hypotheses
made therefore imply that $f_*\Phi$ is constructible.
Let $S\Phi$ be the sheaf of maximal subgerbes of
$\Phi$~\cite[III~2.1.7]{XIII.1}. Since $f_*(S\Phi)$ is constructible, we
can apply (SGA~4 IX~2.7) to it, and the fact that $f_*\Phi$ is
$1$-constructible then follows from the lemma which follows.

\begin{sublemma}
\label{XIII.6.1.1}
Let $S$ be a locally noetherian scheme, $f\colon X\to S$
a morphism, $\Phi$ a stack on~$X$. Suppose given a sheaf
on~$S$, representable by an \'etale $S$-scheme of finite
type $T$, a surjective morphism
$$
a\colon T\to f_*(S\Phi )
$$
and an object $p$ of the fiber $\Phi_{X_T}$ (where $X_T=X\times_ST$),
defining in $f_*(S\Phi )(T)=\allowbreak S\Phi (X_T)$ an element
equal to the image $q$ under $a$ of the identity section of
$T(T)$. Let $f_T\colon X_T\to T$ be the canonical morphism and suppose
that the sheaf $\R^1f_{T*}(\SheafAut_{X_T}(p))$ is constructible;
then so is~$S(f_*\Phi )$.
\end{sublemma}

The canonical morphism $f^*f_*\Phi\to\Phi$ gives a morphism
$$
S(f^*f_*\Phi )\simeq f^*(S(f_*\Phi ))\to S\Phi ,
$$
whence a canonical morphism
$$
\varphi\colon S(f_*\Phi )\to f_*(S\Phi ).
$$
Let
% original p. 437
$F=S(f_*\Phi )$ and let $G$ be the image of~$F$ under $\varphi$; by
(SGA~4 IX,~2.9) $G$ is a constructible sheaf.

It suffices to show that, for every point $s$ of~$S$, there exists a
nonempty open set~$U$ of~$s$ such that $F|U$ is locally constant
constructible. Let $s\in S$, $\overline{s}$ a geometric
point over~$s$, $\overline{q}_1,\dots,\overline{q}_n$ the elements of~$G_{\overline{s}}$. By
definition of~$T$, there exist $S$-morphisms
$h_i\colon\overline{s}\to S'$ such that
$\overline{q}_i=h_i^*(q)$. Let $S'$ be the fiber product over~$S$ of
$n$ schemes isomorphic to $T$, $\overline{s}\to S'$ the fiber
product of the $h_i$, $X'=X\times_SS'$, $q_i$ (\resp $p_i$) the inverse
image of~$q$ (\resp $p$) under the $i$th projection of~$S'$ onto~$T$. If $\Psi_i$ is the maximal subgerbe of~$\Phi|X'$
generated by $p_i$, the sheaf
$F_i=\R^1f_*'(\SheafAut_{X'}(p_i))$ is none other than the sheaf
$S(f_*'\Psi_i)$ of maximal subgerbes of~$f_*'\Psi_i$. In
particular the canonical injection $\Psi_i\to\Phi |X'$ gives a
morphism
$$
\alpha_i\colon F_i\to F|S'.
$$

We shall show that $\alpha_i$ is a bijection of~$F_i$ onto
the inverse image of~$q_i$ in $F|S'$. For every scheme $S''$
\'etale over~$S'$, every section $y$ of~$F_i(S'')$ has as image
$q_i|S''$ in $F(S'')$, because, locally for the \'etale topology
on~$S''$, $y$ is defined by an object~$x$ of~$\Phi_{X''}$ which is
isomorphic to $p_i|X''$. Conversely, if $y\in F(S'')$ has as
image $q_i|S''$ in $F(S'')$, then locally for the \'etale topology
on~$S''$, $y$ is defined by an object~$x$ of~$\Phi_{X''}$ which is
isomorphic to $p_i$; consequently $x$ is an object of~$\Psi_{iX''}$ and
consequently $y\in F_i(S'')$.

The proof is completed by using~\Ref{XIII.6.1.2}
below. Indeed, one can find an open neighborhood $U'$ of~$s$
such that $q_1|U',\dots,q_n|U'$ are sections of~$G(U')$ and
generate this sheaf. Since the $F_i|U'$ and $G|U'$ are
constructible, so is~$F|U'$,
by~\Ref{XIII.6.1.2}; replacing $U$ by a smaller
open set, $F|U$ is locally constant, which completes the
proof.

\begin{sublemma}
\label{XIII.6.1.2}
Let
% original p. 438
$S$ be a locally noetherian scheme, $F\to G$ a surjective
morphism of sheaves of groups on~$S$. Let $q_i$ be a
finite family of sections of~$G$ over~$X$ which generate $G$, and, for
each $i$, let $F_i$ be the subsheaf of~$F$ which is the inverse image of
$q_i$. Then, if $G$ and the $F_i$ are constructible, so is~$F$.
\end{sublemma}

To prove that $F$ is constructible, it suffices to show that, for
every point~$s$ of~$S$, there exists an open neighborhood $U$ of~$s$ such
that $F|U$ is locally constant constructible. So let $s$ be a
point of~$S$. Since the sheaves $F_i$ and $G$ are constructible, one
can find an open neighborhood $U$ of~$s$ such that $F_i|U$ and $G|U$
are locally constant. Let us then show that $F|U$ is locally
constant. By (SGA~4 IX~2.13~(i)), it suffices to see that, if
$\overline{s}$ is a geometric point over~$s$,
$\tilde{s}$ a geometric point of~$U$ and
$\overline{s}\to\tilde{s}$ a specialization morphism, the
canonical morphism
$$
F_{\tilde{s}}\to F_{\overline{s}}
$$
is bijective.

We consider the commutative diagrams
$$
\xymatrix{ (F_i)_{\tilde{s}} \ar[r]^-{\tilde{a}_i} \ar[d]^{\wr} &
F_{\tilde{s}} \ar[r]^-{\tilde{a}} \ar[d]_b & G_{\tilde{s}}
\ar[d]^{\wr}\\ (F_i)_{\overline{s}} \ar[r]^-{\overline{a}_i} &
F_{\overline{s}} \ar[r]^-{\overline{a}} & G_{\overline{s}} }
$$
Let $\overline{q}_i$ (\resp $\tilde{q}_i$) be the inverse image of~$q_i$
in $G_{\overline{s}}$ (\resp $G_{\tilde{s}}$); the morphisms
$\overline{a}$ and $\tilde{a}$ are surjective, and the morphism
$\overline{a}_i$ (\resp $\tilde{a}_i$) induces a bijection of
$(F_i)_{\overline{s}}$ (\resp $(F_i)_{\tilde{s}}$) onto~$\overline{a}^{-1}(q_i)$ (\resp onto~$\tilde{a}^{-1}(q_i)$). It therefore follows
from the diagram above that $b$ is an isomorphism.

\begin{corollary}
\label{XIII.6.2}
Let
% original p. 439
$S$ be a locally noetherian scheme, $f\colon X\to S$
a proper morphism. Let $\Phi$ be a $1$-constructible stack on~$X$;
then $f_*\Phi$ is a \hbox{$1$-constructible} stack.
\end{corollary}

The proof of~\Ref{XIII.6.1} also proves the following
result, taking \Ref{XIII.2.4}~2) into account.

\begin{corollary}
\label{XIII.6.3}
Let $S$ be a locally noetherian scheme, $f\colon X\to S$ a
morphism, $D$ a divisor on~$X$ with normal crossings
relative to~$S$ \eqref{XIII.2.1}, $Y=\Supp D$, $U=X-Y$,
$i\colon U\to X$ the canonical immersion. Let $\Phi$ be a stack on~$U$
given, locally for the \'etale topology on~$X$ and $S$,
as the inverse image of a $1$-constructible stack $\Psi$ on~$S$. Then
the stack $i_*^{\tame}\Phi$ is $1$-constructible.
\end{corollary}

\begin{thebibliography}{0}{XIII.7}

\bibitem[1]{XIII.1} J\ptbl \textsc{Giraud}, Cohomologie non ab\'elienne de
degr\'e~$2$, thesis, Paris (1966).

\bibitem[2]{XIII.2} J\ptbl \textsc{Giraud},
Cohomologie non ab\'elienne, Springer, Berlin-New\kern2pt York (1971).

\bibitem[3]{XIII.3} 
H\ptbl\textsc{Seifert}, W\ptbl\textsc{Threlfall}, Lehrbuch der Topologie, Chelsea, New
York (1934).


\bibitem[4]{XIII.4}
J-P\ptbl\textsc{Serre}, Cohomologie galoisienne, Lecture Notes
\no 5, Springer, Berlin (1964).

\bibitem[5]{XIII.5}
J-P\ptbl\textsc{Serre}, Corps locaux, Hermann, Paris (1962).
\end{thebibliography}

\end{document}
